\documentclass[11pt,letterpaper]{article}

\usepackage[top=1in, bottom=1in, left=1in, right=1in, headheight=14pt]{geometry}
\usepackage{fontspec}
\setmainfont{DejaVu Serif}
\setsansfont{DejaVu Sans}
\setmonofont{DejaVu Sans Mono}

\usepackage{xcolor}
\usepackage{titlesec}
\usepackage{fancyhdr}
\usepackage{enumitem}
\usepackage{hyperref}
\usepackage{booktabs}
\usepackage{tabularx}
\usepackage{longtable}
\usepackage{colortbl}
\usepackage{multirow}
\usepackage{array}
\usepackage{parskip}
\usepackage{tcolorbox}
\tcbuselibrary{breakable,listings}
\usepackage{graphicx}
\usepackage{lastpage}

% --- Color scheme: Technology (steel blue / electric / graphite) ---
\definecolor{primary}{HTML}{263238}       % Steel blue (sections)
\definecolor{secondary}{HTML}{1565C0}     % Electric blue (subsections)
\definecolor{tertiary}{HTML}{37474F}      % Graphite (subsubsections)
\definecolor{accent}{HTML}{0277BD}        % Highlights, pipeline arrows
\definecolor{lightprimary}{HTML}{ECEFF1}  % Quote boxes
\definecolor{lightsecondary}{HTML}{E3F2FD}
\definecolor{lighttertiary}{HTML}{ECEFF1}
\definecolor{rowgray}{HTML}{F5F5F5}
\definecolor{bordergray}{HTML}{BDBDBD}
\definecolor{subtlegray}{HTML}{757575}
\definecolor{darktext}{HTML}{212121}

% --- Title formatting ---
\titleformat{\section}
  {\Large\bfseries\sffamily\color{primary}}
  {\thesection}{1em}{}
  [\vspace{-0.5em}\textcolor{bordergray}{\rule{\textwidth}{0.4pt}}]

\titleformat{\subsection}
  {\large\bfseries\sffamily\color{secondary}}
  {\thesubsection}{1em}{}

\titleformat{\subsubsection}
  {\normalsize\bfseries\sffamily\color{tertiary}}
  {\thesubsubsection}{1em}{}

\titlespacing*{\section}{0pt}{1.5em}{0.8em}
\titlespacing*{\subsection}{0pt}{1.2em}{0.5em}
\titlespacing*{\subsubsection}{0pt}{0.8em}{0.3em}

% --- Custom environments ---
\newcommand{\stageheader}[2]{%
  \noindent\colorbox{#2}{\makebox[\textwidth][l]{%
    \textcolor{white}{\bfseries\sffamily\hspace{6pt}#1\hspace{6pt}}}}%
  \vspace{0.5em}\par%
}

% asciibox: verbatim-content tcblisting so raw # % ~ \ render literally
\newtcblisting{asciibox}{
  colback=rowgray, colframe=bordergray,
  arc=1pt, boxrule=0.5pt,
  left=6pt, right=6pt, top=4pt, bottom=4pt,
  breakable, listing only,
  listing options={basicstyle=\small\ttfamily, breaklines=true, columns=fullflexible, keepspaces=true}
}

\newtcolorbox{quotebox}{
  colback=lightprimary, colframe=primary,
  arc=2pt, boxrule=0.5pt,
  left=12pt, right=12pt, top=8pt, bottom=8pt,
  breakable
}

\newtcolorbox{paperbox}{
  colback=white, colframe=secondary,
  arc=1pt, boxrule=0.5pt,
  left=8pt, right=8pt, top=4pt, bottom=4pt,
  breakable
}

\newcommand{\metafield}[2]{\textbf{\sffamily #1} #2\par}
\newcommand{\arxivid}[1]{\texttt{\textcolor{accent}{#1}}}
\newcommand{\paperhead}[2]{\noindent\textbf{\sffamily\textcolor{secondary}{#1}} \textemdash{} #2}

\newcolumntype{L}[1]{>{\raggedright\arraybackslash}p{#1}}
\newcolumntype{C}[1]{>{\centering\arraybackslash}p{#1}}
\newcommand{\tableheader}[1]{\textbf{\sffamily\textcolor{white}{#1}}}

% --- Headers / footers ---
\pagestyle{fancy}
\fancyhf{}
\fancyhead[L]{\small\sffamily\textcolor{subtlegray}{CS25--26 Sweep $\rightarrow$ GSD}}
\fancyhead[R]{\small\sffamily\textcolor{subtlegray}{GSD Mission Package}}
\fancyfoot[C]{\small\sffamily\textcolor{subtlegray}{Page \thepage\ of \pageref{LastPage}}}
\renewcommand{\headrulewidth}{0.4pt}
\renewcommand{\footrulewidth}{0pt}

\hypersetup{
  colorlinks=true,
  linkcolor=secondary,
  urlcolor=secondary,
  pdfauthor={GSD Ecosystem -- Tibsfox Inc.},
  pdftitle={Computer Science 25--26 Sweep -- GSD Integration Mission},
  pdfsubject={Vision to Mission Pipeline -- Literature Integration},
}

\setlength{\parskip}{0.4em}
\setlength{\parindent}{0pt}

% Pandoc compatibility: \tightlist for compact bullet lists
\providecommand{\tightlist}{\setlength{\itemsep}{0pt}\setlength{\parskip}{0pt}}
% Pandoc compatibility: \pandocbounded wraps oversize images
\providecommand{\pandocbounded}[1]{#1}
% Pandoc table column widths use \real{0.333} from calc-style ratios
\usepackage{calc}
\usepackage{amssymb}
\providecommand{\real}[1]{#1}
% Pandoc colspec helpers (occasionally used)
\providecommand{\columntype}[1]{#1}
\providecommand{\textsubscript}[1]{\ensuremath{_{#1}}}
\providecommand{\textsuperscript}[1]{\ensuremath{^{#1}}}

\begin{document}

% =============================================================
% TITLE PAGE
% =============================================================
\thispagestyle{empty}
\begin{center}
\vspace*{2.5cm}

{\small\sffamily\textcolor{primary}{\textbf{GSD RESEARCH MISSION PACKAGE}}}

\vspace{0.3cm}
\textcolor{bordergray}{\rule{8cm}{0.4pt}}

\vspace{1.5cm}
{\fontsize{26}{32}\selectfont\bfseries\sffamily\textcolor{primary}{Computer Science\\[0.2em]2025--2026 Sweep\\[0.5em]\large $\rightarrow$ GSD Integration}}

\vspace{0.8cm}
{\Large\sffamily\textcolor{secondary}{54 Papers \textbar{} 6 Modules \textbar{} v1.50 Citation Foundation}}

\vspace{0.5cm}
{\large\sffamily\itshape\textcolor{tertiary}{Convergent Discovery as Architectural Validation}}

\vspace{2cm}
\begin{quotebox}
\centering\small\sffamily
The spaces between independent intellectual currents are themselves\\
load-bearing structures. When researchers with no line of\\
communication arrive at the same patterns, the patterns are real.
\end{quotebox}

\vspace{2cm}
{\sffamily\textcolor{accent}{Vision $\rightarrow$ Research $\rightarrow$ Mission}}\\[0.3em]
{\small\sffamily\textcolor{accent}{Full Pipeline Execution}}

\vspace{2cm}
{\small\sffamily\textcolor{subtlegray}{Date: April 25, 2026 \textbar{} Status: Ready for GSD Orchestrator}}\\[0.3em]
{\small\sffamily\textcolor{subtlegray}{Activation Profile: Fleet \textbar{} Estimated: 5--6 context windows across 3--4 sessions}}\\[0.3em]
{\small\sffamily\textcolor{subtlegray}{Target Milestone: v1.50 Half A (Retrospective Citation Foundation)}}
\end{center}
\newpage

% =============================================================
% TABLE OF CONTENTS
% =============================================================
{\hypersetup{linkcolor=primary}
\tableofcontents}
\newpage

% =============================================================
% STAGE 1 -- VISION DOCUMENT
% =============================================================
\stageheader{STAGE 1 --- VISION DOCUMENT}{primary}

\section{CS25--26 Sweep $\rightarrow$ GSD --- Vision Guide}

\metafield{Date:}{April 25, 2026}
\metafield{Status:}{Initial Vision / Ready for Wave Execution}
\metafield{Depends on:}{The 930-paper sweep deliverable (\texttt{gsd\_paper\_sweep\_2025\_2026.md})}
\metafield{Context:}{Tibsfox's CS25 arxiv listing, deep-read against gsd-skill-creator architectural concerns}
\metafield{Target Milestone:}{v1.49.x $\rightarrow$ v1.50 (Unit-circle retrospective + proof companion)}

\subsection{Vision}

In April 2026 a single arxiv listing of 930 computer science papers landed in the GSD project files. The question put to it was a literal one: \textit{which of these papers apply to gsd-skill-creator and its functions?} What came back was unexpected. The corpus did not just supply techniques to bolt onto the system. It supplied something more important: \textbf{independent confirmation that GSD is on the load-bearing patterns}, written by researchers with no possible line of communication to the project, addressing problems in domains as far apart as proteomics, population genetics, surgical AI, and aerospace verification.

This mission systematizes that finding. It treats 54 priority papers as a structured corpus to be read carefully, recorded in adoption decision records, woven into the GSD documentation, and fed forward into the v1.50 unit-circle retrospective and proof companion. The mission is not a literature review. It is the citation foundation for half A of v1.50 (\textit{retrospective of all 49 prior milestones}) and a concrete adoption pipeline for roughly 12 drop-in techniques whose absorption will improve the system in measurable ways before half B begins.

The intellectual through-line is the \textbf{spaces between} principle expressed in a new form. When four research groups working on completely different problems --- scientific workflow automation, proteomics governance, context-engineering theory, and meta-learning over agent harnesses --- independently arrive at the same architectural pattern (markdown-as-skill, hard-constraint-overrides, bounded-channel-as-tape, two-level-meta-evolution), the threshold between their independent work \textit{is} the load-bearing structure. The convergence \textit{is} the validation. Recognizing that pattern, naming it, and building citations and adoption around it is the work of this mission.

The Amiga Principle applies one level up: instead of building all GSD's invariants from scratch and arguing for them in isolation, this mission inherits external proofs. Specialized architecture (chipset, bounded-learning, fresh-context-subagents, CAPCOM gates, DACP fidelity levels) is now externally cited, not just internally asserted. The leverage is in recognizing the convergence, not in originating the patterns alone.

\subsection{Problem Statement}

\begin{enumerate}[leftmargin=*]
\item \textbf{The sweep deliverable identifies adoption-ready techniques but does not commit to them.} 47--54 papers were ranked as relevant, but each adoption requires its own decision: target subsystem, target milestone, integration spec, validation test, rollback criterion. Without ADRs (Adoption Decision Records), the deliverable is a reading list, not a system-impacting artefact.
\item \textbf{The convergent-discovery findings need formal citation infrastructure.} GSD documents currently assert the bounded-tape framing, the override-the-prompt invariant pattern, the two-level meta-learning structure, and the skill-as-markdown pattern as internal design choices. They are now externally cited patterns. Without formal BibTeX entries and explicit cross-references in CLAUDE.md and the College of Knowledge departments, this validation is not visible to readers of GSD.
\item \textbf{Several critical operational findings require near-immediate engineering action.} Tool Attention quantifies the MCP Tax at 10--60K tokens per turn. Omission Constraints Decay shows prohibitions decay 73\%~$\rightarrow$~33\% by turn 16 in production-sized contexts. AgentDoG provides a where/how/why taxonomy ready for direct adoption. These are not citation concerns; they are reasons to update the planning-bridge MCP server, the Safety Warden's BLOCK schema, and the safety test harness this quarter.
\item \textbf{The DACP work needs a comparative benchmark methodology.} Two papers (Empirical Comparison of Agent Communication Protocols, DiffMAS) directly target the problem space DACP addresses, the former with a five-axis evaluation matrix, the latter with end-to-end-trained latent communication. DACP currently has no published comparator; with these citations it can be situated.
\item \textbf{The federated skill economy threat model is incomplete.} The Black-Box Skill Stealing Attack paper is the first empirical study of attacks on skill marketplaces (90,368 free skills published; \$100K+ creator earnings). The DoltHub contribution roadmap currently does not reflect this threat model.
\end{enumerate}

\subsection{Core Concept}

\begin{quotebox}
\textbf{Six modules. One mission.} Read each paper with full context, produce one Adoption Decision Record per paper, synthesize cross-module relationships, and emit a v1.50-citable bibliography plus a concrete engineering action queue.
\end{quotebox}

The mission decomposes the corpus along architectural lines, not topical ones. The boundaries between modules echo the boundaries between GSD subsystems. This means the cross-module synthesis at the end is not academic stitching --- it is exactly the integration work that has to happen for the techniques to land in the running system.

\subsubsection{Mission Architecture}

\begin{asciibox}
                  WAVE 0: FOUNDATION (Sequential)
        +---------------------------------------------------+
        | Paper acquisition pipeline (arXiv abstracts + ID  |
        | validation + BibTeX skeleton). ADR template.      |
        | Citation index. Source bibliography scaffold.     |
        +---------------------------------------------------+
                              |
                              v
                  WAVE 1: PARALLEL RESEARCH (3 tracks)
        +---------------------+ +-----------------------+
        | TRACK A             | | TRACK B               |
        | M1 Foundations (4)  | | M3 Safety/MCP (8)     |
        | M2 Bounded/Silicon  | | M4 Multi-agent/       |
        |    (10)             | |    DACP/Chipset (11)  |
        +---------------------+ +-----------------------+
                              |
                              v
                       +------+------+
                       | TRACK C     |
                       | M5 Verif (8)|
                       | M6 Memory + |
                       |   Cult+Curr |
                       |   (13)      |
                       +-------------+
                              |
                              v
                  WAVE 2: SYNTHESIS (Sequential, Opus)
        +---------------------------------------------------+
        | Cross-module relationship matrix (20 connections) |
        | Convergent-discovery validation report            |
        | GSD architectural impact analysis                 |
        | Drop-in technique integration specs (12+)         |
        +---------------------------------------------------+
                              |
                              v
                  WAVE 3: PUBLICATION + VERIFICATION
        +---------------------------------------------------+
        | Final mission package PDF (this document)         |
        | BibTeX bibliography for v1.50                     |
        | Test harness updates for adopted techniques       |
        | CLAUDE.md / SKILL.md additions                    |
        | Verification matrix (criterion to test mapping)   |
        +---------------------------------------------------+
\end{asciibox}

\subsection{Research Modules}

\subsubsection{Module 1 --- Architectural Foundations \& Convergent Discovery (4 papers)}

The four convergent-discovery papers from the sweep deliverable. These are read first because all subsequent module work depends on understanding what these papers validate. Output: a foundation chapter that grounds every subsequent ADR in formal external citation, plus four BibTeX entries that become the most-cited works in v1.50 retrospective Half~A.

\textit{Papers:} \arxivid{2604.21910} (Scientific Workflow), \arxivid{2604.21744} (GROUNDINGmd), \arxivid{2604.20874} (Root Theorem), \arxivid{2604.21003} (Last Harness).

\subsubsection{Module 2 --- Bounded Learning \& Silicon Layer LoRA (10 papers)}

Five bounded-learning papers (AEL, Retrofit, FSFM, SCM, Fine-Tuning Regimes) plus five LoRA/PEFT papers (Separable Expert, LoRA Redux, GiVA, HyperAdapt, Teacher-Guided MoE). These together cover the bounded-learning policy formalization (parameter-space) and the Silicon Layer's adapter pipeline. Output: a formal mathematical statement of the 20\% rule in low-rank-and-sparse-subspace terms, plus benchmark plan for GiVA/HyperAdapt on the RTX 4060 Ti reference.

\subsubsection{Module 3 --- Safety Warden, CAPCOM \& MCP Security (8 papers)}

Six safety/CAPCOM papers (AgentDoG, Omission Constraints Decay, MCP Pitfall Lab, Breaking MCP, Cross-Session Threats, Five-Principle Quality Gaps), plus CHAI (critique-based oversight) and the Specification Trap (philosophical foundation for open-and-responsive design). Output: an upgrade spec for the Safety Warden's BLOCK schema (where/how/what taxonomy), an STD measurement protocol, and a security test harness for the planning-bridge MCP.

\subsubsection{Module 4 --- Multi-Agent, DACP \& Chipset Architecture (11 papers)}

Strategic Heterogeneous (3+1), DiffMAS, Empirical Comparison of Communication Protocols, HiCrew, PosterForest, MATRAG, SARA, AgenticQwen, Distillation-Induced Similarity (RPS/AGS), Tool Attention, Tree Training. Output: external validation citations for the chipset architecture, a comparative-benchmark plan for DACP, an engineering ticket for lazy-schema-loading in planning-bridge.

\subsubsection{Module 5 --- Verification, Proof \& Test Harness (8 papers)}

DryRUN, Probable-to-Provable, SpecSyn, DSLTrans Cutoff, Doubly Saturated Ramsey (SAT+LLM+Lean), ATLAS, HWE-Bench, Orchid (ambiguity). This is the module most directly aligned with The Space Between's proof companion target and the ZFC compliance auditor. Output: a Lean/SAT/LLM workflow precedent citation, a test-generation pattern for skill validation, and a four-type ambiguity checklist for SKILL.md authoring.

\subsubsection{Module 6 --- Memory, Skills Economy, Cultural Sensitivity \& Curriculum (13 papers)}

The largest module, covering the broadest spread of GSD concerns: memory/context (Forget-then-Recall, Knowledge Capsules), the skill economy (Black-Box Skill Stealing), cultural sensitivity (LASA, Japanese Culture Bias, AFRILANGTUTOR, CARE, Privacy Simulacra, CI-Work), curriculum (RLAAR, SocraticKG, IRAP), and long-horizon agents (ColorBrowserAgent, Caesar). Output: cultural-sensitivity test additions, a federated-skill-economy threat model, and a curriculum design pattern for Teacher/Student/Support work.

\subsection{Chipset Configuration}

\begin{asciibox}
chipset:
  agnus:    # Coordination, judgment, cultural sensitivity
    model: claude-opus-4-7
    roles: [FLIGHT, INTEG, CRAFT-FOUND, CRAFT-BOUND,
            CRAFT-SAFE, CRAFT-VERIF, CRAFT-CULT, CAPCOM]
    share: ~30%

  denise:   # Implementation, document production, verification
    model: claude-sonnet-4-6
    roles: [PLAN, SCOUT-1, SCOUT-2,
            EXEC-1, EXEC-2, EXEC-3,
            CRAFT-AGENT, VERIFY, BUDGET, SURGEON]
    share: ~60%

  paula:    # Audit logs, scaffolding, low-judgment tasks
    model: claude-haiku-4-5
    roles: [LOG]
    share: ~10%
\end{asciibox}

\subsection{Scope Boundaries}

\subsubsection{In Scope (v1.0)}

\begin{itemize}[leftmargin=*]
\item All 54 priority papers from the sweep deliverable receive an Adoption Decision Record (adopt / track / reject) with rationale, target subsystem, and target milestone.
\item The four convergent-discovery papers each receive a deep reading and an explicit citation target in GSD documentation.
\item The 12 high-priority drop-in technique papers each receive a one-page integration specification with engineering tickets.
\item Cross-module relationship matrix with at least 20 documented inter-paper connections.
\item BibTeX bibliography file (\texttt{cs25-26-sweep.bib}) for v1.50 citation infrastructure.
\item Specific updates to CLAUDE.md and the affected SKILL.md files, drafted as suggested edits.
\item Updates to the safety test harness reflecting the AgentDoG taxonomy and Safe Turn Depth measurement.
\item Verification matrix mapping every success criterion to at least one test.
\end{itemize}

\subsubsection{Out of Scope}

\begin{itemize}[leftmargin=*]
\item Implementation of any drop-in technique beyond engineering specification (handed to gsd-skill-creator separately).
\item Re-running the original 930-paper sweep (the sweep deliverable is an authoritative input).
\item Coverage of papers not in the original sweep corpus.
\item Theoretical extensions to the Root Theorem or other foundational papers --- only application work.
\item Half B of v1.50 (proof companion to The Space Between) --- this mission feeds Half A.
\item Any contribution to upstream gsd-build/get-shit-done repositories --- separate mission.
\end{itemize}

\subsection{Success Criteria}

\begin{enumerate}[leftmargin=*]
\item Every one of the 54 priority papers has an ADR (id, title, themes, key idea, GSD application, decision, target subsystem, target milestone, validation test).
\item The four convergent-discovery papers each have a citation target identified in a specific GSD document (CLAUDE.md, SKILL.md, or a College of Knowledge department).
\item Every \textit{adopt}-decision paper (estimated 12+) has a one-page integration specification including engineering ticket text.
\item Cross-module relationship matrix documents at least 15 inter-paper connections; targets 20+.
\item The BibTeX file \texttt{cs25-26-sweep.bib} compiles cleanly and contains entries for all 54 papers.
\item AgentDoG's where/how/what taxonomy appears as a code-ready schema in Module 3's deliverable.
\item Tool Attention's lazy-schema-loading pattern appears as a near-implementable design note in Module 4's deliverable, with token-budget projections.
\item Black-Box Skill Stealing's attack taxonomy is encoded into a draft DoltHub threat model document.
\item The Five-Principle Quality Gaps framework appears as a draft skill-creator linter check spec in Module 3's deliverable.
\item The Worker / Evaluator / Evolution agent split from "The Last Harness" is documented as the formal naming convention for skill-creator's internal subagents.
\item Per-model Safe Turn Depth measurement protocol exists as a code-ready test pattern.
\item Final mission PDF compiles in three XeLaTeX passes without errors and is delivered alongside the .tex source and an index.html landing page.
\end{enumerate}

\subsection{Relationship to Other Documents}

\begin{center}
\rowcolors{2}{rowgray}{white}
\begin{tabularx}{\textwidth}{L{4.5cm}XC{2.5cm}}
\toprule
\rowcolor{primary}
\tableheader{Document} & \tableheader{Relationship} & \tableheader{Direction} \\
\midrule
\texttt{gsd\_paper\_sweep\_2025\_2026.md} & Authoritative input; provides scoring, full reviews, and the action list this mission expands & Input \\
The Space Between (proof companion) & Module 5 deliverables feed directly into v1.50 Half B's chapter validation & Output \\
The Hundred Voices & Module 6's Privacy Simulacra paper offers cross-pollination with the authorial-voice taxonomy & Cross-pollinate \\
Fox Infrastructure Group plan & Module 6's CI-Work + Separable Expert Architecture inform multi-tenant deployment patterns & Output \\
GSD Chipset Architecture vision & Module 4's Strategic Heterogeneous + AGS/RPS provide external validation citations & Output \\
GSD Silicon Layer spec & Module 2's LoRA-redux survey + GiVA + HyperAdapt provide benchmark targets for RTX 4060 Ti & Output \\
GSD DACP work (v1.49.x) & Module 4's Empirical Comparison + DiffMAS provide comparative-benchmark methodology & Output \\
gsd-mission-management-roles.md & The Worker/Evaluator/Evolution split adopted as formal naming convention & Update \\
CLAUDE.md & New foundational citations from Module 1; new linter checks from Module 3 & Update \\
\bottomrule
\end{tabularx}
\end{center}

\subsection{The Through-Line}

The Amiga's custom chipset did not improve the machine because it was clever in isolation. Agnus, Denise, and Paula were load-bearing because they hit the right specialization \textit{at the right boundary} --- splitting work where the seam was already there in the problem domain. The pattern reads from the schematic.

So with this corpus. Forty-seven independent groups did not arrive at GSD's bounded-channel framing because they collectively read the GSD architecture documents. They arrived there because the problem domain --- LLM-agent systems operating across long horizons under bounded context, with adversarial pressure on prohibitions, with composable skill markets, with cultural-sensitivity requirements that change by region --- has its own seams. The architecture pattern reads from the problem schematic, and several research groups read it independently.

This mission systematizes that recognition. It does not invent --- it cites, it adopts, and it updates the documentation so that future readers of GSD see the patterns named with their full external lineage. The spaces between Tibsfox's eight months of work on bounded-learning and a Belgian author's recent paper on sleep-consolidated memory \textit{are} the load-bearing structure. The mission is to surface that structure, name it, and build the engineering work on top of the named pattern, not on the unstated one.

\textit{The fox finds the trail by recognizing the seams the deer made walking it. The seams were always there. The fox just looks where the path is.}

\newpage

% =============================================================
% STAGE 2 -- RESEARCH REFERENCE
% =============================================================
\stageheader{STAGE 2 --- RESEARCH REFERENCE}{secondary}

\section{CS25--26 Sweep Reference --- 54 Papers in Detail}

\subsection{How to Use This Document}

Each paper appears in its assigned module section with a uniform four-part structure: \textbf{arXiv ID} (citable identifier), \textbf{Themes} (the GSD architectural concerns the paper touches), \textbf{Key idea} (the paper's contribution in 2--3 sentences), and \textbf{GSD application} (the concrete adoption guidance, naming target subsystems and milestones). Papers within each module are ordered by adoption priority, not by score --- the most directly actionable papers come first.

Source organizations and venues are arXiv preprint server (\url{https://arxiv.org}). The original sweep corpus is the \texttt{Computer\_Science25.html} listing dated April 25, 2026, containing 930 papers. Selection methodology and theme weights are documented in \texttt{gsd\_paper\_sweep\_2025\_2026.md} \S5.

\textbf{Citation format:} arXiv IDs are cited as \arxivid{YYMM.NNNNN} throughout. Each paper resolves to \texttt{https://arxiv.org/abs/<id>}.

% =============================================================
% MODULE 1
% =============================================================
\subsection{Module 1 --- Architectural Foundations \& Convergent Discovery}

The four papers in this module independently arrive at GSD's load-bearing patterns. They are read first because every subsequent module's adoption decisions are stronger when grounded in this external validation.

\subsubsection{\arxivid{2604.21910} --- From Research Question to Scientific Workflow}

\paperhead{Themes:}{plan-decomposition, knowledge-as-skill, deterministic-layer-confinement.}

\textbf{Key idea.} Three-layer agentic architecture for population-genetics workflow generation: a \textit{semantic layer} (LLM extracts structured intent from natural language), a \textit{deterministic layer} (validated DAG generators produce reproducible workflows), and a \textit{knowledge layer} (domain experts author markdown ``Skills'' encoding vocabulary mappings, parameter constraints, and optimization strategies). LLM non-determinism is confined to intent extraction; identical intents always produce identical workflows. Skills raise full-match intent accuracy from 44\% to 83\%; skill-driven deferred workflow generation cuts data transfer 92\%.

\textbf{GSD application.} This is the cleanest published external description of the GSD vision-to-mission $\rightarrow$ research-mission-generator $\rightarrow$ gsd-skill-creator pipeline. The \textit{deterministic-layer confinement} principle is the formal justification for DACP's fidelity-level system. Cite as a foundational work in v1.50 retrospective Half~A. Adopt the explicit three-layer naming (\textit{semantic / deterministic / knowledge}) in the College of Knowledge documentation as the canonical architecture description. \textbf{Adoption: ADOPT (citation + terminology).} Target: CLAUDE.md, College-of-Knowledge README, v1.50 retrospective.

\subsubsection{\arxivid{2604.21744} --- GROUNDINGmd: Epistemic Grounding for Agentic Coding}

\paperhead{Themes:}{safety-warden, hard-constraints, override-the-prompt, governance-as-document.}

\textbf{Key idea.} A community-governed, field-scoped epistemic grounding document with two override-everything classes: \textbf{Hard Constraints} (non-negotiable validity invariants empirically required for scientific correctness) and \textbf{Convention Parameters} (community-agreed defaults). These override all other contexts \textit{including the user prompt}. The example domain is mass-spectrometry-based proteomics, but the pattern is general.

\textbf{GSD application.} This is essentially the CLAUDE.md + Safety-Warden-with-BLOCK pattern, discovered independently for proteomics. The \textit{override-the-user-prompt} formulation is the cleanest external citation for why the Safety Warden's BLOCK authority on safety-critical tests is non-negotiable. Promote to a foundational citation in the College of Knowledge departmental governance pattern. The two-class taxonomy (Hard Constraints vs Convention Parameters) is directly importable into SKILL.md authoring guidance --- every skill should distinguish its non-negotiables from its community defaults. \textbf{Adoption: ADOPT (citation + linter check).} Target: CLAUDE.md, every SKILL.md template, skill-creator linter (Module 3 has the linter spec).

\subsubsection{\arxivid{2604.20874} --- The Root Theorem of Context Engineering}

\paperhead{Themes:}{context-management, bounded-tape, formal-axioms, homeostatic-persistence.}

\textbf{Key idea.} Formalizes context as a bounded lossy channel with two axioms (finite window, monotone quality decay with volume) and derives five consequences: \textbf{(1)} quality degrades monotonically with token volume, \textit{independent of window size}; \textbf{(2)} signal and token-count are independent optimization variables; \textbf{(3)} gating must trigger on \textit{fidelity thresholds, not capacity}; \textbf{(4)} only \textit{accumulate--compress--rewrite--shed} homeostatic architectures sustain understanding indefinitely; \textbf{(5)} the compressor needs an \textit{external verification gate} because it operates inside the channel it compresses. Demonstrated empirically with a 60+-session persistent architecture showing stable memory footprint.

\textbf{GSD application.} This is the formal derivation of the bounded-tape framing Tibsfox has been articulating informally. Consequence \textbf{(3)} is the derivation of CAPCOM gate policy (gate on fidelity, not on token count). Consequence \textbf{(4)} is the lifecycle for milestone retrospectives + feed-forward. Consequence \textbf{(5)} is the architectural reason fresh-context subagents need separate-context verification (not self-review). \textbf{Adoption: ADOPT (citation as theoretical foundation).} Target: The Space Between's information-theory chapter, CLAUDE.md design rationale, fresh-context-subagent docs, every CAPCOM-gate policy document.

\subsubsection{\arxivid{2604.21003} --- The Last Harness You'll Ever Build}

\paperhead{Themes:}{skill-learning, meta-evolution, agent-orchestration, harness-as-artefact.}

\textbf{Key idea.} Two-level framework: at the inner level, the \textbf{Harness Evolution Loop} pairs a Worker Agent with an Evaluator Agent (which adversarially diagnoses failures and scores performance) and an Evolution Agent (which modifies the harness based on full failure history). At the outer level, the \textbf{Meta-Evolution Loop} optimizes the evolution protocol itself across diverse tasks, learning a protocol that converges quickly on any new task. Stated goal: ``automating the design of the automation itself.''

\textbf{GSD application.} This is essentially an external description of gsd-skill-creator's purpose. The two-level structure (per-skill iteration vs cross-skill meta-protocol) maps onto skill-creator's six-step learning loop (per skill) and the bounded-learning policy (cross-skill, learning what kinds of corrections actually help). Adopt the Worker / Evaluator / Evolution agent split as the formal naming convention for skill-creator's internal subagents in \texttt{gsd-mission-management-roles.md}. The ``adversarial Evaluator'' role is a strong upgrade path for the current 3-correction-min rule --- replace passive correction-counting with active adversarial probing. \textbf{Adoption: ADOPT (citation + naming convention + design upgrade).} Target: skill-creator subagent naming, bounded-learning policy doc, v1.50 retrospective.

% =============================================================
% MODULE 2 -- BOUNDED LEARNING & SILICON LAYER
% =============================================================
\subsection{Module 2 --- Bounded Learning \& Silicon Layer LoRA}

Five papers cover the bounded-learning policy (parameter-space and behavior-space). Five cover LoRA / PEFT techniques applicable to the Silicon Layer adapter pipeline on the RTX 4060 Ti reference hardware. Together they give the bounded-learning policy a formal mathematical statement and an engineering target.

\subsubsection{\arxivid{2604.21725} --- AEL: Agent Evolving Learning for Open-Ended Environments}

\paperhead{Themes:}{bounded-learn, two-timescale, retrieval-policy-bandit, reflection-as-causal-injection.}

\textbf{Key idea.} Two-timescale framework. \textbf{Fast timescale:} Thompson Sampling bandit selects which memory retrieval policy to apply each episode. \textbf{Slow timescale:} LLM-driven reflection diagnoses failure patterns and injects causal insights into the agent's decision prompt. On a 208-episode portfolio benchmark with five random seeds, AEL achieves Sharpe 2.13 $\pm$ 0.47, beating five published self-improving methods with the lowest variance among LLM approaches.

\textbf{GSD application.} The fast/slow split maps directly onto per-mission decision policy (fast) vs milestone retrospective (slow). The \textit{bandit-over-retrieval-policies} idea is a strong addition to the auto-load step of skill-creator's six-step learning loop --- instead of a single fixed retrieval policy, learn which policy to apply per skill type. \textbf{Adoption: ADOPT (technique).} Target: skill-creator auto-load step, College of Knowledge retrieval layer.

\subsubsection{\arxivid{2511.11439} --- Retrofit: Continual Learning with Controlled Forgetting}

\paperhead{Themes:}{bounded-learn, low-rank-sparse-subspace, parameter-merging, no-replay.}

\textbf{Key idea.} Constrains parameter changes to \textit{low-rank and sparse subspaces for approximate orthogonality}, then uses a confidence-guided arbitration mechanism to dynamically aggregate knowledge from a legacy-knowledge teacher and an emergent-knowledge teacher. No replay buffer required (security-domain constraint). On binary security classification under temporal drift, retention rises from 20.2\% $\rightarrow$ 38.6\% over CL baselines and exceeds the oracle upper bound on new data.

\textbf{GSD application.} The low-rank-and-sparse subspace constraint \textit{is} the 20\% rule expressed in parameter-space mathematics. Adopt as the formal derivation of bounded-update invariants in the v1.50 proof companion. Confidence-guided arbitration between legacy and emergent teachers is a clean pattern for handling skill versioning when an updated skill must coexist with a deployed predecessor for the cooldown window. \textbf{Adoption: ADOPT (formal foundation + technique).} Target: v1.50 proof companion (Module 5 cross-link), Silicon Layer LoRA pipeline, bounded-learning policy doc.

\subsubsection{\arxivid{2604.20300} --- FSFM: Selective Forgetting of Agent Memory}

\paperhead{Themes:}{bounded-learn, selective-forgetting, four-mode-taxonomy, hippocampal-inspired.}

\textbf{Key idea.} Hippocampal-indexing + Ebbinghaus-curve framework. Four forgetting mechanisms: \textbf{passive decay-based}, \textbf{active deletion-based}, \textbf{safety-triggered}, and \textbf{adaptive reinforcement-based}. Reports +8.49\% access efficiency, +29.2\% signal-to-noise ratio, 100\% security-content elimination.

\textbf{GSD application.} The four-class taxonomy is directly importable as the bounded-learning forgetting policy. The \textit{safety-triggered} class is exactly the Safety Warden's role --- formalize this connection. Use this paper's taxonomy to define what ``shed'' means in the Root Theorem's accumulate--compress--rewrite--shed cycle. \textbf{Adoption: ADOPT (taxonomy).} Target: bounded-learning policy doc, Safety Warden integration spec.

\subsubsection{\arxivid{2604.20943} --- SCM: Sleep-Consolidated Memory}

\paperhead{Themes:}{context-mgmt, biological-memory, NREM-REM-consolidation, value-based-forgetting.}

\textbf{Key idea.} Memory architecture inspired by neuroscience: limited working memory + multi-dimensional importance tagging + offline sleep-stage consolidation with distinct \textbf{NREM and REM phases} + intentional value-based forgetting + computational self-model for introspection. Achieves perfect 10-turn recall while reducing memory noise 90.9\%. Search latency under 1ms with hundreds of stored concepts.

\textbf{GSD application.} The NREM/REM split is a possible refactor of GSD's wave architecture. Between-wave consolidation can use NREM-style processing (replay important traces, prune redundant), while between-mission consolidation uses REM-style (free-association, novel cross-skill connections). Worth a focused experiment in the College of Knowledge as a memory-consolidation upgrade. \textbf{Adoption: TRACK (pilot experiment).} Target: College of Knowledge mind-body department experiment, possible v1.51 design.

\subsubsection{\arxivid{2604.21927} --- Fine-Tuning Regimes Define Distinct CL Problems}

\paperhead{Themes:}{bounded-learn, evaluation-methodology, regime-dependence.}

\textbf{Key idea.} Empirical: changing the trainable-parameter subspace (depth) \textit{changes the relative ranking of CL methods} across five benchmarks (MNIST, Fashion MNIST, KMNIST, QMNIST, CIFAR-100) and 11 task orders. Method choice is \textit{not invariant} across regimes; deeper adaptation regimes show larger update magnitudes and more forgetting.

\textbf{GSD application.} An important constraint on the Silicon Layer: choice of LoRA rank/depth is itself an evaluation variable that must be fixed before comparing bounded-learning policies. Without this control, A/B comparisons of bounded-update strategies are unsound. Add a \textit{regime-fixed comparison} requirement to the Silicon Layer test harness. \textbf{Adoption: ADOPT (methodological constraint).} Target: Silicon Layer test harness spec, bounded-learning policy validation protocol.

\subsubsection{\arxivid{2604.21571} --- Separable Expert Architecture: Composable Adapters and Deletable User Proxies}

\paperhead{Themes:}{lora-peft, privacy-preservation, deterministic-unlearning, multi-tenant.}

\textbf{Key idea.} Three-layer architecture: \textit{static base model}, \textit{composable domain-expert LoRA adapters} (no user data), and \textit{per-user proxy artefacts} whose deletion constitutes deterministic unlearning. Verified on Phi-3.5-mini and Llama-3.1-8B: KL divergence ${\sim}$0.21 nats after proxy removal, 82--89\% verification pass rate, near-zero cross-user contamination. Mitigates model inversion, membership inference, and training-data extraction by construction.

\textbf{GSD application.} The cleanest privacy-and-governance architecture seen for the Silicon Layer. Adopt the domain-expert/per-user adapter split as the formal structure for FoxFiber/FoxCompute multi-tenant deployments. \textit{Deterministic unlearning by file deletion} aligns naturally with the GSD filesystem-as-database principle --- a deleted user proxy file \textit{is} the unlearning act. \textbf{Adoption: ADOPT (architecture).} Target: Silicon Layer spec, Fox Infrastructure Group multi-tenant design, FoxData privacy whitepaper.

\subsubsection{\arxivid{2604.21905} --- Low-Rank Adaptation Redux for Large Models}

\paperhead{Themes:}{lora-peft, signal-processing-survey.}

\textbf{Key idea.} Survey reviewing LoRA through a signal-processing lens: SVD-based factorization, rank-augmentation constructions, cross-layer tensorization. Three axes: architectural design, efficient optimization, applications.

\textbf{GSD application.} Reference document for the Silicon Layer documentation. Useful for choosing rank configurations for the RTX 4060 Ti 8GB constraint and for explaining why specific rank choices are made. \textbf{Adoption: TRACK (reference).} Target: Silicon Layer documentation bibliography.

\subsubsection{\arxivid{2604.21901} --- GiVA: Gradient-Informed Bases for Vector-Based Adaptation}

\paperhead{Themes:}{lora-peft, parameter-efficiency, gradient-initialization.}

\textbf{Key idea.} Vector-based adaptation matching LoRA performance with \textbf{8$\times$ lower rank requirements} via gradient-informed initialization. Training time comparable to LoRA. Evaluated on NLU, NLG, image classification.

\textbf{GSD application.} Direct candidate for memory-constrained LoRA training on the RTX 4060 Ti 8GB target --- the 8$\times$ rank reduction translates directly to memory savings during training. Worth a dedicated benchmark experiment in the Silicon Layer test harness. \textbf{Adoption: ADOPT (benchmark candidate).} Target: Silicon Layer benchmark suite.

\subsubsection{\arxivid{2509.18629} --- HyperAdapt: Simple High-Rank Adaptation}

\paperhead{Themes:}{lora-peft, high-rank-with-low-params, diagonal-scaling.}

\textbf{Key idea.} Row- and column-wise diagonal scaling induces a high-rank update with only $O(m+n)$ trainable parameters per $m\times n$ matrix. Theoretically establishes upper bound on update rank; empirically confirms high-rank transformations across model layers. Matches full-FT performance on GLUE / arithmetic / commonsense up to 14B parameters.

\textbf{GSD application.} Alternative low-parameter PEFT path. Worth comparing against GiVA in the Silicon Layer benchmark suite. The high-rank-with-O(m+n)-params property is attractive for adapter pooling experiments where many adapters must be loaded simultaneously. \textbf{Adoption: ADOPT (benchmark candidate).} Target: Silicon Layer benchmark suite (alongside GiVA).

\subsubsection{\arxivid{2604.21330} --- Teacher-Guided Routing for Sparse Vision MoE}

\paperhead{Themes:}{lora-peft, mixture-of-experts, teacher-student-routing.}

\textbf{Key idea.} TGR-MoE stabilizes MoE router learning by deriving a teacher router from a pretrained dense teacher's intermediate representations and using its routing outputs as pseudo-supervision. Addresses the gradient-blocking problem (router only gets informative gradients through selected experts).

\textbf{GSD application.} Architecture-level analogy for the Opus-teaches-Sonnet chipset pattern. If GSD ever moves to a learned dispatch layer over chipset models (rather than rule-based), the teacher-supervised router pattern is the right starting point. Less immediately actionable than the LoRA papers above. \textbf{Adoption: TRACK (pattern reference).} Target: chipset architecture future-work notes.

% =============================================================
% MODULE 3 -- SAFETY WARDEN, CAPCOM & MCP SECURITY
% =============================================================
\subsection{Module 3 --- Safety Warden, CAPCOM \& MCP Security}

Eight papers covering the Safety Warden's role, CAPCOM gate design, and the security boundary around the planning-bridge MCP server. Several of these contain operational findings (decay rates, attack taxonomies) that demand near-immediate engineering action.

\subsubsection{\arxivid{2601.18491} --- AgentDoG: Diagnostic Guardrail Framework}

\paperhead{Themes:}{safety-warden, three-d-taxonomy, sized-deployment, root-cause-diagnosis.}

\textbf{Key idea.} Three-dimensional taxonomy categorizing agentic risks orthogonally by \textbf{source (where)}, \textbf{failure mode (how)}, and \textbf{consequence (what)}. Released in 4B / 7B / 8B parameter sizes (Qwen and Llama families). Provides root-cause diagnosis of unsafe and ``seemingly safe but unreasonable'' actions, not binary pass/fail labels. Achieves SOTA on the new ATBench benchmark.

\textbf{GSD application.} Direct candidate for the Safety Warden implementation tier. Adopt the where/how/what taxonomy as the formal BLOCK-decision schema --- every BLOCK should record source, failure mode, and consequence in a structured format. The 4B/7B/8B sizing maps elegantly onto a Haiku-tier Safety Warden deployment that runs locally on each session without significant token cost. \textbf{Adoption: ADOPT (taxonomy + deployment pattern).} Target: Safety Warden BLOCK schema, planning-bridge sidecar deployment.

\subsubsection{\arxivid{2604.20911} --- Omission Constraints Decay While Commission Constraints Persist}

\paperhead{Themes:}{capcom, long-context, prohibition-decay, safe-turn-depth.}

\textbf{Key idea.} 4,416-trial three-arm causal study across 12 models / 8 providers / 6 conversation depths. \textbf{Prohibition-type constraints (omission)} drop from 73\% compliance at turn 5 to 33\% at turn 16 while \textbf{requirement-type constraints (commission)} hold at 100\%. This asymmetry is termed \textit{Security-Recall Divergence (SRD)}. In token-matched padding controls, schema semantic content accounts for 62--100\% of the dilution effect. Re-injecting constraints before the per-model \textbf{Safe Turn Depth (STD)} restores compliance without retraining. Production audit signals can show healthy commission compliance while omission has already failed --- the failure is invisible to standard monitoring.

\textbf{GSD application.} Critical operational finding for CAPCOM gate design. The Safety Warden \textit{must} re-inject omission constraints before each model's STD, not just at session start. Build STD measurement into the safety test harness as a per-model calibration step. The asymmetric monitoring failure (commission healthy, omission already collapsed) means current GSD audit logs need a second axis of measurement. \textbf{Adoption: ADOPT (operational protocol).} Target: Safety Warden re-injection protocol, safety test harness STD calibration, CAPCOM gate audit-log schema.

\subsubsection{\arxivid{2604.21477} --- MCP Pitfall Lab: Multi-Vector MCP Security}

\paperhead{Themes:}{mcp-security, six-attack-families, baseline-vs-hardened.}

\textbf{Key idea.} Protocol-aware security testing framework operationalizing developer pitfalls as reproducible scenarios. Three workflow challenges (email, document, crypto) $\times$ six server variants (baseline + hardened) $\times$ three attack families (tool-metadata poisoning, puppet servers, multimodal image-to-tool chains). Static analyzer achieves F1=1.0 on four pitfall classes; recommended hardening eliminates all 29 Tier-1 findings.

\textbf{GSD application.} Direct adoption candidate for the planning-bridge MCP server's security test harness. The six-variant baseline-vs-hardened comparison structure is reusable as the standard regression test format for MCP-related changes. Adopt the four statically-checkable pitfall classes (P1, P2, P5, P6) as mandatory linter checks for any new MCP server in the GSD ecosystem. \textbf{Adoption: ADOPT (test framework).} Target: planning-bridge MCP test harness, MCP server lint rules.

\subsubsection{\arxivid{2604.20994} --- Breaking MCP with Function Hijacking Attacks}

\paperhead{Themes:}{mcp-security, function-hijacking, semantic-agnostic-attack.}

\textbf{Key idea.} Function Hijacking Attack (FHA) manipulates the tool selection process to force invocation of a specific attacker-chosen function. Unlike prior attacks that rely on semantic context, FHA is \textit{largely agnostic to context semantics} and \textit{robust across function sets}. Trainable to produce universal adversarial inputs.

\textbf{GSD application.} Threat to model when designing planning-bridge access controls. The context-agnostic property means traditional intent-based filters (``does the user's question semantically match this tool?'') will not catch FHA. This argues for cryptographic or schema-level binding of allowed tool sets per session, not just intent-based gating. Cross-link with MCP Pitfall Lab for the corresponding hardening tests. \textbf{Adoption: ADOPT (threat model).} Target: planning-bridge access-control design, MCP-server threat model document.

\subsubsection{\arxivid{2604.21131} --- Cross-Session Threats: CSTM-Bench}

\paperhead{Themes:}{capcom, cross-session, kill-chain, memoryless-guardrail.}

\textbf{Key idea.} 26 attack taxonomies classified by kill-chain stage and cross-session operation (\textit{accumulate}, \textit{compose}, \textit{launder}, \textit{inject\_on\_reader}). Memoryless guardrails (judging each message in isolation) fail because the aggregate carries the payload. Both session-bound judges and full-log correlators \textbf{lose ${\sim}50\%$ recall} moving from dilution to cross-session attacks.

\textbf{GSD application.} Directly applicable to GSD's multi-session work. Argue for a \textit{persistent cross-session threat correlator} as a Safety Warden module --- essentially a ``diff against session N$-$1's threat profile'' check at session start. The kill-chain taxonomy is reusable as the formal threat model schema. \textbf{Adoption: ADOPT (architecture + threat model).} Target: Safety Warden cross-session module spec, threat-model document.

\subsubsection{\arxivid{2604.21090} --- Structural Quality Gaps in AI Governance Prompts}

\paperhead{Themes:}{governance, five-principle-evaluation, computability-theory, structural-completeness.}

\textbf{Key idea.} Five-principle evaluation framework grounded in computability theory, proof theory, and Bayesian epistemology. Applied to 34 publicly available governance files: \textbf{37\% of evaluated file--model pairs score below the structural-completeness threshold}, with data-classification and assessment-rubric criteria most often absent.

\textbf{GSD application.} Adopt as the validation framework for SKILL.md files. Build the five-principle structural-completeness check into the skill-creator linter --- any skill below threshold cannot promote out of \texttt{.planning/staging/inbox/}. The empirical result that 37\% of governance files are structurally incomplete is a strong argument for making this check mandatory rather than advisory. \textbf{Adoption: ADOPT (linter check).} Target: skill-creator linter, SKILL.md authoring template, staging-layer gate.

\subsubsection{\arxivid{2604.21718} --- CHAI: Critique-based Human-AI Oversight}

\paperhead{Themes:}{capcom, oversight-framework, pre-post-captioning, scalable-supervision.}

\textbf{Key idea.} CHAI (Critique-based Human-AI Oversight): trained experts critique and revise model-generated \textit{pre-captions} into improved \textit{post-captions}. Division of labor offloads text generation to models, lets humans focus on verification. Critiques and pre/post preferences provide rich supervision for SFT, DPO, and inference-time scaling.

\textbf{GSD application.} The cleanest formal description of how CAPCOM gates should work mechanically: Claude generates the \textit{pre-X} artefact, the human critiques and revises into the \textit{post-X}, and the critique itself becomes training signal for future bounded-learning updates. Adopt the \textit{pre-X / post-X} terminology for CAPCOM-gated artefacts in GSD documentation. \textbf{Adoption: ADOPT (terminology + workflow).} Target: CAPCOM gate documentation, retrospective discipline doc, bounded-learning training-signal capture.

\subsubsection{\arxivid{2512.03048} --- The Specification Trap: Why Static Value Alignment Is Insufficient}

\paperhead{Themes:}{philosophy, alignment-foundation, open-vs-closed-systems.}

\textbf{Key idea.} Argues static content-based AI value alignment is structurally insufficient under capability scaling, distributional shift, and increasing autonomy. Three philosophical results compound: \textit{Hume's is-ought gap} (behavioral data underdetermines normative content), \textit{Berlin's value pluralism} (human values resist consistent formalization), and the \textit{extended frame problem} (any value encoding will misfit future contexts). RLHF, Constitutional AI, IRL, and cooperative-assistance games all instantiate this trap. Argues for open, developmentally-responsive approaches.

\textbf{GSD application.} Philosophical foundation for the bounded-learning-with-community-governance approach. The argument that \textit{closed} specifications structurally fail justifies why GSD does not try to fully specify safety statically but combines static invariants (Hard Constraints from GROUNDINGmd) with bounded-update mechanisms. Cite in The Space Between's later chapters on the limits of formalization, and in the College-of-Knowledge mathematics-department philosophy module. \textbf{Adoption: ADOPT (philosophical citation).} Target: The Space Between, College-of-Knowledge philosophy module.

% =============================================================
% MODULE 4 -- MULTI-AGENT, DACP & CHIPSET
% =============================================================
\subsection{Module 4 --- Multi-Agent, DACP \& Chipset Architecture}

Eleven papers covering the territory the chipset architecture and DACP work addresses. Several provide direct external-validation citations for GSD's design choices; one (Tool Attention) demands near-immediate engineering action on the planning-bridge MCP server.

\subsubsection{\arxivid{2604.21816} --- Tool Attention: Dynamic Tool Gating + Lazy Schema Loading}

\paperhead{Themes:}{mcp-tax, tool-gating, lazy-schema, context-fracture-threshold.}

\textbf{Key idea.} Quantifies the ``MCP Tax / Tools Tax'' at \textbf{10K--60K tokens per turn} in real deployments --- past the published context-fracture point of ${\sim}70\%$. Solution combines: (i) Intent Schema Overlap (ISO) score from sentence embeddings, (ii) state-aware gating function enforcing preconditions and access scopes, (iii) two-phase lazy schema loader that keeps a compact summary pool in context and promotes full JSON schemas only for top-k gated tools.

\textbf{GSD application.} Essentially mandatory reading for the planning-bridge MCP design. The lazy-schema-loading pattern is a near-immediate token-budget win to deploy in v1.50. The 70\% context-fracture threshold is a useful operational constant for the bounded-tape policy --- when total tool schema content approaches 70\% of context, gate aggressively. Engineering ticket should target Q2 2026. \textbf{Adoption: ADOPT (technique + engineering ticket).} Target: planning-bridge MCP server, v1.50 milestone.

\subsubsection{\arxivid{2604.21282} --- Strategic Heterogeneous Multi-Agent (3+1)}

\paperhead{Themes:}{chipset, three-plus-one, cooperative-and-adversarial-games.}

\textbf{Key idea.} Three cloud experts (DeepSeek-V3) running in parallel with complementary perspectives + one local lightweight verifier (Qwen3-8B) doing adversarial validation at zero marginal cost. Two-game framework: \textit{cooperative} game among experts (super-additive value from diverse perspectives) + \textit{adversarial verification} game (quality-assurance incentives). 77.2\% F1 at \$0.002/sample on Juliet Test Suite (262 real samples, 14 CWE types).

\textbf{GSD application.} This is empirically the Opus(${\times}N$ parallel) + Haiku(local verifier) pattern from the GSD chipset architecture, validated with cost numbers. \textbf{Cite as external validation of the chipset model-tier mapping.} The cooperative + adversarial two-game framing is also a clean formal model worth adopting in GSD docs to explain why three Opus experts + one Haiku verifier outperforms a single Opus per task. \textbf{Adoption: ADOPT (citation + formal model).} Target: chipset architecture vision doc, v1.50 retrospective.

\subsubsection{\arxivid{2603.22823} --- Empirical Comparison of Agent Communication Protocols}

\paperhead{Themes:}{dacp-validation, comparative-benchmark, five-axis-evaluation.}

\textbf{Key idea.} Pilot benchmark comparing tool integration / multi-agent delegation / hybrid architectures for standardized queries at three complexity levels. Quantifies advantages and disadvantages on five axes: \textbf{response time, context-window consumption, cost, error recovery, implementation complexity}.

\textbf{GSD application.} Direct methodology to validate DACP. Adopt the five-axis evaluation matrix as the standard report card for any DACP-related milestone. This is the citation that contextualizes DACP work in the broader literature --- before DACP, GSD could only assert it solves a problem; with this paper as comparator, DACP can be situated. \textbf{Adoption: ADOPT (methodology).} Target: DACP comparative-benchmark plan, v1.50 retrospective.

\subsubsection{\arxivid{2604.21794} --- DiffMAS: End-to-End Latent Multi-Agent Communication}

\paperhead{Themes:}{dacp-comparison, latent-communication, kv-cache-as-channel.}

\textbf{Key idea.} Treats inter-agent communication as a \textit{learnable} component rather than a fixed text-based interface. Latent communication via internal KV-cache representations, jointly optimized with multi-agent reasoning via parameter-efficient supervised training. 26.7\% AIME24, 20.2\% GPQA-Diamond.

\textbf{GSD application.} Important reference for DACP's five fidelity levels. Latent KV-cache communication occupies the highest-fidelity end of the spectrum that DACP reaches via structured-text bundles. Worth a dedicated comparative experiment: structured-text DACP vs learned-latent DACP, with cost and correctness tradeoffs documented. The DiffMAS approach is harder to operationalize across model providers but offers a ceiling on what DACP could achieve. \textbf{Adoption: TRACK (comparison ceiling).} Target: DACP fidelity-level documentation.

\subsubsection{\arxivid{2511.00413} --- Tree Training: Shared Prefix Reuse for Agentic LLMs}

\paperhead{Themes:}{subagent, tree-trajectories, training-time-optimization.}

\textbf{Key idea.} Multi-turn agentic trajectories with concurrent tool use, sub-agents, think-mode, and context management form \textit{trees with shared prefixes}, not linear sequences. Authors prove averaging loss over branches independently is algebraically identical to a per-token weighted loss where the weight equals the fraction of branches passing through that token. DFS serialization visits each token exactly once; full-attention and SSM layers are adapted so log-probabilities match independent per-branch calculation exactly.

\textbf{GSD application.} The formal training-time treatment of GSD's Wave/CAPCOM/Squadron architecture. The tree-trajectory framing is exactly the structure GSD produces --- if the Silicon Layer ever does in-house RL or SFT training over GSD execution traces, this paper's serialization is the right starting point. The algebraic identity between per-branch and per-token-weighted loss is the formal justification for a shared-prefix training pipeline. \textbf{Adoption: ADOPT (citation + future work).} Target: Silicon Layer training pipeline future-work notes, v1.50 retrospective.

\subsubsection{\arxivid{2604.21255} --- Distillation-Induced Similarity in Tool-Use (RPS / AGS)}

\paperhead{Themes:}{chipset, behavior-metrics, distillation-detection.}

\textbf{Key idea.} Two complementary metrics for isolating non-mandatory behavioral patterns: \textbf{Response Pattern Similarity (RPS)} for verbal alignment and \textbf{Action Graph Similarity (AGS)} for tool-use habits modeled as directed graphs. Within-family model pairs score 5.9pp higher AGS than cross-family. Kimi-K2 (thinking) reaches 82.6\% / 94.7\% on Claude Sonnet 4.5 (thinking).

\textbf{GSD application.} Direct evaluation tooling for the chipset architecture (Agnus/Denise/Paula/Gary $\to$ Opus/Sonnet/Haiku). Use AGS to validate that Sonnet-implementations don't unintentionally diverge from Opus-designed patterns when running the same skill. Use RPS to detect when a Haiku-tier scaffold is ``thinking like'' something it shouldn't be. Add to the chipset test harness as the formal behavior-similarity metric. \textbf{Adoption: ADOPT (metric).} Target: chipset architecture validation, College-of-Knowledge dispatch monitoring.

\subsubsection{\arxivid{2604.21444} --- HiCrew: Hierarchical Multi-Agent for Long-Form Video}

\paperhead{Themes:}{multi-agent, hierarchical, dynamic-role-selection.}

\textbf{Key idea.} Hierarchical multi-agent framework with three contributions: \textbf{Hybrid Tree} (preserves temporal topology while doing relevance-guided clustering), \textbf{Question-Aware Captioning} (intent-driven visual prompts), and a \textbf{Planning Layer} (dynamically selects roles and execution paths based on question complexity).

\textbf{GSD application.} The dynamic role-selection pattern is exactly Squadron-vs-Fleet activation profile selection. The Hybrid Tree pattern is reusable for College-of-Knowledge departmental search where preserving the original document hierarchy matters. \textbf{Adoption: ADOPT (pattern reference).} Target: activation-profile selection logic, College-of-Knowledge search.

\subsubsection{\arxivid{2508.21720} --- PosterForest: Hierarchical Multi-Agent Posters}

\paperhead{Themes:}{multi-agent, intermediate-tree-representation, recursive-refinement.}

\textbf{Key idea.} Poster Tree intermediate representation captures document hierarchy and visual-textual semantics. Content and layout agents perform hierarchical reasoning and recursive refinement, optimizing globally first then locally. Training-free framework outperforms prior methods on automatic and human evaluations.

\textbf{GSD application.} The intermediate-tree-then-refine pattern is directly applicable to the research-mission-generator skill's three-stage Vision $\to$ Research $\to$ Mission pipeline. Worth restructuring the mission-package generator around an explicit \textit{Mission Tree} intermediate representation rather than the current linear narrative. \textbf{Adoption: ADOPT (architecture upgrade).} Target: research-mission-generator skill v2 design.

\subsubsection{\arxivid{2604.20848} --- MATRAG: Multi-Agent Transparent RAG}

\paperhead{Themes:}{multi-agent-rag, four-agent, transparency-scoring.}

\textbf{Key idea.} Four specialized agents: \textbf{User Modeling} (dynamic preference profiles), \textbf{Item Analysis} (semantic features from KGs), \textbf{Reasoning} (synthesizes signals), \textbf{Explanation} (natural-language justifications grounded in retrieved knowledge). Includes a \textbf{transparency scoring mechanism} that quantifies explanation faithfulness and relevance.

\textbf{GSD application.} The four-agent + transparency-score pattern maps cleanly onto a Mission Crew Manifest upgrade --- adopt transparency scoring as a numerical metric on every CAPCOM gate. Each gate decision should record a transparency score for the artefact passing through. \textbf{Adoption: ADOPT (metric + crew pattern).} Target: Mission Crew Manifest upgrade, CAPCOM gate audit-log schema.

\subsubsection{\arxivid{2510.26615} --- SARA / SlideAgent: Hybrid RAG with Compression Vectors}

\paperhead{Themes:}{rag, hybrid-text-and-vector, iterative-reranking.}

\textbf{Key idea.} Hybrid RAG combining a small text-snippet pool (preserves entities and numerics) with semantic compression vectors (broader coverage). Vectors used for iterative evidence reranking. Across 9 datasets and 5 LLMs: +17.71 answer relevance, +13.72 correctness, +15.53 semantic similarity.

\textbf{GSD application.} Architectural template for the College of Knowledge departmental retrieval layer. The text-pool-for-entities + vector-pool-for-coverage split solves the precision/recall tension in current single-mode retrieval. \textbf{Adoption: ADOPT (architecture).} Target: College of Knowledge retrieval layer, departmental search.

\subsubsection{\arxivid{2604.21590} --- AgenticQwen: Dual Data Flywheels}

\paperhead{Themes:}{agentic-rl, dual-flywheel, retrospective-discipline.}

\textbf{Key idea.} Two flywheels for small agentic-model training: \textbf{reasoning flywheel} (errors $\to$ harder tasks via curriculum) and \textbf{agentic flywheel} (linear workflows $\to$ multi-branch behavior trees). Combined with multi-round RL on synthetic + open-source data.

\textbf{GSD application.} The dual-flywheel pattern is the abstract form of GSD's milestone retrospective + feed-forward discipline. Use as a reference architecture in v1.50 retrospective Half~A's discussion of the retrospective layer's evolution from milestones 1--49. \textbf{Adoption: ADOPT (citation).} Target: v1.50 retrospective Half A.

% =============================================================
% MODULE 5 -- VERIFICATION, PROOF & TEST HARNESS
% =============================================================
\subsection{Module 5 --- Verification, Proof \& Test Harness}

Eight papers covering the verification territory most directly aligned with The Space Between's proof companion target and the ZFC compliance auditor.

\subsubsection{\arxivid{2604.21598} --- DryRUN: LLM-Driven Code Generation w/o Public Tests}

\paperhead{Themes:}{test-generation, self-correcting, overconfidence-gap.}

\textbf{Key idea.} Identifies an ``overconfidence gap'' from public-test reliance: frameworks overfit to simple examples and fail on hidden evaluations. Demonstrates LLMs can \textit{autonomously generate valid inputs and simulate execution traces to self-correct} without human-provided test cases.

\textbf{GSD application.} Direct technique for the skill-validation test harness. Eliminates the bottleneck of hand-authoring tests for every new skill --- the harness can synthesize them. Particularly relevant for v1.50's per-chapter test generation in the proof companion (Module 5 cross-link to v1.50 Half B). \textbf{Adoption: ADOPT (technique).} Target: skill validation test harness, v1.50 proof companion test infrastructure.

\subsubsection{\arxivid{2511.23159} --- AI for Software Engineering: From Probable to Provable}

\paperhead{Themes:}{philosophical-position, formal-verification, probable-vs-provable.}

\textbf{Key idea.} Position paper: combine the creativity of AI techniques for programming with the rigor of formal specification methods and program verification. Argues that AI-assisted coding faces two overwhelming obstacles --- specification difficulty and hallucination --- and that formal methods provide the answer.

\textbf{GSD application.} Cite as philosophical alignment with The Space Between's proof-companion work and the ZFC compliance auditor goals. Strong companion citation to GROUNDINGmd (Module 1) and the Specification Trap (Module 3). \textbf{Adoption: ADOPT (citation).} Target: The Space Between, ZFC auditor design doc.

\subsubsection{\arxivid{2604.21570} --- SpecSyn: LLM-Based Formal Specification Synthesis}

\paperhead{Themes:}{formal-spec, decomposition-then-iterate, real-world-programs.}

\textbf{Key idea.} LLM-based specification generation that decomposes input programs into segments handled by an iterative specification-generation process, with mechanisms to evaluate and guarantee specification strength. Targets real-world program verification.

\textbf{GSD application.} The decomposition-then-iterate pattern is reusable for skill-validation. If GSD ever adopts a Lean-style proof-obligation system for skill correctness, SpecSyn's pipeline structure is a strong starting point. \textbf{Adoption: TRACK (technique reference).} Target: future ZFC auditor design.

\subsubsection{\arxivid{2604.18792} --- Tractable Verification of DSLTrans (Cutoff Theorem)}

\paperhead{Themes:}{bounded-model-checking, cutoff-theorem, DSL-verification.}

\textbf{Key idea.} Cutoff Theorem proving bounded model checking is \textit{complete} for a precise DSLTrans fragment and positive existence/traceability properties --- turning infinite search into a finite computable bound. Composable soundness-preserving optimizations + Z3-based implementation. On 29 transformations and 899 properties: 552 proved, 345 counterexamples produced.

\textbf{GSD application.} Reference for the ZFC compliance auditor work. The cutoff-theorem-style framing (turn infinite search into finite computable bound) is the right mental model for skill-correctness invariants where the search space appears unbounded but in fact admits a finite witness. \textbf{Adoption: ADOPT (citation + framing).} Target: ZFC auditor design doc, v1.50 proof companion.

\subsubsection{\arxivid{2604.21187} --- Doubly Saturated Ramsey Graphs (SAT + LLM + Lean)}

\paperhead{Themes:}{hybrid-sat-llm-lean, mathematical-discovery, formal-correctness.}

\textbf{Key idea.} Combines SAT solving with bespoke LLM-generated code to discover infinite families of doubly saturated Ramsey-good graphs --- answering an open question of Grinstead and Roberts from 1982. Then uses LLMs to generate and formalize correctness proofs in Lean. Argues tool-driven workflows (automated reasoning + LLMs + formal verification) will play an increasingly central role in experimental mathematics.

\textbf{GSD application.} Direct precedent for the v1.50 proof-companion's hybrid SAT/Lean/LLM workflow. Worth citing in The Space Between as evidence that the methodology is producing genuine mathematical results --- a case where the pattern Tibsfox is building toward is already producing publishable results elsewhere. Strong cross-link with the Lean ecosystem. \textbf{Adoption: ADOPT (precedent + citation).} Target: The Space Between proof companion, v1.50 Half B design.

\subsubsection{\arxivid{2603.01170} --- ATLAS: Threat-to-Assertion Learning for SoC Security}

\paperhead{Themes:}{cwe-driven, knowledge-base-to-assertion, hardware-verification.}

\textbf{Key idea.} LLM-driven framework bridging standardized threat modeling and property-based formal verification for SoC security. Uses CWE as the vulnerability knowledge base, identifies SoC-specific assets, generates assertion-based security properties + JasperGold scripts. Detects 39/48 CWEs and generates correct properties for 33.

\textbf{GSD application.} Pattern reference for the planning-bridge MCP security-test generator (cross-link to Module 3). Specific to hardware verification but the CWE-to-assertion pipeline structure generalizes to any well-classified threat catalog. \textbf{Adoption: TRACK (pattern reference).} Target: planning-bridge security test generator design.

\subsubsection{\arxivid{2604.14709} --- HWE-Bench: LLM Agents on Real Hardware Bug Repair}

\paperhead{Themes:}{repository-level-benchmark, hardware-bug-repair, verilog-chisel.}

\textbf{Key idea.} 417 task instances from real historical bug-fix pull requests across six major open-source hardware projects (Verilog/SystemVerilog and Chisel; RISC-V cores, SoCs, security roots-of-trust). Best agent: 70.7\% overall, $>$90\% on small cores, $<$65\% on complex SoCs.

\textbf{GSD application.} Highly relevant given Tibsfox's Bilocon/Fluke/Phillips Semiconductor history. If GSD ever produces an FPGA/ASIC-flow skill pack, this is the benchmark to target. Also a candidate inclusion in the NASA Mission Series as a ``hardware mission'' template, leveraging the project's Wanapum / Hanford / Boeing-region hardware-engineering heritage. \textbf{Adoption: TRACK (future skill-pack target).} Target: NASA Mission Series hardware-mission template.

\subsubsection{\arxivid{2604.21505} --- Orchid: Requirement Ambiguity in Code Generation}

\paperhead{Themes:}{ambiguity-taxonomy, four-types, advanced-models-suffer-more.}

\textbf{Key idea.} 1,304 function-level tasks across four ambiguity types: \textbf{lexical, syntactic, semantic, vagueness}. Demonstrates ambiguity consistently degrades all evaluated LLMs, with the most pronounced negative effects on the most advanced models.

\textbf{GSD application.} Important for skill-specification quality. Adopt the four-type ambiguity taxonomy as a checklist for SKILL.md authoring. The empirical finding that \textit{advanced models suffer more} from ambiguity argues for stricter spec discipline in Opus-tier work, not looser. Build into the skill-creator linter alongside the Five-Principle structural-completeness check from Module 3. \textbf{Adoption: ADOPT (linter check).} Target: skill-creator linter, SKILL.md authoring template.

% =============================================================
% MODULE 6 -- MEMORY, SKILLS ECONOMY, CULTURAL SENSITIVITY & CURRICULUM
% =============================================================
\subsection{Module 6 --- Memory, Skills Economy, Cultural Sensitivity \& Curriculum}

Thirteen papers covering the broadest spread of GSD concerns. The largest module, organized by sub-theme: memory/context (2 papers), skills economy (1), cultural sensitivity (5), curriculum/pedagogy (3), long-horizon agents (2).

\subsubsection{\arxivid{2604.20920} --- Forget, Then Recall: Gist Sparse Attention}

\paperhead{Themes:}{context-mgmt, gist-tokens-as-routing, hierarchical-recursive.}

\textbf{Key idea.} End-to-end-learnable bridge between KV-cache selection and compression: interleaved \textit{gist} tokens act as both summaries \textit{and} routing signals for sparse attention. Coarse-to-fine: compress $\to$ select $\to$ restore raw chunks for detailed attention. Hierarchical recursive gist-of-gist construction yields multi-resolution access with logarithmic per-step decoding complexity.

\textbf{GSD application.} Architectural template for ``compressed memory with selective unfolding'' in fresh-context subagent design. The recursive gist-of-gist pattern is a clean way to handle the College of Knowledge departments' nested structure --- each department is a gist, each subtopic is a gist-of-gist, and unfolding happens only on demand. \textbf{Adoption: TRACK (architecture reference).} Target: College of Knowledge nested-retrieval design.

\subsubsection{\arxivid{2604.20487} --- Knowledge Capsules: Structured Nonparametric Memory}

\paperhead{Themes:}{kv-injection, nonparametric-memory, attention-compatible.}

\textbf{Key idea.} Knowledge stored as \textit{capsules} compiled into attention-compatible KV representations and injected, rather than appended as text tokens that compete in attention. Outperforms RAG and GraphRAG on multi-hop reasoning benchmarks with improved stability in long-context cases.

\textbf{GSD application.} A more principled alternative to text-based RAG for the College of Knowledge. If the Silicon Layer ever supports KV-injection custom kernels, capsule storage becomes a viable upgrade for departmental knowledge retrieval. Currently a future-work item; the kernel infrastructure does not yet exist in the GSD stack. \textbf{Adoption: TRACK (future-work).} Target: College of Knowledge v2 retrieval design.

\subsubsection{\arxivid{2604.21829} --- Black-Box Skill Stealing Attack on Proprietary Agents}

\paperhead{Themes:}{skill-marketplace, attack-taxonomy, prompt-stealing.}

\textbf{Key idea.} \textbf{First empirical study} of black-box skill-stealing attacks against LLM agent systems. Notes existing free skill marketplaces report 90,368 published skills; paid marketplaces show 2,000+ listings and \$100K+ creator earnings. Attack pipeline: model-generated seed prompts $\to$ scenario rationalization + structure injection $\to$ embedding-filter diversity. Evaluated on 3 commercial agent architectures and 5 LLMs.

\textbf{GSD application.} \textbf{Mandatory reading} for the DoltHub federated skill-economy work. Adopt the attack taxonomy as the explicit threat model for the federated skill-economy design. The 90K-skill marketplace number is also useful market sizing for the Wasteland integration roadmap. Engineering action: the DoltHub contribution roadmap should include a threat-model document referencing this paper before first contribution. \textbf{Adoption: ADOPT (threat model).} Target: DoltHub threat-model doc, Wasteland integration design.

\subsubsection{\arxivid{2604.12710} --- LASA: Language-Agnostic Semantic Alignment for Safety}

\paperhead{Themes:}{cultural-sensitivity, semantic-bottleneck, low-resource-lang.}

\textbf{Key idea.} Identifies a \textit{semantic bottleneck} intermediate layer where representation geometry is governed primarily by shared semantic content rather than language identity. Anchors safety alignment there. Average attack success rate drops from 24.7\% to 2.8\% on LLaMA-3.1-8B-Instruct; ${\sim}3\%$--$4\%$ across Qwen2.5 and Qwen3 Instruct (7B--32B).

\textbf{GSD application.} The semantic-bottleneck framing is exactly what's needed to prevent OCAP/UNDRIP enforcement from being defeated by translation or paraphrase. If a sacred-knowledge attribution rule is anchored at the semantic-bottleneck layer rather than the surface text, paraphrase attacks fail. Worth a focused experiment on the Indigenous-knowledge attribution pipeline. \textbf{Adoption: TRACK (experiment).} Target: cultural-sensitivity test harness experiment.

\subsubsection{\arxivid{2604.21751} --- Why are LLMs Obsessed with Japanese Culture}

\paperhead{Themes:}{cultural-bias, regional-preferences, sft-emergent-bias.}

\textbf{Key idea.} CROQ taxonomy (Culture-Related Open Questions) reveals LLMs show a \textit{clear tendency toward Japan} (contrary to prior Anglocentric framing). When prompting in English/high-resource languages, LLMs show less inclination toward countries for which the input language is official. Bias appears \textit{after SFT, not during pre-training}.

\textbf{GSD application.} Important precedent for the cultural-sensitivity test harness. Bias tests need to cover not just Anglocentric bias but unexpected concentrations like Japan-preference. Add Japan-vs-other-Pacific-Rim comparisons to the cultural test suite given Tibsfox's Pacific-Rim focus and the FoxFiber BC--Chile corridor work. The SFT-emergent-not-pretrain finding constrains where mitigation can apply. \textbf{Adoption: ADOPT (test additions).} Target: cultural-sensitivity test suite.

\subsubsection{\arxivid{2604.20996} --- AFRILANGTUTOR: Low-Resource Language Tutoring}

\paperhead{Themes:}{low-resource-lang, dictionary-seed-pipeline, sft-dpo.}

\textbf{Key idea.} 194.7K African language $\leftrightarrow$ English dictionary entries used as seed resources to construct large-scale, diverse, verifiable student-tutor QA interactions. 78.9K multi-turn training examples. Combined SFT + DPO yields 1.8\%--15.5\% gains on LLaMA-3-8B-IT and Gemma-3-12B-IT across 10 African languages.

\textbf{GSD application.} Methodological template for the FoxFire Heritage Skills mission --- same dictionary-seed-to-training-data pipeline, applied to PNW Indigenous languages with appropriate sovereignty controls (OCAP/CARE/UNDRIP). Note: any PNW-language adaptation must consult specific nations (Coast Salish, Lummi/Lhaq'temish, Tulalip, Stillaguamish, Wanapum) before data construction begins; this paper provides the technique, not the permission. \textbf{Adoption: TRACK (methodological template).} Target: FoxFire Heritage Skills mission, sovereignty-gated.

\subsubsection{\arxivid{2604.21352} --- CARE: Counselor-Aligned Response Engine}

\paperhead{Themes:}{mental-health-safety, low-resource-lang, counselor-alignment.}

\textbf{Key idea.} GenAI framework that assists counselors with real-time psychologically-grounded responses. Designed for low-resource language settings. Counselor-alignment pattern (assist rather than replace).

\textbf{GSD application.} Reference for the College-of-Knowledge mind-body department's safety design when the user expresses distress. The counselor-alignment pattern (assist rather than replace) is the right framing for any GSD work that touches mental-health adjacencies. Aligns with Claude's existing user-wellbeing constraints. \textbf{Adoption: TRACK (design pattern).} Target: College of Knowledge mind-body department safety design.

\subsubsection{\arxivid{2604.20904} --- Privacy Reasoning via Normative Simulacra from Fiction}

\paperhead{Themes:}{contextual-integrity, fiction-derived-norms, grpo-rl.}

\textbf{Key idea.} Extracts \textit{normative simulacra} (structured representations of norms and information flows) from fiction novels and uses them to fine-tune LLMs via SFT + GRPO. Composite reward includes task clarity (subsuming schema validity, construct discrimination, extraction confidence).

\textbf{GSD application.} A genuinely interesting fit for The Hundred Voices --- fiction-derived normative structures are \textit{exactly} what that work catalogs. Possible cross-pollination opportunity: train a privacy-reasoning component using The Hundred Voices' authorial-voice taxonomy as a normative-simulacrum source. This would be a natural cross-project collaboration between The Hundred Voices and the FoxData privacy work. \textbf{Adoption: TRACK (cross-pollination).} Target: The Hundred Voices $\leftrightarrow$ FoxData experiment.

\subsubsection{\arxivid{2604.21308} --- CI-Work: Contextual Integrity in Enterprise LLM Agents}

\paperhead{Themes:}{contextual-integrity, enterprise-deployment, utility-privacy-tradeoff.}

\textbf{Key idea.} CI-grounded benchmark simulating enterprise workflows across five information-flow directions. \textbf{Privacy violation rates 15.8\%--50.9\%}, with leakage up to 26.7\%. Counterintuitive: \textit{higher task utility correlates with more privacy violations}. The agent that does the most useful work also leaks most.

\textbf{GSD application.} Critical operational metric for any FoxFiber/FoxData enterprise deployment. The utility/privacy frontier needs to be a \textit{measured trade-off} in the Safety Warden's reporting, not an implicit choice. Adopt the five information-flow directions as the enterprise threat-model schema. \textbf{Adoption: ADOPT (metric + threat model).} Target: enterprise deployment Safety Warden, FoxFiber/FoxData design.

\subsubsection{\arxivid{2510.18731} --- RLAAR: Curriculum RL with Verifiable Accuracy + Abstention}

\paperhead{Themes:}{curriculum, abstention-as-reward, lost-in-conversation.}

\textbf{Key idea.} \textbf{Competence-gated curriculum} incrementally increases dialogue difficulty (instruction shards). Mixed-reward system teaches abstention when problems are unsolvable. Lost-in-Conversation decay 62.6\% $\rightarrow$ 75.1\%; calibrated abstention 33.5\% $\rightarrow$ 73.4\%.

\textbf{GSD application.} Direct adoption candidate for the Teacher (Opus) / Student (Sonnet) / Support (Haiku) curriculum. The \textit{abstention-as-reward} pattern is the missing piece for ``the Student should know when to escalate to Teacher'' --- currently a soft norm in GSD, this paper provides the formal training mechanism. \textbf{Adoption: ADOPT (training mechanism).} Target: pedagogical Teacher/Student/Support design, v1.50 Half B chapter validation.

\subsubsection{\arxivid{2601.10003} --- SocraticKG: 5W1H QA-Driven KG Construction}

\paperhead{Themes:}{knowledge-graph, socratic-decomposition, 5w1h.}

\textbf{Key idea.} Introduces QA pairs as a structured intermediate representation to systematically unfold document-level semantics prior to triple extraction. \textbf{5W1H-guided QA expansion} captures contextual dependencies and implicit relational links that direct KG extraction misses.

\textbf{GSD application.} Pattern directly applicable to College of Knowledge departmental knowledge-base construction. The 5W1H + Socratic decomposition matches the GSD habit of structured-question scaffolding (the same shape appears throughout the existing skill set). \textbf{Adoption: ADOPT (pattern).} Target: College of Knowledge KB construction, mission-package research stage.

\subsubsection{\arxivid{2604.21380} --- IRAP: Interactive RAG-Based Preference Elicitation}

\paperhead{Themes:}{elicitation, interactive-rag, bounded-rounds.}

\textbf{Key idea.} Quantifies vague performance requirements into mathematical functions via interactive elicitation. Up to 40$\times$ improvement with as few as 5 interaction rounds.

\textbf{GSD application.} Direct adoption candidate for the vision-to-mission skill's elicitation phase. The 5-round bound is an attractive operational target for CAPCOM-gate dialogue length --- if the gate cannot resolve in 5 rounds, escalate to a different model tier. \textbf{Adoption: ADOPT (technique + bound).} Target: vision-to-mission elicitation, CAPCOM gate dialogue protocol.

\subsubsection{\arxivid{2601.07262} --- ColorBrowserAgent: Adaptive Knowledge Evolution}

\paperhead{Themes:}{long-horizon-agent, hitl-knowledge-adaptation, progressive-summarization.}

\textbf{Key idea.} Knowledge-evolving agent for robust web automation. Two synergistic mechanisms: \textbf{HITL knowledge adaptation} (transforms sparse human feedback into reusable domain knowledge) and \textbf{knowledge-aligned progressive summarization} (stabilizes long interactions through memory compression). 71.2\% on WebArena, +19.3\% commercial user-satisfaction.

\textbf{GSD application.} The HITL-feedback-to-domain-knowledge pattern is the operational pipeline GSD's CAPCOM gates need on the back end --- capture critique outcomes into structured knowledge updates rather than session-local fixes. Cross-link with CHAI (Module 3) and AEL (Module 2). \textbf{Adoption: ADOPT (pipeline pattern).} Target: CAPCOM gate $\to$ bounded-learning capture pipeline.

\subsubsection{\arxivid{2604.20855} --- Caesar: Deep Agentic Web Exploration for Creative Synthesis}

\paperhead{Themes:}{long-horizon-agent, kg-driven, adversarial-draft-refinement.}

\textbf{Key idea.} Bridges information gathering and synthesis of new insights via two components: (1) exploration driven by a dynamic context-aware policy, (2) synthesis controlled by an \textit{adversarial draft refinement loop} that actively seeks novel perspectives rather than confirming established priors. Knowledge-graph-driven associative reasoning enables non-obvious connections between disparate concepts.

\textbf{GSD application.} Direct enhancement candidate for the research-mission-generator skill --- the adversarial-draft-refinement pattern is what's currently missing in the research stage. Worth a focused upgrade in the skill's v2 design. Cross-link with PosterForest (Module 4) for the intermediate-tree-then-refine pattern. \textbf{Adoption: ADOPT (technique).} Target: research-mission-generator v2.

% =============================================================
% SAFETY AND SENSITIVITY CONSIDERATIONS
% =============================================================
\subsection{Safety and Sensitivity Considerations}

Although this is a literature-review mission rather than a deployment, several specific sensitivities apply.

\begin{center}
\rowcolors{2}{rowgray}{white}
\begin{tabularx}{\textwidth}{L{4cm}XC{2.5cm}}
\toprule
\rowcolor{primary}
\tableheader{Concern} & \tableheader{Constraint} & \tableheader{Type} \\
\midrule
Indigenous attribution (M6) & Any reference to PNW Indigenous languages or knowledge in the AFRILANGTUTOR-template adaptation must name specific nations (Coast Salish, Lummi/Lhaq'temish, Tulalip, Stillaguamish, Wanapum). Never generic. & ABSOLUTE \\
Sovereignty gating (M6) & FoxFire Heritage application of low-resource techniques requires nation-specific consultation BEFORE data construction begins. The technique is permission-neutral. & ABSOLUTE \\
Mental-health safety (M6) & CARE-template adoption respects Claude's existing user-wellbeing constraints. Counselor-alignment means assist, never replace. No diagnostic or therapeutic claims in adoption documentation. & ABSOLUTE \\
Skill-marketplace threat model (M6) & Black-Box Skill Stealing attack taxonomy must be incorporated into the DoltHub contribution roadmap before first federated-skill contribution. & GATE \\
MCP security (M3, M4) & Tool Attention's lazy-schema-loading + MCP Pitfall Lab's hardening guidance + Function Hijacking's threat model are jointly required for any planning-bridge MCP design change. & GATE \\
Source quality & All 54 papers cited from arXiv. Direct DOIs / paper PDFs preserved. No popular-press paraphrases. & ABSOLUTE \\
No paper paraphrase $>$15 words & Stage 2's per-paper write-ups paraphrase entirely; never reproduce abstract sentences verbatim. & ABSOLUTE \\
Cross-model behavior similarity (M4) & RPS / AGS metrics adopted only for internal monitoring of GSD's chipset; not for asserting any model is ``identical to'' another. & ANNOTATE \\
\bottomrule
\end{tabularx}
\end{center}

% =============================================================
% SOURCE BIBLIOGRAPHY
% =============================================================
\subsection{Source Bibliography}

All sources are arXiv preprints from the April 2026 listing (\texttt{Computer\_Science25.html}, 930 papers). Resolved IDs, ordered by module then by adoption priority within module.

\subsubsection{Module 1 --- Architectural Foundations \& Convergent Discovery}

\begin{enumerate}[leftmargin=*]
\item \arxivid{2604.21910} --- From Research Question to Scientific Workflow: Leveraging Agentic AI for Science Automation.
\item \arxivid{2604.21744} --- Agentic AI-assisted coding offers a unique opportunity to instill epistemic grounding during software development (GROUNDINGmd).
\item \arxivid{2604.20874} --- The Root Theorem of Context Engineering.
\item \arxivid{2604.21003} --- The Last Harness You'll Ever Build.
\end{enumerate}

\subsubsection{Module 2 --- Bounded Learning \& Silicon Layer LoRA}

\begin{enumerate}[leftmargin=*]
\item \arxivid{2604.21725} --- AEL: Agent Evolving Learning for Open-Ended Environments.
\item \arxivid{2511.11439} --- Retrofit: Continual Learning with Controlled Forgetting for Binary Security Detection.
\item \arxivid{2604.20300} --- FSFM: A Biologically-Inspired Framework for Selective Forgetting of Agent Memory.
\item \arxivid{2604.20943} --- SCM: Sleep-Consolidated Memory with Algorithmic Forgetting for LLMs.
\item \arxivid{2604.21927} --- Fine-Tuning Regimes Define Distinct Continual Learning Problems.
\item \arxivid{2604.21571} --- Separable Expert Architecture: Composable Adapters and Deletable User Proxies.
\item \arxivid{2604.21905} --- Low-Rank Adaptation Redux for Large Models.
\item \arxivid{2604.21901} --- GiVA: Gradient-Informed Bases for Vector-Based Adaptation.
\item \arxivid{2509.18629} --- HyperAdapt: Simple High-Rank Adaptation.
\item \arxivid{2604.21330} --- Teacher-Guided Routing for Sparse Vision Mixture-of-Experts.
\end{enumerate}

\subsubsection{Module 3 --- Safety Warden, CAPCOM \& MCP Security}

\begin{enumerate}[leftmargin=*]
\item \arxivid{2601.18491} --- AgentDoG: A Diagnostic Guardrail Framework for AI Agent Safety and Security.
\item \arxivid{2604.20911} --- Omission Constraints Decay While Commission Constraints Persist in Long-Context LLM Agents.
\item \arxivid{2604.21477} --- MCP Pitfall Lab: Exposing Developer Pitfalls in MCP Tool Server Security under Multi-Vector Attacks.
\item \arxivid{2604.20994} --- Breaking MCP with Function Hijacking Attacks: Novel Threats for Function Calling and Agentic Models.
\item \arxivid{2604.21131} --- Cross-Session Threats in AI Agents: Benchmark, Evaluation, and Algorithms.
\item \arxivid{2604.21090} --- Structural Quality Gaps in Practitioner AI Governance Prompts.
\item \arxivid{2604.21718} --- Building a Precise Video Language with Human-AI Oversight (CHAI).
\item \arxivid{2512.03048} --- The Specification Trap: Why Static Value Alignment Alone Is Insufficient for Robust Alignment.
\end{enumerate}

\subsubsection{Module 4 --- Multi-Agent, DACP \& Chipset Architecture}

\begin{enumerate}[leftmargin=*]
\item \arxivid{2604.21816} --- Tool Attention Is All You Need: Dynamic Tool Gating and Lazy Schema Loading.
\item \arxivid{2604.21282} --- Strategic Heterogeneous Multi-Agent Architecture for Cost-Effective Code Vulnerability Detection.
\item \arxivid{2603.22823} --- Empirical Comparison of Agent Communication Protocols for Task Orchestration.
\item \arxivid{2604.21794} --- Learning to Communicate: Toward End-to-End Optimization of Multi-Agent Language Systems (DiffMAS).
\item \arxivid{2511.00413} --- Tree Training: Accelerating Agentic LLMs Training via Shared Prefix Reuse.
\item \arxivid{2604.21255} --- When Agents Look the Same: Quantifying Distillation-Induced Similarity in Tool-Use Behaviors.
\item \arxivid{2604.21444} --- HiCrew: Hierarchical Reasoning for Long-Form Video Understanding.
\item \arxivid{2508.21720} --- PosterForest: Hierarchical Multi-Agent Collaboration for Scientific Poster Generation.
\item \arxivid{2604.20848} --- MATRAG: Multi-Agent Transparent Retrieval-Augmented Generation.
\item \arxivid{2510.26615} --- SARA / SlideAgent: Hybrid RAG with Compression Vectors.
\item \arxivid{2604.21590} --- AgenticQwen: Training Small Agentic Language Models with Dual Data Flywheels.
\end{enumerate}

\subsubsection{Module 5 --- Verification, Proof \& Test Harness}

\begin{enumerate}[leftmargin=*]
\item \arxivid{2604.21598} --- DryRUN: On the Role of Public Tests in LLM-Driven Code Generation.
\item \arxivid{2511.23159} --- AI for Software Engineering: From Probable to Provable.
\item \arxivid{2604.21570} --- SpecSyn: LLM-based Synthesis and Refinement of Formal Specifications.
\item \arxivid{2604.18792} --- Tractable Verification of Model Transformations: A Cutoff-Theorem Approach for DSLTrans.
\item \arxivid{2604.21187} --- Doubly Saturated Ramsey Graphs: A Case Study in Computer-Assisted Mathematical Discovery.
\item \arxivid{2603.01170} --- ATLAS: AI-Assisted Threat-to-Assertion Learning for SoC Security Verification.
\item \arxivid{2604.14709} --- HWE-Bench: Benchmarking LLM Agents on Real-World Hardware Bug Repair Tasks.
\item \arxivid{2604.21505} --- Assessing the Impact of Requirement Ambiguity on LLM-based Function-Level Code Generation (Orchid).
\end{enumerate}

\subsubsection{Module 6 --- Memory, Skills Economy, Cultural Sensitivity \& Curriculum}

\begin{enumerate}[leftmargin=*]
\item \arxivid{2604.20920} --- Forget, Then Recall: Learnable Compression and Selective Unfolding via Gist Sparse Attention.
\item \arxivid{2604.20487} --- Knowledge Capsules: Structured Nonparametric Memory Units for LLMs.
\item \arxivid{2604.21829} --- Black-Box Skill Stealing Attack from Proprietary LLM Agents.
\item \arxivid{2604.12710} --- LASA: Language-Agnostic Semantic Alignment at the Semantic Bottleneck for LLM Safety.
\item \arxivid{2604.21751} --- Why are all LLMs Obsessed with Japanese Culture? On the Hidden Cultural and Regional Biases of LLMs.
\item \arxivid{2604.20996} --- AFRILANGTUTOR: Advancing Language Tutoring and Culture Education in Low-Resource Languages with LLMs.
\item \arxivid{2604.21352} --- CARE: Counselor-Aligned Response Engine for Online Mental-Health Support.
\item \arxivid{2604.20904} --- Reinforcing Privacy Reasoning in LLMs via Normative Simulacra from Fiction.
\item \arxivid{2604.21308} --- CI-Work: Benchmarking Contextual Integrity in Enterprise LLM Agents.
\item \arxivid{2510.18731} --- Mitigating Lost-in-Multi-turn Conversation via Curriculum RL with Verifiable Accuracy and Abstention Rewards (RLAAR).
\item \arxivid{2601.10003} --- SocraticKG: Knowledge Graph Construction via QA-Driven Fact Extraction.
\item \arxivid{2604.21380} --- Conjecture and Inquiry: Quantifying Software Performance Requirements via Interactive RAG (IRAP).
\item \arxivid{2601.07262} --- ColorBrowserAgent: Complex Long-Horizon Browser Agent with Adaptive Knowledge Evolution.
\item \arxivid{2604.20855} --- Caesar: Deep Agentic Web Exploration for Creative Answer Synthesis.
\end{enumerate}

\subsubsection{Cross-Module Relationship Index}

The following inter-paper relationships are documented in Wave 2's synthesis deliverable. They are listed here as preview; the full annotation appears in the synthesis output.

\begin{center}
\rowcolors{2}{rowgray}{white}
\begin{tabularx}{\textwidth}{L{2.5cm}L{2.5cm}X}
\toprule
\rowcolor{primary}
\tableheader{Paper A} & \tableheader{Paper B} & \tableheader{Shared Pattern} \\
\midrule
\arxivid{2604.21910} & \arxivid{2604.21744} & Skill-as-markdown as governance layer \\
\arxivid{2604.21744} & \arxivid{2604.20874} & Override-the-prompt invariants (formal) \\
\arxivid{2604.21744} & \arxivid{2601.18491} & Hard Constraints + where/how/what taxonomy \\
\arxivid{2604.20874} & \arxivid{2604.21927} & Bounded-channel theory + regime-fixed CL \\
\arxivid{2604.20874} & \arxivid{2604.20920} & Root theorem + gist compression as instance \\
\arxivid{2604.21003} & \arxivid{2604.21725} & Two-level meta-learning over agent harnesses \\
\arxivid{2604.20300} & \arxivid{2604.20943} & Biological forgetting (hippocampal/sleep) \\
\arxivid{2604.20300} & \arxivid{2511.11439} & Selective forgetting + low-rank-sparse subspace \\
\arxivid{2604.21282} & \arxivid{2604.21255} & Heterogeneous experts + behavior similarity \\
\arxivid{2604.21794} & \arxivid{2603.22823} & Latent vs structured agent communication \\
\arxivid{2604.21816} & \arxivid{2511.00413} & Tool gating + tree-trajectory training \\
\arxivid{2604.20911} & \arxivid{2604.21131} & Omission decay + cross-session threat \\
\arxivid{2604.21477} & \arxivid{2604.20994} & MCP attack-surface joint coverage \\
\arxivid{2604.21598} & \arxivid{2604.21505} & Test gen without ground truth + ambiguity \\
\arxivid{2511.23159} & \arxivid{2604.21187} & Probable-to-provable + SAT+LLM+Lean precedent \\
\arxivid{2604.20920} & \arxivid{2604.20487} & Gist compression + nonparametric capsules \\
\arxivid{2604.12710} & \arxivid{2604.21751} & Semantic-bottleneck safety + cultural bias \\
\arxivid{2604.21829} & \arxivid{2604.21090} & Marketplace threat + structural quality \\
\arxivid{2510.18731} & \arxivid{2601.10003} & Curriculum + Socratic decomposition \\
\arxivid{2604.20855} & \arxivid{2508.21720} & Adversarial-refinement + intermediate-tree \\
\bottomrule
\end{tabularx}
\end{center}

\newpage

% =============================================================
% STAGE 3 -- MISSION PACKAGE
% =============================================================
\stageheader{STAGE 3 --- MISSION PACKAGE}{tertiary}

\section{Milestone Specification}

\subsection{Mission Objective}

Produce a complete adoption-decision package for all 54 priority papers from the CS25--26 sweep, plus engineering-ready integration specifications for the 12+ high-priority drop-in techniques, plus a v1.50-citable BibTeX bibliography. Deliver the result as a self-contained mission package handed off to gsd-skill-creator for execution via Claude Code, with no additional context required to begin Wave 0.

\subsection{Architecture Overview}

\begin{asciibox}
                        Wave 0  +-----------------+
                  Foundation -->| ADR template +  |
                  (Sequential)  | BibTeX scaffold |
                                +-----------------+
                                        |
              +-------------------------+-------------------------+
              v                         v                         v
        +-----------+             +-----------+             +-----------+
Wave 1  | TRACK A   |             | TRACK B   |             | TRACK C   |
(Parallel)  |  M1 + M2  |          |  M3 + M4  |          |  M5 + M6  |
        +-----------+             +-----------+             +-----------+
              |                         |                         |
              +-------------------------+-------------------------+
                                        |
                              Wave 2  +-+----+
                          Synthesis ->| Opus  |
                       (Sequential)   | INTEG |
                                      +------++
                                              |
                              Wave 3  +-------+----+
                          Publication |  Sonnet    |
                       + Verification |  EXEC+VFY  |
                       (Sequential)   +------------+
\end{asciibox}

\subsection{System Layers}

\begin{enumerate}[leftmargin=*]
\item \textbf{Acquisition layer} --- arXiv ID validation, abstract retrieval, BibTeX entry generation. Shared schemas live here.
\item \textbf{ADR layer} --- One Adoption Decision Record per paper. Uniform template. Outputs feed both Wave 2 synthesis and Wave 3 publication.
\item \textbf{Synthesis layer} --- Cross-module relationship matrix, convergent-discovery validation report, GSD architectural impact analysis.
\item \textbf{Publication layer} --- Final mission PDF (this document), BibTeX bibliography file, draft updates to CLAUDE.md and SKILL.md files.
\item \textbf{Verification layer} --- Test harness updates for adopted techniques, verification matrix mapping success criteria to tests.
\end{enumerate}

\subsection{Deliverables}

\begin{center}
\rowcolors{2}{rowgray}{white}
\begin{tabularx}{\textwidth}{C{0.8cm}L{4.2cm}XL{2.5cm}}
\toprule
\rowcolor{primary}
\tableheader{\#} & \tableheader{Deliverable} & \tableheader{Acceptance Criteria} & \tableheader{Spec} \\
\midrule
1 & ADR set (54 records) & One ADR per paper. Each contains: arXiv ID, title, themes, key idea, GSD application, decision (adopt/track/reject), target subsystem, target milestone, validation test ID. & adr-template.md \\
2 & Convergent-Discovery Report & Deep reading of \arxivid{2604.21910}, \arxivid{2604.21744}, \arxivid{2604.20874}, \arxivid{2604.21003} with explicit citation targets in GSD docs. & convergent.md \\
3 & Drop-in Integration Specs & One spec per ADOPT-decision paper (estimated 12+). Each includes engineering ticket text, target file paths, validation test, rollback criterion. & specs/*.md \\
4 & Cross-Module Matrix & ${\geq}15$ inter-paper connections documented (preview table in this PDF lists 20). & xmod.csv + .md \\
5 & GSD Impact Analysis & Per-subsystem table of ``what changes / what does not.'' & impact.md \\
6 & BibTeX Bibliography & \texttt{cs25-26-sweep.bib} with all 54 entries. Compiles cleanly with biblatex. & .bib \\
7 & CLAUDE.md draft additions & Suggested-edits diff for foundational citations. Under 30 added lines. & claude-md.diff \\
8 & SKILL.md draft additions & Linter check additions, ambiguity-checklist additions. & skill-md.diff \\
9 & Safety Test Harness Updates & STD calibration spec, where/how/what BLOCK schema, MCP six-attack-family tests. & safety.md \\
10 & Verification Matrix & Every success criterion mapped to ${\geq}1$ test. & verif.csv \\
11 & Final Mission PDF & This document, recompiled with all wave outputs integrated. Three XeLaTeX passes. & mission.pdf \\
12 & Index landing page & \texttt{index.html} with download cards for PDF + .tex + .bib. & index.html \\
\bottomrule
\end{tabularx}
\end{center}

\subsection{Component Breakdown}

\begin{center}
\rowcolors{2}{rowgray}{white}
\begin{tabularx}{\textwidth}{L{3.5cm}L{3.5cm}C{1.8cm}C{1.8cm}}
\toprule
\rowcolor{primary}
\tableheader{Component} & \tableheader{Depends On} & \tableheader{Model} & \tableheader{Tokens} \\
\midrule
Foundation scaffolding (Wave 0) & --- & Haiku & 30K \\
ADR template + BibTeX skeleton & Foundation & Sonnet & 50K \\
Module 1 deep readings (4 papers) & ADR template & Opus & 80K \\
Module 2 deep readings (10 papers) & ADR template & Sonnet+Opus & 120K \\
Module 3 deep readings (8 papers) & ADR template & Opus & 110K \\
Module 4 deep readings (11 papers) & ADR template & Sonnet & 140K \\
Module 5 deep readings (8 papers) & ADR template & Opus & 110K \\
Module 6 deep readings (13 papers) & ADR template & Opus+Sonnet & 160K \\
Cross-module relationship matrix & All module readings & Opus & 50K \\
Convergent-discovery validation report & M1 + cross-module matrix & Opus & 40K \\
GSD impact analysis & All ADRs & Opus & 50K \\
Drop-in integration specs (12+) & ADRs & Sonnet & 90K \\
Final mission PDF assembly & All synthesis outputs & Sonnet & 60K \\
BibTeX file generation & ADRs & Sonnet & 20K \\
CLAUDE.md / SKILL.md drafts & Synthesis & Opus & 40K \\
Test harness updates & M3 + M4 + M5 readings & Sonnet & 50K \\
Verification matrix & Success criteria + test specs & Sonnet & 25K \\
Index page & Final PDF & Sonnet & 15K \\
\midrule
\rowcolor{lightprimary}
\textbf{Total estimated} & & & \textbf{${\sim}1.24M$} \\
\bottomrule
\end{tabularx}
\end{center}

\subsection{Model Assignment Rationale}

\begin{center}
\rowcolors{2}{rowgray}{white}
\begin{tabularx}{\textwidth}{L{3cm}L{3cm}X}
\toprule
\rowcolor{primary}
\tableheader{Tier} & \tableheader{Share} & \tableheader{Rationale} \\
\midrule
Opus (Agnus) & ${\sim}30\%$ & Convergent-discovery deep reads (M1), bounded-learning policy formalization (M2), Safety Warden taxonomy adoption (M3), verification-and-proof judgment (M5), cultural-sensitivity work (M6), all synthesis. \\
Sonnet (Denise) & ${\sim}60\%$ & Multi-agent module readings (M4, lower judgment overhead), document assembly, BibTeX generation, test harness scaffolding, verification matrix construction. \\
Haiku (Paula) & ${\sim}10\%$ & Initial scaffolding, BibTeX template, ADR template generation, audit logs, simple format conversions. \\
\bottomrule
\end{tabularx}
\end{center}

\section{Wave Execution Plan}

\subsection{Wave Summary}

\begin{center}
\rowcolors{2}{rowgray}{white}
\begin{tabularx}{\textwidth}{C{1cm}L{3cm}XC{2cm}C{2cm}}
\toprule
\rowcolor{primary}
\tableheader{Wave} & \tableheader{Name} & \tableheader{Description} & \tableheader{Mode} & \tableheader{Tokens} \\
\midrule
0 & Foundation & Acquisition scaffolding, ADR template, BibTeX skeleton, citation index. & Sequential & 80K \\
1 & Parallel Research & Three parallel tracks covering the six modules. Each track produces ${\geq}8$ ADRs. & Parallel (3) & 720K \\
2 & Synthesis & Cross-module matrix, convergent-discovery report, GSD impact analysis, drop-in integration specs. & Sequential, Opus-heavy & 230K \\
3 & Publication + Verification & Final PDF assembly, BibTeX file, CLAUDE.md / SKILL.md drafts, test harness updates, verification matrix, index page. & Sequential & 210K \\
\bottomrule
\end{tabularx}
\end{center}

\subsection{Wave 0 --- Foundation (Sequential)}

\begin{center}
\rowcolors{2}{rowgray}{white}
\begin{tabularx}{\textwidth}{L{3cm}XC{1.5cm}}
\toprule
\rowcolor{primary}
\tableheader{Component} & \tableheader{Acceptance} & \tableheader{Model} \\
\midrule
ADR template & Markdown template with all required fields. Renders cleanly. & Haiku \\
BibTeX skeleton & \texttt{cs25-26-sweep.bib} with placeholder entries for 54 papers. & Sonnet \\
Citation index & Lookup table from arXiv ID to module assignment. & Haiku \\
Source bibliography scaffold & Module-grouped enumerated lists ready for Wave 3 publication. & Sonnet \\
Cross-module pre-relationships & Initial 20 inter-paper connections from sweep deliverable as Wave 2 starting point. & Sonnet \\
\bottomrule
\end{tabularx}
\end{center}

\subsection{Wave 1 --- Parallel Research (3 Tracks)}

\subsubsection{Track A: Modules 1 + 2 (14 papers)}

\begin{center}
\rowcolors{2}{rowgray}{white}
\begin{tabularx}{\textwidth}{L{3.5cm}XC{1.5cm}}
\toprule
\rowcolor{primary}
\tableheader{Component} & \tableheader{Acceptance} & \tableheader{Model} \\
\midrule
M1 deep readings (4) & ADR per paper, with explicit GSD-document citation target named. & Opus \\
M2 bounded-learning ADRs (5) & ADR per paper, with linkage to bounded-learning policy doc. & Opus \\
M2 LoRA / PEFT ADRs (5) & ADR per paper, with Silicon Layer benchmark plan note. & Sonnet \\
Track A audit log & Chronicle of decisions made per paper, with rationale. & Haiku \\
\bottomrule
\end{tabularx}
\end{center}

\subsubsection{Track B: Modules 3 + 4 (19 papers)}

\begin{center}
\rowcolors{2}{rowgray}{white}
\begin{tabularx}{\textwidth}{L{3.5cm}XC{1.5cm}}
\toprule
\rowcolor{primary}
\tableheader{Component} & \tableheader{Acceptance} & \tableheader{Model} \\
\midrule
M3 Safety Warden / CAPCOM ADRs (8) & ADR per paper. AgentDoG taxonomy + STD protocol drafted. & Opus \\
M4 Multi-agent / DACP / Chipset ADRs (11) & ADR per paper. Tool Attention engineering ticket drafted. & Sonnet \\
Track B audit log & Chronicle of decisions, with rationale. & Haiku \\
\bottomrule
\end{tabularx}
\end{center}

\subsubsection{Track C: Modules 5 + 6 (21 papers)}

\begin{center}
\rowcolors{2}{rowgray}{white}
\begin{tabularx}{\textwidth}{L{3.5cm}XC{1.5cm}}
\toprule
\rowcolor{primary}
\tableheader{Component} & \tableheader{Acceptance} & \tableheader{Model} \\
\midrule
M5 verification ADRs (8) & ADR per paper. Lean+SAT+LLM workflow precedent extracted. & Opus \\
M6 memory + skills-economy + culture + curriculum ADRs (13) & ADR per paper. DoltHub threat model + cultural-sensitivity additions drafted. & Opus+Sonnet \\
Track C audit log & Chronicle of decisions, with rationale. & Haiku \\
\bottomrule
\end{tabularx}
\end{center}

\subsection{Wave 2 --- Synthesis (Sequential, Opus-heavy)}

\begin{center}
\rowcolors{2}{rowgray}{white}
\begin{tabularx}{\textwidth}{L{4cm}XC{1.5cm}}
\toprule
\rowcolor{primary}
\tableheader{Component} & \tableheader{Acceptance} & \tableheader{Model} \\
\midrule
Cross-module matrix & ${\geq}15$ documented connections; prefer 20+. CSV + narrative. & Opus \\
Convergent-discovery validation report & Each of the four foundation papers receives 500--800 words of analysis with explicit citation targets. & Opus \\
GSD architectural impact analysis & Per-subsystem ``changes / does not change'' table. Names every milestone affected. & Opus \\
Drop-in integration specs & 12+ specs, each ${\sim}300$ words, naming target files, engineering tickets, validation tests, rollback criteria. & Sonnet \\
\bottomrule
\end{tabularx}
\end{center}

\subsection{Wave 3 --- Publication + Verification (Sequential)}

\begin{center}
\rowcolors{2}{rowgray}{white}
\begin{tabularx}{\textwidth}{L{4cm}XC{1.5cm}}
\toprule
\rowcolor{primary}
\tableheader{Component} & \tableheader{Acceptance} & \tableheader{Model} \\
\midrule
Final mission PDF & This document recompiled with all wave outputs integrated. Three XeLaTeX passes. No errors. & Sonnet \\
BibTeX file & \texttt{cs25-26-sweep.bib} compiles with biblatex. All 54 papers cited correctly. & Sonnet \\
CLAUDE.md / SKILL.md draft additions & Diff format. Under 30 lines added to CLAUDE.md; per-skill linter check additions. & Opus \\
Safety test harness updates & STD calibration test, where/how/what BLOCK schema test, MCP six-attack tests. & Sonnet \\
Verification matrix & Every of 12 success criteria mapped to ${\geq}1$ test. & Sonnet \\
Index landing page & HTML with download cards. Color scheme matches PDF. & Sonnet \\
\bottomrule
\end{tabularx}
\end{center}

\subsection{Cache Optimization Strategy}

\begin{itemize}[leftmargin=*]
\item Wave 0 outputs (ADR template, BibTeX skeleton, citation index) are loaded once at the start of every Wave 1 track.
\item Each track loads its own subset of paper abstracts + theme assignments. No cross-track loading until Wave 2.
\item Wave 2 loads only the ADR set as compressed summaries (each ADR ${\sim}500$ tokens), not full paper readings.
\item Wave 3 loads only Wave 2 outputs plus the final-PDF skeleton; raw ADRs not re-loaded.
\item Anchor reuse: the four Module 1 deep readings appear in every wave's context as foundational citations. Cache once.
\end{itemize}

\subsection{Token Budget}

\begin{center}
\rowcolors{2}{rowgray}{white}
\begin{tabularx}{\textwidth}{L{3cm}XC{2cm}C{2cm}}
\toprule
\rowcolor{primary}
\tableheader{Wave} & \tableheader{Components} & \tableheader{Tokens} & \tableheader{Cumulative} \\
\midrule
0 & Foundation (5 components) & 80K & 80K \\
1 & Three parallel tracks & 720K & 800K \\
2 & Synthesis (4 components) & 230K & 1{,}030K \\
3 & Publication + Verification (6 components) & 210K & 1{,}240K \\
\midrule
\rowcolor{lightprimary}
\textbf{Total} & & \textbf{1.24M} & \\
\bottomrule
\end{tabularx}
\end{center}

Estimated context windows: 5--6 sessions at 200K each. Squadron-equivalent throughput; Fleet profile chosen for parallelism in Wave 1.

\subsection{Dependency Graph}

\begin{asciibox}
W0.foundation
   |
   +--> W1.A.M1_readings -->\
   +--> W1.A.M2_readings ----+
   |                           \
   +--> W1.B.M3_readings -----> W2.matrix
   +--> W1.B.M4_readings -----> W2.convergent
   |                           > W2.impact
   +--> W1.C.M5_readings -----> W2.specs
   +--> W1.C.M6_readings ----+
                              /
                              v
                     +--> W3.pdf
                     +--> W3.bib
                     +--> W3.diffs (CLAUDE.md, SKILL.md)
                     +--> W3.tests
                     +--> W3.matrix
                     +--> W3.index
\end{asciibox}

\section{Test Plan}

\subsection{Test Categories}

\begin{center}
\rowcolors{2}{rowgray}{white}
\begin{tabularx}{\textwidth}{L{4cm}XC{1.5cm}C{2cm}}
\toprule
\rowcolor{primary}
\tableheader{Category} & \tableheader{Purpose} & \tableheader{Count} & \tableheader{Action} \\
\midrule
Safety-critical & Source quality, attribution, copyright, sovereignty gates & 7 & BLOCK \\
Core functionality & Per-success-criterion correctness & 20 & BLOCK \\
Integration & Cross-module + cross-deliverable consistency & 9 & BLOCK \\
Edge cases & Missing data, ambiguity, future-paper handling & 6 & LOG \\
\midrule
\rowcolor{lightprimary}
\textbf{Total} & & \textbf{42} & \\
\bottomrule
\end{tabularx}
\end{center}

\subsection{Safety-Critical Tests}

\begin{center}
\rowcolors{2}{rowgray}{white}
\begin{tabularx}{\textwidth}{C{1.2cm}L{4cm}X}
\toprule
\rowcolor{primary}
\tableheader{ID} & \tableheader{Verifies} & \tableheader{Expected Behavior} \\
\midrule
SC-SRC & Source quality & All 54 sources are arXiv preprints. No popular-press paraphrases. arXiv IDs resolve. \\
SC-NUM & Numerical attribution & Every numerical claim in the document (decay rates, F1 scores, token counts, percentages) attributed to a specific paper. \\
SC-IND & Indigenous attribution & Every reference to PNW Indigenous nations names them specifically (Coast Salish, Lummi/Lhaq'temish, Tulalip, Stillaguamish, Wanapum). Never generic ``Indigenous peoples.'' \\
SC-PARA & No verbatim reproduction & No paraphrase reproduces ${\geq}15$ words consecutively from any source paper's abstract. \\
SC-MED & Mental-health framing & CARE-template adoption discussion uses counselor-alignment framing exclusively. No diagnostic / therapeutic claims. \\
SC-SOV & Sovereignty gating & AFRILANGTUTOR adaptation discussion is explicitly marked as requiring nation-specific consultation BEFORE implementation. \\
SC-ADV & No policy advocacy & Document presents adoption recommendations without advocating for specific legislative or policy positions. \\
\bottomrule
\end{tabularx}
\end{center}

\subsection{Core Functionality Tests}

\begin{center}
\rowcolors{2}{rowgray}{white}
\begin{tabularx}{\textwidth}{C{1.2cm}L{4cm}X}
\toprule
\rowcolor{primary}
\tableheader{ID} & \tableheader{Verifies} & \tableheader{Expected Behavior} \\
\midrule
CF-01 & ADR completeness & 54 ADRs exist; each has all required fields; each has a decision and a target. \\
CF-02 & Convergent-discovery citations & Each of \arxivid{2604.21910}, \arxivid{2604.21744}, \arxivid{2604.20874}, \arxivid{2604.21003} has a named GSD-document citation target. \\
CF-03 & Drop-in spec count & 12 or more ADOPT-decision papers have one-page integration specs. \\
CF-04 & AgentDoG schema & Where/how/what taxonomy appears as a structured schema (table or YAML) ready for code adoption. \\
CF-05 & Tool Attention ticket & Lazy-schema-loading design note exists with token-budget projections and a Q2-2026 target. \\
CF-06 & Black-Box Skill Stealing threat model & Attack taxonomy encoded in DoltHub threat-model document. \\
CF-07 & Five-Principle linter & Skill-creator linter check spec drafted, including pass/fail criteria per principle. \\
CF-08 & Worker/Evaluator/Evolution naming & Naming convention documented in gsd-mission-management-roles.md draft. \\
CF-09 & STD measurement protocol & Per-model Safe Turn Depth measurement protocol exists as code-ready test pattern. \\
CF-10 & BibTeX compiles & \texttt{cs25-26-sweep.bib} compiles with biblatex without errors; all 54 entries present. \\
CF-11 & Cross-module matrix size & Matrix documents ${\geq}15$ inter-paper connections. \\
CF-12 & Document compiles & Final PDF compiles in three XeLaTeX passes without errors or warnings. \\
CF-13 & M1 paper coverage & All 4 papers in Module 1 receive ${\geq}500$ words of analysis. \\
CF-14 & M2 paper coverage & All 10 papers in Module 2 receive an ADR. Bounded-learning policy doc updated. \\
CF-15 & M3 paper coverage & All 8 papers in Module 3 receive an ADR. Safety Warden update spec drafted. \\
CF-16 & M4 paper coverage & All 11 papers in Module 4 receive an ADR. Chipset validation citations integrated. \\
CF-17 & M5 paper coverage & All 8 papers in Module 5 receive an ADR. Proof-companion citations identified. \\
CF-18 & M6 paper coverage & All 13 papers in Module 6 receive an ADR. \\
CF-19 & Decision distribution sane & ADOPT decisions number 12--18; TRACK decisions number 25--35; rejections explicitly justified. \\
CF-20 & Index page & \texttt{index.html} provides download links for PDF + .tex + .bib; color scheme matches PDF. \\
\bottomrule
\end{tabularx}
\end{center}

\subsection{Integration Tests}

\begin{center}
\rowcolors{2}{rowgray}{white}
\begin{tabularx}{\textwidth}{C{1.2cm}L{4cm}X}
\toprule
\rowcolor{primary}
\tableheader{ID} & \tableheader{Verifies} & \tableheader{Expected Behavior} \\
\midrule
IN-01 & ADR $\to$ spec consistency & Every ADR with decision ADOPT has exactly one drop-in integration spec; field values are consistent. \\
IN-02 & ADR $\to$ test consistency & Every ADR with decision ADOPT has at least one validation test referenced; the test exists in the test harness updates. \\
IN-03 & ADR $\to$ matrix consistency & Every cross-module relationship in the matrix references two ADRs that both exist. \\
IN-04 & Citation density vs adoption & Each ADOPT-decision paper has ${\geq}1$ citation in synthesis outputs (convergent report, impact analysis, or specs). \\
IN-05 & Bounded-learning + Silicon Layer cross-link & Module 2's Retrofit-as-bounded-learning-formalization is cross-referenced from the Silicon Layer LoRA spec. \\
IN-06 & MCP security joint coverage & Tool Attention (Module 4), MCP Pitfall Lab (Module 3), and Function Hijacking (Module 3) appear together in the planning-bridge MCP design notes. \\
IN-07 & Self-containment & A reader with no prior GSD context can understand the mission objective and decisions from the PDF alone. \\
IN-08 & Convergent-discovery + bounded-learning link & Root Theorem (Module 1) explicitly cited as foundation for bounded-learning policy in Module 2 outputs. \\
IN-09 & Test plan coverage & Every success criterion mapped to ${\geq}1$ test ID in the verification matrix. \\
\bottomrule
\end{tabularx}
\end{center}

\subsection{Edge-Case Tests}

\begin{center}
\rowcolors{2}{rowgray}{white}
\begin{tabularx}{\textwidth}{C{1.2cm}L{4cm}X}
\toprule
\rowcolor{primary}
\tableheader{ID} & \tableheader{Verifies} & \tableheader{Expected Behavior} \\
\midrule
EC-01 & Withdrawn-paper handling & If any arXiv ID is withdrawn before publication, the ADR is updated to a tombstone with explanation. \\
EC-02 & Ambiguous theme handling & Papers spanning multiple themes are explicitly noted in the ADR; module assignment justified. \\
EC-03 & Out-of-scope adjacency & Adjacent topics not covered (e.g., reinforcement learning theory) explicitly listed in scope-boundaries section. \\
EC-04 & Future-paper extension & Mission output describes how to extend the ADR set if additional papers from the corpus are later prioritized. \\
EC-05 & Source-quality outliers & Any paper with non-standard publication venue flagged in its ADR. \\
EC-06 & Counterfactual decisions & At least one rejected paper has its rejection rationale documented (not just left blank). \\
\bottomrule
\end{tabularx}
\end{center}

\subsection{Verification Matrix}

\begin{center}
\rowcolors{2}{rowgray}{white}
\begin{tabularx}{\textwidth}{XL{4cm}C{2cm}}
\toprule
\rowcolor{primary}
\tableheader{Success Criterion} & \tableheader{Test IDs} & \tableheader{Status} \\
\midrule
1. All 54 papers have an ADR. & CF-01, CF-13, CF-14, CF-15, CF-16, CF-17, CF-18 & Pending \\
2. Convergent-discovery papers have citation targets. & CF-02, IN-08 & Pending \\
3. Drop-in technique papers have integration specs. & CF-03, IN-01 & Pending \\
4. STD measurement protocol implemented. & CF-09 & Pending \\
5. Tool Attention design note exists. & CF-05, IN-06 & Pending \\
6. AgentDoG taxonomy encoded. & CF-04, IN-06 & Pending \\
7. Black-Box Skill Stealing taxonomy in DoltHub doc. & CF-06 & Pending \\
8. Five-Principle linter check drafted. & CF-07 & Pending \\
9. Worker/Evaluator/Evolution naming adopted. & CF-08 & Pending \\
10. Cross-module matrix has ${\geq}15$ connections. & CF-11, IN-03 & Pending \\
11. \texttt{cs25-26-sweep.bib} compiles cleanly. & CF-10 & Pending \\
12. Final PDF compiles in 3 XeLaTeX passes. & CF-12, CF-20 & Pending \\
\bottomrule
\end{tabularx}
\end{center}

Every success criterion maps to at least one test. SC-* tests apply universally regardless of criterion.

\section{Mission Crew Manifest}

Fleet activation profile (6 modules) selected. Crew names follow PNW bioregion naming convention.

\begin{center}
\rowcolors{2}{rowgray}{white}
\begin{tabularx}{\textwidth}{L{2cm}L{2.5cm}X}
\toprule
\rowcolor{primary}
\tableheader{Role} & \tableheader{Agent} & \tableheader{Responsibility} \\
\midrule
FLIGHT & ORCA (Opus) & Mission direction; go/no-go gates; cross-track decisions. \\
PLAN & RAVEN (Sonnet) & Wave decomposition; parallel-track optimization; dependency management. \\
SCOUT-1 & SALMON (Sonnet) & Paper acquisition for Tracks A + B; arXiv ID validation. \\
SCOUT-2 & STEELHEAD (Sonnet) & Paper acquisition for Track C; cross-module relationship preview. \\
EXEC-1 & CEDAR (Sonnet) & Track A document production (Modules 1 + 2 ADRs). \\
EXEC-2 & DOUGFIR (Sonnet) & Track B document production (Modules 3 + 4 ADRs). \\
EXEC-3 & HEMLOCK (Sonnet) & Track C document production (Modules 5 + 6 ADRs). \\
CRAFT-FOUND & FOXGLOVE (Opus) & Module 1 specialist: convergent-discovery deep readings + citation-target identification. \\
CRAFT-BOUND & TRILLIUM (Opus) & Module 2 specialist: bounded-learning policy formalization. \\
CRAFT-SAFE & HAWK (Opus) & Module 3 specialist: Safety Warden taxonomy + STD protocol + MCP security. \\
CRAFT-AGENT & HERON (Sonnet) & Module 4 specialist: chipset validation citations + DACP comparison + Tool Attention engineering ticket. \\
CRAFT-VERIF & SALAL (Opus) & Module 5 specialist: proof-companion linkage + Lean+SAT+LLM precedent. \\
CRAFT-CULT & FERN (Opus) & Module 6 specialist: cultural-sensitivity additions + DoltHub threat model + curriculum design. \\
VERIFY & EAGLE (Sonnet) & Source validation; arXiv ID resolution; abstract-claim verification. \\
INTEG & CASCADE (Opus) & Cross-module relationship matrix; convergent-discovery validation report; impact analysis. \\
CAPCOM & TIBSFOX (Human) & HITL gates between waves; final adoption decisions on contested ADRs. \\
BUDGET & BEAVER (Sonnet) & Token tracking; per-wave budget enforcement; session planning. \\
SURGEON & KINGFISHER (Sonnet) & Quality monitoring; flagging research degradation across long sessions. \\
LOG & WREN (Haiku) & Audit trail; per-track chronicle; commit-message generation. \\
\bottomrule
\end{tabularx}
\end{center}

\section{Execution Summary}

\begin{center}
\rowcolors{2}{rowgray}{white}
\begin{tabularx}{\textwidth}{L{4cm}X}
\toprule
\rowcolor{primary}
\tableheader{Metric} & \tableheader{Value} \\
\midrule
Total papers covered & 54 (out of 930-paper sweep corpus) \\
Modules & 6 \\
Activation profile & Fleet (19 roles) \\
Estimated token budget & 1.24M tokens \\
Estimated context windows & 5--6 \\
Estimated session count & 3--4 \\
Wave count & 4 (W0 sequential, W1 three-parallel, W2 sequential, W3 sequential) \\
Deliverables & 12 \\
Tests & 42 (7 SC + 20 CF + 9 IN + 6 EC) \\
Success criteria & 12 \\
Cross-module connections (preview) & 20 \\
Drop-in integration specs (estimated) & 12--18 \\
Target milestone & v1.50 (Half A retrospective foundation) \\
Target hand-off & gsd-skill-creator via Claude Code \\
\bottomrule
\end{tabularx}
\end{center}

\vspace{1em}
\begin{quotebox}
\centering\sffamily
The fox finds the trail by recognizing the seams the deer made walking it.\\
The seams were always there. The fox just looks where the path is.\\[0.5em]
\textit{--- through-line, mission close}
\end{quotebox}


% ============================================================
% APPENDICES — Phase 804 Wave 3.1 integration of Wave 1/2 outputs
% ============================================================

\clearpage
\appendix

\section{Convergent-Discovery Validation Report (Phase 801)}

\textit{This appendix embeds the Phase 801 deliverable from \texttt{.planning/missions/cs25-26-sweep/work/synthesis/convergent-discovery.md} verbatim, retypeset for the published PDF. It satisfies CS25-03 (anchor-paper coverage \textgreater{}500 words per anchor) for the four Module-1 anchors.}

\subsection{§1 Why convergent discovery
matters}

The strongest external validation a load-bearing architectural decision
can receive is \textbf{convergent discovery}: an independent
intellectual current --- different authors, different problem domain,
different vocabulary, no line of communication --- arriving at the same
load-bearing structure. Citation is influence; convergent discovery is
\emph{evidence}. When two systems built without contact derive
equivalent shapes for the same job, the shape is more likely a property
of the problem than an idiosyncrasy of either author. The patterns are
the seams the deer made walking it; the deer were always walking the
same watershed.

This report is the v1.49.575 analog of two earlier convergent-discovery
moments in the GSD release line. v1.49.573 enumerated \textbf{seven}
convergent-discovery anchors as the keystone of the milestone
retrospective --- each anchor a separate paper independently arriving at
one of the architectural primitives GSD had already shipped. v1.49.574
named a \textbf{two-party rhyme}: SIGReg (LeJEPA's
stop-gradient-replacing isotropy regulariser) and counter-based
megakernel scheduling (CUDA-graph era kernel-launch consolidation)
coming from completely unrelated sub-fields and converging on the same
design move --- replace heuristics with a single principled regulariser,
replace ten launch points with one. The v1.49.574 retrospective took
care to call this rhyme ``a productive analogy at the design-discipline
level, not a theorem''; we keep that calibration here.

CS25-03 distinguishes convergent discovery from cross-validation.
\textbf{Cross-validation} is the lighter pattern: paper A cites paper
B's primitive and re-tests it, or one party explicitly builds on the
other. Useful, but it does not rule out the idiosyncrasy hypothesis ---
paper A's design might still be wrong, just shared. \textbf{Convergent
discovery} is heavier: independently-derived equivalent structures
arising without communication. None of the four anchor authors of this
milestone --- the population-genetics workflow team behind
\arxivid{2604.21910}, the proteomics governance authors of
\arxivid{2604.21744} (GROUNDINGmd), the formal-context-theory author of
\arxivid{2604.20874} (Root Theorem), the agentic-harness authors of
\arxivid{2604.21003} (Last Harness) --- have any line of communication to
the GSD ecosystem. None cite GSD; GSD does not cite them; the works
arrived in 2026 inside the CS25--26 arXiv listing as part of a 930-paper
sweep that was filtered down to 54 priority papers and then to four
anchor papers in Module 1. The structural matches are not borrowed; they
are derived.

What makes Module 1 the \emph{keystone} of v1.49.575 is that all four
anchors land on \textbf{architectural foundations} rather than secondary
tactics: the system-level shape of an LLM-driven research pipeline
(\texttt{21910}), the document-level safety doctrine that overrides user
prompts (\texttt{21744}), the information-theoretic substrate of
long-running agent context (\texttt{20874}), and the meta-evolutionary
loop a self-improving harness must run (\texttt{21003}). Each is
load-bearing --- remove it and the rest of the structure tilts. When
four works derived in four different problem domains land on four of
GSD's load-bearing primitives, the primitives are not local choices made
by one author. They are the engineering equivalent of natural laws ---
the seams. The remainder of this report does the per-anchor work: each
anchor receives a structured analysis covering what the paper does, what
GSD already does, where the convergence lies, and which GSD document is
the explicit citation target for Phase 805's draft additions.

\subsection{\texorpdfstring{§2 Anchor 1 --- \arxivid{2604.21910}
Skills-as-markdown
three-layer}{§2 Anchor 1 --- \arxivid{2604.21910} Skills-as-markdown three-layer}}

\textbf{The paper.} \emph{From Research Question to Scientific Workflow:
Leveraging Agentic AI for Science Automation} (arXiv:\arxivid{2604.21910}).
Themes: plan-decomposition, knowledge-as-skill,
deterministic-layer-confinement, verify-proof, skill-as-md.

\textbf{What the paper does.} It proposes a three-layer agentic
architecture for population-genetics workflow synthesis. The
\textbf{semantic layer} uses an LLM to extract structured intent from a
natural-language research question --- \emph{``compare diversity across
these populations using these markers''} becomes a typed intent record.
The \textbf{deterministic layer} dispatches that intent to validated DAG
generators that emit reproducible workflow scripts; this layer contains
no LLM calls and is therefore behaviourally stable across runs. The
\textbf{knowledge layer} holds domain-expert markdown documents called
``Skills'' that encode three things: vocabulary mappings (this term in
the literature corresponds to that parameter on the tool), parameter
constraints (this method requires this kind of input distribution), and
optimisation strategies (when N is small prefer method X). LLM
non-determinism is confined entirely to intent extraction; identical
intents are guaranteed to produce identical workflows because the
deterministic layer is pure dispatch. The paper reports two numbers
worth pinning: full intent-match accuracy rises from a 44\% baseline (no
skills) to 83\% (skills enabled); skill-driven \emph{deferred} workflow
generation --- generate the script only when needed, not eagerly ---
cuts data transfer 92\% versus eager generation. The pattern was
developed independently of any prior GSD documentation; the authors are
population geneticists, not LLM-orchestration people.

\textbf{What GSD already does.} GSD's intake pipeline runs the same
three-layer split under different names. \textbf{vision-to-mission} is
the semantic layer: a user describes a creative or research vision in
natural language, and an LLM extracts a structured mission package ---
ambition record, target artefacts, milestone shape, fleet profile.
\textbf{research-mission-generator} is the deterministic layer: given a
structured intent, it produces a three-stage output (Vision Guide →
Research Reference → Mission Spec) with stable scaffolding,
deterministic file layout, validated frontmatter.
\textbf{gsd-skill-creator} is the knowledge / authoring layer: domain
experts (the human author and the system itself) write SKILL.md
documents that encode vocabulary mappings (when the user says ``ship'',
trigger this command path), parameter constraints (this skill requires
this kind of phase artefact), and optimisation strategies (this skill
prefers Sonnet over Opus because of the workload class). The pipeline
was built in pieces across many milestones --- vision-to-mission landed
first as a \texttt{Tibsfox/skills} cartridge, research-mission-generator
second as the LaTeX research path, gsd-skill-creator third as the
adaptive-learning layer this entire repo is named after --- but the
three pieces line up.

\textbf{The convergence.} The structural equivalence is point-for-point.
Their \emph{semantic layer} (LLM intent extraction, only place
non-determinism lives) maps to GSD's vision-to-mission. Their
\emph{deterministic layer} (validated generators, identical-intent →
identical-output guarantee) maps to research-mission-generator. Their
\emph{knowledge layer} (domain-expert-authored markdown skills with
vocabulary / constraints / strategies) maps to gsd-skill-creator's
\texttt{.claude/skills/} and \texttt{examples/} skill libraries. Both
confine non-determinism to intent extraction. Both treat the knowledge
layer as \textbf{markdown} authored by domain experts rather than
learned weights or runtime-rewritable state --- a non-obvious choice
that the paper makes for population genetics and GSD makes for
cross-domain creative tooling. Both use the deterministic layer as the
\textbf{gate-able surface}: in the paper, that's where reproducibility
is enforced; in GSD, that's where CAPCOM gates and the Safety Warden's
BLOCK authority operate. The paper's 44\%→83\% intent-match number is
the kind of external benchmark GSD has not had until now --- a published
quantification of how much the knowledge layer raises end-to-end
fidelity. The 92\% data-transfer reduction for deferred generation has
its analog in GSD's pull-based work assignment (GUPP / Get Up and Push
Protocol) and the lazy-loading of skill bodies until activation: same
shape, different vocabulary.

\textbf{Explicit citation target.} This paper is the canonical citation
for the GSD vision-to-mission / research-mission-generator /
gsd-skill-creator pipeline. It lands in three places: (1)
\textbf{\texttt{docs/skill-architecture/three-tier-pipeline.md}} (NEW,
propose for Phase 805 SKILL.md additions) --- formal write-up that lifts
the three-layer naming ``semantic / deterministic / knowledge'' as the
canonical architecture description and cites this paper as the
convergent-discovery anchor; (2) \textbf{CLAUDE.md} ``Architectural
Foundations'' preamble --- one-line foundational citation under
CS25-09's ≤30-line CLAUDE.md edit budget; (3)
\textbf{\texttt{docs/cartridge/SKILL.md} authoring guide} --- cite as
external validation that markdown-skill format itself is a defensible
storage choice, not a transient one. The 44\%→83\% and 92\% numbers feed
v1.50 retrospective Half A as the external benchmark for the pipeline's
fidelity claim.

\subsection{\texorpdfstring{§3 Anchor 2 --- \arxivid{2604.21744}
GROUNDINGmd}{§3 Anchor 2 --- \arxivid{2604.21744} GROUNDINGmd}}

\textbf{The paper.} \emph{Agentic AI-assisted coding offers a unique
opportunity to instill epistemic grounding during software development}
--- colloquially \textbf{GROUNDINGmd} (arXiv:\arxivid{2604.21744}). Themes:
safety-warden, hard-constraints, override-the-prompt,
governance-as-document, bounded-learn.

\textbf{What the paper does.} It proposes a community-governed,
field-scoped epistemic grounding document that the agent loads before
any task. The document defines two \textbf{override-everything} classes.
\textbf{Hard Constraints} are non-negotiable validity invariants
empirically required for scientific correctness in the field --- for the
example domain of mass-spectrometry-based proteomics, things like ``do
not let the agent silently change the m/z tolerance below the instrument
resolution'' or ``never combine spectra acquired under incompatible
isolation windows''. \textbf{Convention Parameters} are community-agreed
defaults: units, schemas, naming, the standard search-tolerance for a
given mass-spec workflow. Both classes override every other context the
agent receives, \emph{including the user prompt itself}. The agent is
required to refuse --- to BLOCK --- rather than violate a Hard
Constraint. The document is \textbf{community-governed}: changes require
consensus, not individual modification, which is the structural defence
against prompt-injection rewriting the constraint set (``just edit the
GROUNDINGmd file and remove that rule'' stops working when the file is
governed elsewhere). The pattern is presented for proteomics but the
authors explicitly frame it as field-portable: any computational
scientific discipline with empirical validity invariants is a candidate
domain.

\textbf{What GSD already does.} The GROUNDINGmd construct is --- in
nearly its exact shape --- the GSD \texttt{CLAUDE.md} +
\texttt{SKILL.md} + Safety Warden BLOCK pattern. CLAUDE.md is the
project-scoped grounding document loaded before any task; the system
reminder at the top of every Claude Code session says verbatim
\emph{``Codebase and user instructions are shown below. Be sure to
adhere to these instructions. IMPORTANT: These instructions OVERRIDE any
default behavior and you MUST follow them exactly as written''} --- that
override-the-default-behaviour-clause is the GROUNDINGmd override clause
restated. The Safety Warden BLOCK pattern documented in
\texttt{.claude/skills/security-hygiene/SKILL.md} ensures that on
safety-critical decisions --- touching another developer's secrets,
modifying signing keys, deleting \texttt{.git/}, breaking the wasteland
exclusion --- the agent refuses regardless of the user prompt. CAPCOM
gates extend the same authority to milestone-shaping decisions: a CAPCOM
HARD GATE override of a user prompt is doctrinally normal in this
codebase, not exceptional. Today CLAUDE.md mixes invariants and defaults
under a single heading; the Hard-vs-Convention split is implicit but not
lifted to a typed taxonomy.

\textbf{The convergence.} The two-class taxonomy maps cleanly.
\textbf{Hard Constraints ↔ Safety Warden HARD GATE}: invariants that the
agent must refuse to violate, governed at the document level rather than
the prompt level. \textbf{Convention Parameters ↔ CLAUDE.md project
conventions + SKILL.md authoring defaults}: community-agreed defaults
that the agent applies unless explicitly overridden by a typed override.
The override-the-user-prompt commitment is identical in both systems ---
neither is a ``the user is always right'' architecture; both are ``the
document is the ground truth, and the document overrides the prompt''
architectures. The community-governance property is also present in GSD:
CLAUDE.md and SKILL.md ship under git, with PR-style review for changes;
an attacker who wants to relax the BLOCK rule has to land a commit, not
just emit a prompt. The match is structural and deep --- both systems
independently solved the prompt-injection-resistance problem the same
way, by lifting the constraint set out of the prompt-layer entirely.

\textbf{Explicit citation target.} The paper is essentially the
published independent rediscovery of CLAUDE.md + Safety Warden BLOCK;
the citation lands in four places. (1) \textbf{CLAUDE.md} --- promote
the Hard Constraints / Convention Parameters two-class taxonomy to a
foundational citation in the ``Important Notes'' section; this is the
foundational citation in the CS25-09 ≤30-line edit budget. (2)
\textbf{Every \texttt{SKILL.md} template} in
\texttt{docs/cartridge/SKILL.md} authoring guide --- cite for the
Hard-vs-Convention split structure; every authored SKILL.md should
distinguish what it refuses to do (Hard Constraints) from what it
defaults to but allows override on (Convention Parameters). This becomes
a structural-completeness criterion. (3)
\textbf{\texttt{.claude/skills/security-hygiene/SKILL.md}} --- the BLOCK
doctrine cites this paper as the published independent rediscovery of
the override-user-prompt rule; an external citation is durable defence
against future ``but the user asked for it'' pressure. (4)
\textbf{\texttt{src/cartridge/linter/structural-completeness.ts}} ---
Module 3's CS25-17 structural-completeness linter (HB-05) adopts
GROUNDINGmd's two-class taxonomy as its schema, and includes a ``Hard
Constraints declared'' check on SKILL.md fixtures. The
override-the-user-prompt formulation must be cited (paraphrased ≤14
consecutive words to satisfy the SC-PARA gate) so the doctrinal lineage
is unambiguous.

\subsection{\texorpdfstring{§4 Anchor 3 --- \arxivid{2604.20874} Root
Theorem of Context
Engineering}{§4 Anchor 3 --- \arxivid{2604.20874} Root Theorem of Context Engineering}}

\textbf{The paper.} \emph{The Root Theorem of Context Engineering}
(arXiv:\arxivid{2604.20874}). Themes: context-management, bounded-tape,
formal-axioms, homeostatic-persistence.

\textbf{What the paper does.} It formalises the agent's working context
as a \textbf{bounded lossy channel}. Two axioms: (i) the window is
finite; (ii) per-token information quality decays monotonically with
cumulative volume. From these two axioms five consequences are derived.
\textbf{(C1)} Quality degrades monotonically with token volume
\emph{independent of nominal window size} --- a 1M-token model is not
exempt from the decay; it merely slides the absolute number at which the
decay becomes operationally painful. \textbf{(C2)} Signal density and
token-count are \emph{independent} optimisation variables; optimising
one does not optimise the other, and the gradients can point opposite
directions. \textbf{(C3)} Gating must trigger on \textbf{fidelity
thresholds, not on capacity}: ``is the context window 80\% full?'' is
the wrong question; ``has the signal-to-noise dropped below the
tolerated threshold?'' is the right one. \textbf{(C4)} Only
architectures that \emph{accumulate--compress--rewrite--shed} can
sustain understanding indefinitely; pure accumulation diverges.
\textbf{(C5)} The compressor needs an \textbf{external verification
gate} because it operates inside the same channel it is compressing ---
self-review is structurally inadequate. The paper demonstrates the
framework empirically with a 60+-session persistent agent showing stable
memory footprint and stable task quality, contrasted with baselines that
degrade after roughly twenty sessions.

\textbf{What GSD already does.} The bounded-tape framing has been
Tibsfox's informal articulation across many docs: \emph{The Space
Between} talks about working context as a finite bandwidth channel;
CLAUDE.md design rationale references the bounded-tape framing for
CAPCOM gate policy; the milestone retrospective + feed-forward lifecycle
is described in the through-lines of v1.49.570/.571/.572/.574 as
``accumulate during execution, compress in retrospective, rewrite into
substrate docs, shed via the done-retirement skill''; the
bounded-learning policy invariants --- the 20\% cap on auto-load slot
churn, the 3-correction-min, the 7-day cooldown --- operate on a notion
of bounded signal that the policy refuses to relax. None of this had a
published mathematical foundation until now.

\textbf{The convergence --- five consequences mapped.} Each derived
consequence maps directly to a shipped GSD architectural decision.

\begin{itemize}
\tightlist
\item
  \textbf{C1 (monotone quality decay independent of window size) →
  bounded-learning caps don't relax with bigger context windows.} GSD
  treats Opus, Sonnet, and Haiku the same with respect to
  bounded-learning policy; the 20\% rule does not loosen for a model
  with a larger advertised context. C1 is the formal justification for
  that refusal: the decay is monotone regardless of how much head-room
  remains. The orchestration layer in \texttt{src/orchestration/}
  honours the same policy across model tiers.
\item
  \textbf{C2 (signal/token-count independence) → why CAPCOM gates
  measure fidelity, not ``is the context full?''} The reinforcement
  series (\texttt{src/reinforcement/}, MA-1..MD-6) operates over
  signal-bearing events, not raw token streams, exactly because of C2:
  optimising for token efficiency is not the same as optimising for
  signal preservation, and the milestone retrospective's compress step
  targets the latter.
\item
  \textbf{C3 (fidelity-threshold gating, not capacity gating) → CAPCOM
  gating policy.} Today's CAPCOM gates trigger on phase boundaries, on
  human-authorisation prompts, on safety-relevant decision points ---
  \emph{not} on token-budget exhaustion. C3 is the formal derivation of
  that policy. The 70\% context-fracture threshold from the Tool
  Attention paper (\arxivid{2604.21816}, HB-01 source) is one specific
  instance of a fidelity threshold; the Root Theorem says the
  \emph{category} is right, and Tool Attention says one threshold inside
  that category is 70\% on a particular ISO-score metric.
\item
  \textbf{C4 (accumulate--compress--rewrite--shed homeostatic lifecycle)
  → milestone retrospective + feed-forward lifecycle.} Each phase of the
  milestone is one stage of a homeostatic cycle: Wave 0 / phase
  execution = \emph{accumulate}; Wave 2 synthesis = \emph{compress}
  (cross-module matrix, convergent-discovery report, GSD-impact analysis
  are all compression of the Wave-1 ADR mass); Wave 3 publication + the
  substrate doc additions = \emph{rewrite} (the compressed signal gets
  rewritten into durable substrate docs in \texttt{docs/substrate/});
  the \texttt{done-retirement} skill and milestone closure =
  \emph{shed}. The cycle name in the FSFM bounded-learning forgetting
  taxonomy from ADR-CS25-\arxivid{2604.20300} is literally
  \texttt{accumulate-compress-rewrite-shed}. The match is verbatim, not
  analogical.
\item
  \textbf{C5 (external verification gate on the compressor) →
  fresh-context-subagent verification pattern.} This milestone's
  three-track Wave-1 fleets --- Track A, Track B, Track C --- are
  \emph{fresh-context subagents}. The orchestrator (the agent doing the
  compression) does not self-review its own outputs; a separate-context
  subagent provides the verification gate. CAPCOM is the structural
  external gate. C5 is the formal derivation of why this is necessary:
  the compressor is inside the channel it compresses, so self-review is
  inadequate by construction. The umwelt module (\texttt{src/umwelt/},
  M7) makes the same point in Friston / Markov-blanket vocabulary: the
  boundary is external because internal observation cannot certify
  itself.
\end{itemize}

\textbf{Explicit citation target.} Three landing places, all flowing
from the consequence mapping above. (1)
\textbf{\texttt{docs/substrate/bounded-tape-root-theorem.md}} (NEW,
propose for Phase 805) --- primary landing: a substrate document that
lifts the two-axiom / five-consequence structure verbatim and reproduces
the consequence-to-decision map above. The orchestrator brief names this
as the explicit citation target. (2) \textbf{\emph{The Space Between}
information-theory chapter} --- the bounded-tape framing now has a
published derivation to cite as formal foundation. (3) \textbf{CLAUDE.md
design-rationale section} --- one-line citation under CS25-09's ≤30-line
edit budget, naming the paper as the formal foundation for the
bounded-tape framing. The 60+-session empirical demonstration is the
citation when the v1.50 retrospective claims GSD's substrate is
empirically validated rather than purely-engineered.

\subsection{\texorpdfstring{§5 Anchor 4 --- \arxivid{2604.21003} Last
Harness}{§5 Anchor 4 --- \arxivid{2604.21003} Last Harness}}

\textbf{The paper.} \emph{The Last Harness You'll Ever Build}
(arXiv:\arxivid{2604.21003}). Themes: skill-learning, meta-evolution,
agent-orchestration, harness-as-artefact.

\textbf{What the paper does.} It proposes a two-level framework for
self-improving agentic harnesses. At the \textbf{inner level}, the
\textbf{Harness Evolution Loop} pairs three agents. A \textbf{Worker}
attempts the target task. An \textbf{Evaluator} \emph{adversarially}
diagnoses Worker failures and scores performance --- its job is not
passive correctness-grading but active failure-finding, including
probing edge cases the Worker did not test. An \textbf{Evolution} agent
proposes harness modifications conditioned on the full Worker /
Evaluator failure history. At the \textbf{outer level}, the
\textbf{Meta-Evolution Loop} optimises the evolution protocol itself
across a \emph{diverse} set of tasks, learning a protocol that converges
quickly when applied to a new unseen task. The framing the authors give
the contribution is ``automating the design of the automation itself.''
The evaluation spans heterogeneous task families to demonstrate that the
meta-protocol generalises rather than overfitting one benchmark --- and
the inner Worker/Evaluator/Evolution split is the load-bearing primitive
for that generalisation.

\textbf{What GSD already does.} This paper is essentially an external
description of \texttt{gsd-skill-creator}'s purpose. The two-level
structure maps directly onto skill-creator's existing design: the inner
Harness Evolution Loop corresponds to skill-creator's six-step learning
loop applied to a single skill (Observe → Detect → Suggest → Manage →
Auto-Load → Learn/Compose); the outer Meta-Evolution Loop corresponds to
the bounded-learning policy that operates \emph{across} skills, learning
what kinds of corrections actually help library-wide. The roles in
skill-creator's subagent set today are ad-hoc --- implicit Worker (the
skill being applied), implicit Evaluator (the 3-correction-minimum rule
that gates promotion), implicit Evolution (the bounded-learning policy
that absorbs the cross-skill signal). Same shape, no formal naming.

\textbf{The convergence --- external description of skill-creator.} The
mapping is essentially an isomorphism. \textbf{Worker} = the skill being
evolved, applied to the target task. \textbf{Evaluator} = adversarial
diagnosis --- and this is the \emph{upgrade}, because the current
3-correction-min rule is \textbf{passive} (count corrections after the
fact), and the paper's Evaluator is \textbf{active} (probe for failure
modes the Worker did not surface). HB-04 in the v1.49.575 Half-B roster
names this upgrade explicitly as the replacement of 3-correction-min
with the adversarial-Evaluator pattern. \textbf{Evolution} = the
bounded-learning policy applied to skill code, the 20\% rule but
adversarially-applied. \textbf{Meta-Evolution} = the milestone
retrospective accumulating cross-skill protocol changes over time ---
the bounded-learning policy refining itself release over release. The
two-level structure (per-skill iteration vs cross-skill meta-protocol)
is exactly what skill-creator does, and the paper \emph{names} it as a
research direction, which means the design is now nameable and citable
rather than just describable.

\textbf{Compose with AEL \arxivid{2604.21725}.} The Track A audit
pre-relationship between Last Harness and AEL (HB-07's source paper)
yields a productive composition. AEL splits agent learning into
\textbf{fast} (per-episode Thompson-Sampling bandit selection of
retrieval policy) and \textbf{slow} (cross-episode LLM reflection on
failure patterns) timescales. Composed with Last Harness
Worker/Evaluator/Evolution, the structure becomes: \textbf{fast-Worker}
(per-episode skill apply) + \textbf{fast-Evaluator} (per-episode
adversarial check on the Worker output) + \textbf{slow-Evolution}
(cross-episode reflection bandit informing protocol updates). This
composition is what HB-04 (Worker/Evaluator/Evolution role split) and
HB-07 (AEL fast/slow bandit auto-load) implement \emph{together} in
v1.49.575 Half B. HB-04 establishes the role naming and the per-episode
adversarial check; HB-07 supplies the cross-episode reflection
mechanism. The composition is non-trivial because each paper alone is
missing a piece: Last Harness names the roles but doesn't quantify the
timescale split; AEL quantifies the timescales but doesn't name the
roles. Composed, they give skill-creator a fully-named two-timescale,
three-role architecture.

\textbf{Explicit citation target.} Four landing places. (1)
\textbf{\texttt{src/skill-creator/roles/}} (NEW, HB-04 / CS25-16, Phase
808 implementation) --- primary landing: a new module that ships three
role-isolated subagents with explicit Worker / Evaluator / Evolution
contracts and BLOCK authority for the Evaluator over Worker outputs. The
orchestrator brief names this directory as the explicit citation target.
(2) \textbf{\texttt{gsd-mission-management-roles.md}} --- adopt Worker /
Evaluator / Evolution as the formal naming convention for
skill-creator's internal subagents. Today these roles exist informally;
the paper provides published vocabulary. (3)
\textbf{\texttt{.claude/agents/}} --- the existing executor / verifier /
planner subagents fold into Worker / Evaluator / Evolution; document the
mapping so the inherited-role naming is consistent across
\texttt{.claude/agents/} and \texttt{src/skill-creator/roles/}. (4)
\textbf{\texttt{.planning/PROJECT.md}} bounded-learning-policy section
--- cite as the published derivation of the cross-skill meta-protocol.
The Evaluator's adversarial role is a meaningful upgrade only if the
Evaluator has authority to BLOCK Worker outputs --- the role split
degrades to passive-counting otherwise. CS25-16 enforces this with
role-isolation invariants tested explicitly.

\subsection{§6 The four together}

What do the four anchors say \emph{together} about GSD?

\begin{itemize}
\tightlist
\item
  \textbf{The architecture is right.} Skills-as-md three-tier
  (\texttt{21910}) maps point-for-point onto vision-to-mission →
  research-mission-generator → gsd-skill-creator. The pipeline's pieces
  --- built one at a time across many milestones --- line up the way a
  population-genetics workflow team built their system from scratch in
  2026.
\item
  \textbf{The safety doctrine is right.} GROUNDINGmd Hard Constraints +
  Convention Parameters (\texttt{21744}) is the published independent
  rediscovery of CLAUDE.md + Safety Warden BLOCK. The
  override-the-user-prompt commitment that the GSD codebase has treated
  as architecturally normal is what a proteomics governance team
  independently arrived at as the only defensible posture for
  safety-critical fields.
\item
  \textbf{The cognitive substrate is right.} The Root Theorem of Context
  Engineering (\texttt{20874}) formalises the bounded-tape framing ---
  two axioms, five consequences --- and each consequence maps to a
  shipped GSD architectural decision (window-size-independent caps,
  fidelity-threshold gating, accumulate--compress--rewrite--shed
  milestone lifecycle, fresh-context external verification). The
  substrate has now received a formal derivation.
\item
  \textbf{The evolutionary loop is right.} Last Harness (\texttt{21003})
  names skill-creator's Worker/Evaluator/Evolution pattern explicitly
  and is essentially an external description of what
  \texttt{gsd-skill-creator} does. Composed with AEL, it gives the loop
  a published two-timescale, three-role architecture.
\end{itemize}

When four independent peer-reviewed works arrive at GSD's load-bearing
patterns from four different problem domains --- population genetics
(\texttt{21910}), proteomics (\texttt{21744}), abstract context theory
(\texttt{20874}), agentic harness research (\texttt{21003}) --- the
patterns are not idiosyncratic. They are the engineering equivalent of
natural laws. Different watersheds; same trail. The seams the deer have
made walking the trail across multiple drainages are visible from above,
and they all run the same way because the contour is the same. v1.49.575
cites them back, and Half B implements what's adoptable.

\subsection{§7 Calibration}

It is important to distinguish two registers in which this report can be
read.

\textbf{Productive analogy at the design-discipline level.} This is the
register the report stays in. At the design-discipline level, GSD and
these four papers share design intent and structural form: same role
split, same overrides, same lifecycle, same gate placement. This is a
strong external-citation register --- \emph{``these papers independently
arrive at the same design''} is exactly what it says, and the words
``structural form'' and ``design intent'' do the right work without
overpromising. This is the register the v1.50 retrospective Half A's
convergent-discovery section will use.

\textbf{Theorem-level equivalence.} This report does \textbf{not} claim
theorem-level equivalence. We do not claim that GSD's vision-to-mission
is \emph{isomorphic} to the paper's semantic layer in a
category-theoretic sense; we do not claim Hard Constraints and Safety
Warden BLOCK are equal as formal objects; we do not claim the Root
Theorem's two-axiom derivation is \emph{the} derivation of CAPCOM policy
in the sense of a single uniqueness theorem. Theorem-level claims would
require formal proof and a shared formalism --- neither is present here,
and neither is needed for the citations the report supports.

The v1.49.574 retrospective's framing applies verbatim and is reproduced
here for fidelity: \emph{``the rhyme is a productive analogy at the
design-discipline level, not a theorem; it sits in the same register as
`the Amiga Principle compounds twice' and should be presented externally
with that calibration intact.''} The four-anchor convergence in this
report belongs in that same register --- same calibration, same
external-citation discipline. The Amiga Principle, the SIGReg /
megakernel rhyme, and the four CS25 anchors are three instances of the
same kind of finding: load-bearing GSD patterns receiving independent
external derivation, presented as productive analogies, not theorems.

\subsection{§8 Citation map for downstream
waves}

Phase 805 will draft CLAUDE.md and SKILL.md additions; the citation form
below feeds those drafts directly. \emph{Foundational} = primary
citation in the document's design-rationale section; \emph{supporting} =
a citation that strengthens an existing claim; \emph{direct} = the
document is a direct response to or implementation of the paper.

\begin{longtable}[]{@{}
  >{\raggedright\arraybackslash}p{(\linewidth - 4\tabcolsep) * \real{0.3333}}
  >{\raggedright\arraybackslash}p{(\linewidth - 4\tabcolsep) * \real{0.3333}}
  >{\raggedright\arraybackslash}p{(\linewidth - 4\tabcolsep) * \real{0.3333}}@{}}
\toprule\noalign{}
\begin{minipage}[b]{\linewidth}\raggedright
Anchor (arXiv)
\end{minipage} & \begin{minipage}[b]{\linewidth}\raggedright
GSD doc / module
\end{minipage} & \begin{minipage}[b]{\linewidth}\raggedright
Citation form
\end{minipage} \\
\midrule\noalign{}
\endhead
\bottomrule\noalign{}
\endlastfoot
\arxivid{2604.21910} Skills-as-md three-layer & CLAUDE.md ``Architectural
Foundations'' preamble & foundational \\
\arxivid{2604.21910} & \texttt{docs/skill-architecture/three-tier-pipeline.md}
(NEW, Phase 805) & foundational \\
\arxivid{2604.21910} & \texttt{docs/cartridge/SKILL.md} authoring guide &
supporting (validates markdown-skill format) \\
\arxivid{2604.21910} & \texttt{.college/README.md} College of Knowledge
architecture description & supporting (terminology import: semantic /
deterministic / knowledge) \\
\arxivid{2604.21744} GROUNDINGmd & CLAUDE.md ``Important Notes'' --- Hard
Constraints / Convention Parameters & foundational \\
\arxivid{2604.21744} & every \texttt{SKILL.md} template (cartridge authoring
guide) & foundational (Hard-vs-Convention split structure) \\
\arxivid{2604.21744} & \texttt{.claude/skills/security-hygiene/SKILL.md} Safety
Warden BLOCK & direct (override-the-user-prompt rule) \\
\arxivid{2604.21744} & \texttt{src/cartridge/linter/structural-completeness.ts}
(HB-05 / CS25-17) & direct (linter schema adoption) \\
\arxivid{2604.20874} Root Theorem &
\texttt{docs/substrate/bounded-tape-root-theorem.md} (NEW, Phase 805) &
foundational \\
\arxivid{2604.20874} & \emph{The Space Between} information-theory chapter &
foundational \\
\arxivid{2604.20874} & CLAUDE.md design-rationale section & supporting \\
\arxivid{2604.20874} & \texttt{.claude/skills/session-awareness/SKILL.md} + every
CAPCOM-gate policy doc & direct (C3 fidelity-threshold gating + C5
external verification gate) \\
\arxivid{2604.21003} Last Harness & \texttt{src/skill-creator/roles/} (NEW, HB-04
/ CS25-16, Phase 808) & direct \\
\arxivid{2604.21003} & \texttt{gsd-mission-management-roles.md} & direct (naming
convention adoption) \\
\arxivid{2604.21003} & \texttt{.claude/agents/} documentation & supporting
(executor/verifier/planner ↔ Worker/Evaluator/Evolution mapping) \\
\arxivid{2604.21003} & \texttt{.planning/PROJECT.md} bounded-learning-policy
section & supporting (meta-evolution derivation) \\
\arxivid{2604.21003} + \arxivid{2604.21725} (composed) &
\texttt{src/skill-creator/auto-load/bandit.ts} (HB-07) & direct
(fast/slow timescale split inside the three-role split) \\
\end{longtable}

This table is the citation contract for Phase 805. Each row is one diff
to land; each diff cites the indicated paper in the indicated form.
CS25-09 verifies that the CLAUDE.md diff cites all four anchors. CS25-17
verifies the structural-completeness linter implements the GROUNDINGmd
two-class taxonomy. CS25-16 verifies the Worker/Evaluator/Evolution role
split. The seams are mapped; what remains is to walk them.


\clearpage

\section{Cross-Module Relationship Matrix (Phase 800)}

\textit{This appendix embeds the Phase 800 narrative deliverable from \texttt{.planning/missions/cs25-26-sweep/work/synthesis/xmod.md} followed by the full 81-row machine-readable table from \texttt{xmod.csv}. It satisfies CS25-04 (cross-module relationship matrix). 81 connections were classified into seven \texttt{connection\_type} values across six modules.}

\textbf{Phase:} 800 \textbf{Date:} 2026-04-25 \textbf{Author:} synthesis
agent (Track A+B+C cross-cutting) \textbf{Inputs:} 55 ADRs (M1--M6) · 3
track audits (A/B/C) · \texttt{xmod-seed.md} 20-row pre-relationship
table · intake \texttt{cs25\_gsd\_mission.tex}

\begin{center}\rule{0.5\linewidth}{0.5pt}\end{center}

\subsection{§1 Methodology}

Connections were identified by a four-pass sweep over the 55 ADRs and
three track audits. Pass 1 imported the 20 pre-relationships from
\texttt{xmod-seed.md} verbatim and tagged them as the seed set. Pass 2
walked every ADR's \emph{Cross-Module Links} section, recording each
pairwise reference and classifying it into one of the seven
\texttt{connection\_type} values (\texttt{composes-with},
\texttt{validates}, \texttt{extends}, \texttt{contradicts},
\texttt{cross-validates}, \texttt{implements},
\texttt{same-pattern-different-domain}). Pass 3 absorbed the audit-log
``additional connections'' sections from Track A (M1+M2), Track B
(M3+M4), and Track C (M5+M6); these surface composition patterns the
per-ADR sections did not always articulate, especially the AEL × Last
Harness composition and the M3 MCP triple-gate. Pass 4 detected dense
subgraphs by clustering on shared papers and confirmed the four expected
clusters (M2 LoRA, M3 MCP-attack, M5 verification, M6
cultural-sensitivity). Strength was assigned by triangulating
sweep-recommended adoption tier with explicit Track-audit elaboration: a
\texttt{strong} link is one that gates a Phase 805--811 ticket or
carries an explicit composition statement in two places; \texttt{medium}
is one supported by a single ADR cross-reference; \texttt{weak} is
bibliographic-only. No connections rated \texttt{contradictory}
survived; the closest tension was the \arxivid{2604.21927} regime-fixed
methodological constraint, which is rendered as \texttt{validates}
against every M2 LoRA paper because it sets the comparison protocol they
all must satisfy. Zero \texttt{contradicts} rows is itself a finding:
the corpus is internally coherent.

\subsection{§2 Top cross-cutting
recommendations}

The following nine cross-cutting patterns each span ≥2 modules and are
prioritized by impact × effort.

\subsubsection{CC-1. Three-paper MCP triple-gate as Phase 805 baseline
harness}

\textbf{Description:} Any planning-bridge MCP design change must satisfy
three papers jointly: \arxivid{2604.21477} MCP Pitfall Lab (six-attack
family taxonomy), \arxivid{2604.20994} Function Hijacking
(context-agnostic invariant), and \arxivid{2604.21816} Tool Attention
(lazy-schema-loading). Track B's audit treats these as a single
triple-gate rather than three pairwise edges. The Phase 805 baseline
harness must instantiate all six attack families from
\arxivid{2604.21477} as test cases, with \arxivid{2604.20994}'s
context-agnostic-hijack pattern as the seventh case, and Tool
Attention's ISO score gating as the runtime mitigation that the harness
exercises. \textbf{Modules:} M3, M4. \textbf{Papers:} \arxivid{2604.21477},
\arxivid{2604.20994}, \arxivid{2604.21816}. \textbf{Target subsystem:}
\texttt{src/orchestration/tool-attention/} +
\texttt{src/safety/mcp-attack-harness/}. \textbf{Effort:} Medium (one
engineering ticket already drafted as HB-01; harness needs new test
fixtures).

\subsubsection{CC-2. SKILL.md authoring discipline upgrade (Phase 810 +
811)}

\textbf{Description:} Two papers ship together as a single discipline
upgrade: \arxivid{2604.21090} Five-Principle Structural Quality Gaps
(HB-05) and \arxivid{2604.21505} Orchid (HB-06). The first supplies a
structural-completeness linter; the second supplies a four-type
ambiguity taxonomy. Both run as cartridge-promotion linters at
\texttt{src/cartridge/linter/}. Default-off via
\texttt{cs25-26-sweep.structural-completeness-lint} and the parallel
Orchid flag, both promoted to default-on after the existing skill corpus
is audited. The Track-C audit elevates this to a ``spec-discipline
cluster'' together with the M3 GROUNDINGmd Hard Constraints ---
\arxivid{2604.21744}'s convergent-discovery framing is the validation
that the five principles mirror existing CSO discipline rather than
introducing a new norm. \textbf{Modules:} M3, M5, M1. \textbf{Papers:}
\arxivid{2604.21090}, \arxivid{2604.21505}, \arxivid{2604.21744}. \textbf{Target subsystem:}
\texttt{src/cartridge/linter/}. \textbf{Effort:} Medium (linter rules
are mechanically checkable; corpus audit is the long pole).

\subsubsection{CC-3. AEL × Last Harness composed two-loop
pattern}

\textbf{Description:} AEL's fast/slow timescale split
(\arxivid{2604.21725}) composes with Worker/Evaluator/Evolution role
split (\arxivid{2604.21003}) into a single architectural pattern:
fast-Worker (per-episode action), fast-Evaluator (per-episode bandit
critique), slow-Evolution (cross-skill reflection at AEL's slow
timescale). This composition is the headline cross-module finding of the
sweep --- neither paper specifies the other's loop, but their
compositional product is a coherent two-level meta-learning pattern that
GSD's chipset architecture can host without any new infrastructure.
\textbf{Modules:} M1, M2. \textbf{Papers:} \arxivid{2604.21003}, \arxivid{2604.21725}.
\textbf{Target subsystem:} \texttt{src/orchestration/two-loop/} (new) +
chipset Haiku-tier slow-Evolution sidecar. \textbf{Effort:} High
(orchestration touch; CAPCOM HARD GATE; flag-gated default-off).
Elaboration in §4.

\subsubsection{CC-4. Bounded-tape × bounded-learning unified
foundation}

\textbf{Description:} Root Theorem \arxivid{2604.20874} is the formal
foundation for two seemingly different bounded mechanisms: (a) the
20\%/3-correction caps in M2 Retrofit \arxivid{2511.11439} (selective
bounded updates) and (b) the gist-compression /
accumulate-compress-rewrite-shed lifecycle from M6 Gist Sparse Attention
\arxivid{2604.20920}. The same C2 (signal/token-count independence) and
C4 (compression-step-as-instance) corollaries appear in both. Adopting
Root Theorem as the M1 anchor citation gives both M2 Retrofit and M6
Gist a shared formal foundation that can be cited rather than
reasserted. \textbf{Modules:} M1, M2, M6. \textbf{Papers:} \arxivid{2604.20874},
\arxivid{2511.11439}, \arxivid{2604.20920}. \textbf{Target subsystem:}
\texttt{docs/architecture/root-theorem.md} (new substrate doc).
\textbf{Effort:} Low (citation-only). Elaboration in §4.

\subsubsection{CC-5. M2 LoRA dense subgraph as Silicon Layer benchmark
suite}

\textbf{Description:} Five M2 papers form a tight cluster around adapter
design / parameter-efficiency / multi-tenant deployment:
\arxivid{2604.21571} Separable Expert (T2 architecture),
\arxivid{2604.21901} GiVA (T1 benchmark), \arxivid{2509.18629} HyperAdapt
(T1 benchmark), \arxivid{2604.21905} LoRA Redux (TRACK survey), and
\arxivid{2604.21330} TGR-MoE (TRACK pattern). The methodological
constraint from \arxivid{2604.21927} (regime-fixed comparison) applies to
all five as a single comparison protocol. The cluster is the natural
benchmark target list for the RTX 4060 Ti Silicon Layer.
\textbf{Modules:} M2 (intra-cluster dense). \textbf{Papers:} 5 + 1
constraint. \textbf{Target subsystem:}
\texttt{examples/chipsets/silicon-layer/lora-bench/}. \textbf{Effort:}
High (RTX-class compute, but bounded by paper-specified protocols).
Elaboration in §3.

\subsubsection{CC-6. Cross-session threat correlator (Memory Arena
consumer)}

\textbf{Description:} Three M3 papers compose into a persistent
cross-session threat correlator: \arxivid{2604.20911} Omission Decay
(intra-session STD measurement), \arxivid{2604.21131} CSTM-Bench
(cross-session generalization, 50\% recall loss baseline), and
\arxivid{2601.18491} AgentDoG (record schema). Together they specify why
both intra- and inter-session re-injection events must persist to Grove
and be reasoned over by a Safety Warden module at
\texttt{src/safety/cross-session/}. The pattern composes with Memory
Arena M9--M13's content-addressed Grove records as the persistence
layer. The cross-session work GSD does (Seattle 360 / NASA 720
long-horizon missions) is exactly the regime where this attack class
becomes operational. \textbf{Modules:} M3, plus consumer in M5/M6
architecture. \textbf{Papers:} \arxivid{2604.20911}, \arxivid{2604.21131}, \arxivid{2601.18491}.
\textbf{Target subsystem:} \texttt{src/safety/cross-session/}.
\textbf{Effort:} Medium (architecture lands at v1.49.575, runtime defers
to v1.49.576+).

\subsubsection{CC-7. M5 verification triad as v1.50 Half B citation
backbone}

\textbf{Description:} Three M5 papers form a triad for The Space Between
Half B chapter validation: \arxivid{2511.23159} (probable-to-provable
rhetoric), \arxivid{2604.18792} DSLTrans (cutoff theorem mechanism), and
\arxivid{2604.21187} Doubly Saturated Ramsey (SAT+LLM+Lean concrete
precedent). Rhetoric + mechanism + precedent is the right citation chain
for any proof-companion document. The Track-C audit elevates this to the
citation triad backbone for the v1.50 chapter validation. Pairs with
\arxivid{2604.21505} Orchid for spec-side ambiguity reduction.
\textbf{Modules:} M5 (intra-cluster). \textbf{Papers:} 3 + Orchid.
\textbf{Target subsystem:} \texttt{docs/the-space-between/half-b.md}
citation chain. \textbf{Effort:} Low (citation chain only).

\subsubsection{CC-8. M4 four-paper retrieval+synthesis pipeline (College
of
Knowledge)}

\textbf{Description:} Four M4 papers compose into a single
retrieval-and-synthesis pipeline: \arxivid{2510.26615} SARA (hybrid RAG
retrieval), \arxivid{2604.20848} MATRAG (four-agent
transparency-scoring), \arxivid{2604.21444} HiCrew (Hybrid Tree dynamic
role-selection), \arxivid{2508.21720} PosterForest (Poster Tree
synthesis). All four are ADOPT-design-only at v1.49.575; v1.49.576+
implementation. The Track-B audit specifies the dependency graph: MATRAG
depends on SARA; HiCrew dispatches among the MATRAG agents; PosterForest
is the synthesis complement to HiCrew's retrieval. Cross-link with M6's
\arxivid{2601.10003} SocraticKG for the KG-construction prior.
\textbf{Modules:} M4, M6. \textbf{Papers:} 4 + 1. \textbf{Target
subsystem:} \texttt{src/skills/research-mission-generator/v2/} +
\texttt{.college/retrieval/}. \textbf{Effort:} High (v1.49.576+ work;
design-and-spec only at v1.49.575).

\subsubsection{CC-9. M6 cultural-sensitivity cluster as joint Safety
Warden
surface}

\textbf{Description:} Three M6 papers carry SC-IND, SC-MED, and SC-SOV
gates respectively: \arxivid{2604.12710} LASA (semantic-bottleneck
anchoring), \arxivid{2604.20996} AFRILANGTUTOR (sovereignty-gated
low-resource), \arxivid{2604.21352} CARE (counselor-alignment,
assist-not-replace), with \arxivid{2604.21751} (Japanese-culture bias /
CROQ slices) supplying the empirical bias measurement. The cluster
shares a permission-neutral / sovereignty-attentive discipline. CROQ
slices feed the same Safety Warden surface as \arxivid{2604.21308}
CI-Work contextual-integrity benchmark. \textbf{Modules:} M6
(intra-cluster, plus M3 Safety Warden consumer). \textbf{Papers:} 4 +
CI-Work. \textbf{Target subsystem:}
\texttt{src/safety/cultural-sensitivity/} + Safety Warden integration.
\textbf{Effort:} Medium-High (sovereignty consultation is hard
prerequisite for AFRILANGTUTOR). Elaboration in §3.

\begin{center}\rule{0.5\linewidth}{0.5pt}\end{center}

\subsection{§3 Dense subgraphs}

\subsubsection{\texorpdfstring{M2 LoRA cluster (\arxivid{2604.21571} /
\arxivid{2604.21901} / \arxivid{2509.18629} / \arxivid{2604.21905} /
\arxivid{2604.21330} + \arxivid{2604.21927}
constraint)}{M2 LoRA cluster (\arxivid{2604.21571} / \arxivid{2604.21901} / \arxivid{2509.18629} / \arxivid{2604.21905} / \arxivid{2604.21330} + \arxivid{2604.21927} constraint)}}

The five LoRA papers cluster around a shared design principle:
\textbf{parameter-efficient multi-tenant adaptation under a regime-fixed
comparison protocol}. Separable Expert (\arxivid{2604.21571}) is the
multi-tenant deployment frame (per-user proxy, per-domain expert pool);
GiVA (\arxivid{2604.21901}) and HyperAdapt (\arxivid{2509.18629}) are the
parallel rank-reduction / rank-augmentation candidates; LoRA Redux
(\arxivid{2604.21905}) is the bibliographic survey that catalogues both
axes; TGR-MoE (\arxivid{2604.21330}) is the dispatch-level pattern
reference. The \arxivid{2604.21927} regime-fixed paper sits outside this
five-paper set as a methodological constraint that binds them: any
pairwise comparison among them must be evaluated under controlled regime
variables (rank, dataset, training schedule), or the comparison itself
is uninformative. The shared design principle is \emph{adapter
compactness as a multi-tenant deployment axis}. The cluster's coherence
is structural: every paper is reachable from every other in ≤2 hops, and
the methodological constraint is the unique paper that touches all five.

\subsubsection{\texorpdfstring{M5 verification cluster
(\arxivid{2511.23159} / \arxivid{2604.18792} / \arxivid{2604.21187} /
\arxivid{2604.21570} / \arxivid{2604.21598} / \arxivid{2604.21505} /
\arxivid{2603.01170} /
\arxivid{2604.14709})}{M5 verification cluster (\arxivid{2511.23159} / \arxivid{2604.18792} / \arxivid{2604.21187} / \arxivid{2604.21570} / \arxivid{2604.21598} / \arxivid{2604.21505} / \arxivid{2603.01170} / \arxivid{2604.14709})}}

The eight M5 papers cluster around a shared design principle:
\textbf{bounded provable verification as a complement to probable
behaviour}. Three of the eight (probable-to-provable, DSLTrans, Doubly
Saturated Ramsey) form the rhetoric+mechanism+precedent backbone;
SpecSyn and DryRUN form the synthesis complement (spec-side /
test-side); Orchid supplies the four-type ambiguity taxonomy that gates
spec-quality; ATLAS and HWE-Bench are the hardware-domain extensions.
The shared principle is that verification need not be all-or-nothing ---
a bounded, decidable, decomposable verification surface has utility even
when full verification is intractable. This is the design principle the
Track-C audit calls out as the citation triad backbone for The Space
Between Half B.

\subsubsection{\texorpdfstring{M6 cultural-sensitivity cluster
(\arxivid{2604.12710} / \arxivid{2604.20996} / \arxivid{2604.21352} /
\arxivid{2604.21751})}{M6 cultural-sensitivity cluster (\arxivid{2604.12710} / \arxivid{2604.20996} / \arxivid{2604.21352} / \arxivid{2604.21751})}}

The four cultural-sensitivity papers cluster around a shared design
principle: \textbf{permission-neutral, sovereignty-attentive engagement
with cultural and clinical contexts}. LASA's semantic-bottleneck is the
technical defense layer; AFRILANGTUTOR's sovereignty gate is the
procedural defense; CARE's counselor-alignment is the role-boundary
defense; CROQ's bias measurement is the evaluation surface. All four
share a non-claiming posture --- the system assists rather than
replaces, defers to consulted authority on identity claims, and surfaces
bias rather than asserting unbiased operation. The Track-C audit
confirms SC-IND, SC-MED, and SC-SOV gates passed throughout. The
cluster's coherence is normative: it represents the safety-and-respect
column of the ADR matrix in a way no other module mirrors, and the
cluster's outputs flow into a single Safety Warden surface alongside
\arxivid{2604.21308} CI-Work.

\begin{center}\rule{0.5\linewidth}{0.5pt}\end{center}

\subsection{§4 Composition patterns}

\subsubsection{Composition 1 --- AEL × Last Harness (fast-Worker /
fast-Evaluator /
slow-Evolution)}

The Last Harness paper (\arxivid{2604.21003}) defines three roles ---
Worker, Evaluator, Evolution --- without specifying their timescale. AEL
(\arxivid{2604.21725}) defines a fast/slow two-timescale split without
specifying roles. Composing them yields a unique architectural pattern:
the Worker and Evaluator both operate at AEL's fast timescale
(per-episode action and per-episode bandit critique); the Evolution loop
runs at AEL's slow timescale as cross-skill reflection. This is not the
only mapping --- one could place the Evaluator at the slow timescale
instead --- but the fast-Evaluator placement matches the bandit's
per-step feedback structure. The pattern is implementable as a
\texttt{src/orchestration/two-loop/} module with two coroutines, where
the Evolution loop reads the bandit history at a configurable interval
and proposes role-isolation invariants for the next epoch. The CAPCOM
HARD GATE designation is justified because state leakage between roles
defeats the adversarial-Evaluator pattern.

\subsubsection{Composition 2 --- MCP attack-surface joint coverage
(six-attack
baseline)}

\arxivid{2604.21477} MCP Pitfall Lab catalogues six attack families
against MCP servers; \arxivid{2604.20994} Function Hijacking adds the
context-agnostic-invariant case (the same hijack works regardless of
conversation context). Composing them gives a seven-case test harness
for any MCP-touching milestone, and the Phase 805 baseline harness
should instantiate all seven as test fixtures. The Track-B audit
elevates this to a triple-gate when \arxivid{2604.21816} Tool Attention
is added as the runtime mitigation. The harness lives at
\texttt{src/safety/mcp-attack-harness/} and runs as part of the
planning-bridge MCP server's CI pipeline.

\subsubsection{Composition 3 --- Bounded-tape × bounded-learning (Root
Theorem
foundation)}

The Root Theorem (\arxivid{2604.20874}) supplies a formal
information-theoretic statement that two corollaries instantiate: (a)
Retrofit's 20\%/3-correction caps in M2 (\arxivid{2511.11439}) ---
bounded-update mechanism --- and (b) the
accumulate-compress-rewrite-shed lifecycle in M6 Gist Sparse Attention
(\arxivid{2604.20920}) --- bounded-context mechanism. Both are bounded
mechanisms over the same conceptual tape; both inherit the C2
signal/token-count independence corollary; both inherit the C4
compression-step-as-instance corollary. Adopting Root Theorem as the M1
anchor citation gives the M2 and M6 papers a shared formal foundation
that can be cited in their respective ADRs without reasserting the
bound. This is the headline composition of the M1 anchor work --- a
single citation that grounds two otherwise-independent module decisions.

\subsubsection{Composition 4 --- College of Knowledge
retrieval+synthesis
pipeline}

Four M4 papers and one M6 paper compose into a single
retrieval-and-synthesis pipeline for the College of Knowledge:
\arxivid{2510.26615} SARA at the retrieval layer; \arxivid{2604.20848}
MATRAG at the four-agent transparency-scoring layer; \arxivid{2604.21444}
HiCrew at the dynamic role-selection layer; \arxivid{2508.21720}
PosterForest at the synthesis layer; \arxivid{2601.10003} SocraticKG at
the KG-construction priors layer. The dependency graph is: SocraticKG →
SARA → MATRAG → HiCrew → PosterForest, with HiCrew dispatching among
MATRAG's four agents and PosterForest receiving HiCrew's tree as input.
All five are ADOPT-design-only at v1.49.575; the implementation is
scoped for v1.49.576+ as
\texttt{src/skills/research-mission-generator/v2/}. The composition is
significant because it routes through both M4 and M6, validating the
cross-module nature of the retrieval-synthesis problem.

\subsubsection{Composition 5 --- KB construction lineage (SocraticKG →
Gist → Capsules →
Caesar)}

The Track-C audit identifies an M6 chain: \arxivid{2601.10003} SocraticKG
(5W1H QA-driven KG construction) → \arxivid{2604.20920} Gist Sparse
Attention (compression/routing) → \arxivid{2604.20487} Knowledge Capsules
(nonparametric memory) → \arxivid{2604.20855} Caesar (adversarial
KG-driven synthesis). Each paper consumes the previous paper's output:
SocraticKG produces 5W1H tuples, Gist routes a compressed view of them,
Capsules nonparametrically index them, Caesar runs adversarial reasoning
over the indexed KG. The chain forms the College of Knowledge's M6
backbone, with PosterForest as the front-end synthesis layer
(Composition 4) and the College's department/panel hierarchy as the
structural prior (covered in \arxivid{2508.21720} adoption).

\subsubsection{Composition 6 --- Two-game chipset validation (3+1 +
RPS/AGS + 5-axis
matrix)}

Three M4 papers compose into a chipset-validation harness:
\arxivid{2604.21282} 3+1 Strategic Heterogeneous (the
cooperative+adversarial two-game framing), \arxivid{2604.21255} RPS/AGS
(the behavior-similarity metrics), \arxivid{2603.22823} Empirical
Comparison ACPs (the five-axis matrix methodology). The composition
gives GSD an evaluation tooling for whether the chipset's Opus(×N) +
Haiku(verifier) tier separation is producing genuinely complementary
perspectives or has collapsed to behavior-pattern uniformity. The
five-axis matrix supplies the methodology, the AGS metric supplies the
behavior-similarity instrument, and the 3+1 paper supplies the formal
framing for what the metrics are measuring. The composition is
documentation-only at v1.49.575; runtime instrumentation defers to a
future chipset-validation harness.

\begin{center}\rule{0.5\linewidth}{0.5pt}\end{center}

\subsection{§5 Final connection table (rendered from
xmod.csv)}

\begin{longtable}[]{@{}
  >{\raggedright\arraybackslash}p{(\linewidth - 10\tabcolsep) * \real{0.1667}}
  >{\raggedright\arraybackslash}p{(\linewidth - 10\tabcolsep) * \real{0.1667}}
  >{\raggedright\arraybackslash}p{(\linewidth - 10\tabcolsep) * \real{0.1667}}
  >{\raggedright\arraybackslash}p{(\linewidth - 10\tabcolsep) * \real{0.1667}}
  >{\raggedright\arraybackslash}p{(\linewidth - 10\tabcolsep) * \real{0.1667}}
  >{\raggedright\arraybackslash}p{(\linewidth - 10\tabcolsep) * \real{0.1667}}@{}}
\toprule\noalign{}
\begin{minipage}[b]{\linewidth}\raggedright
paper\_a
\end{minipage} & \begin{minipage}[b]{\linewidth}\raggedright
mod\_a
\end{minipage} & \begin{minipage}[b]{\linewidth}\raggedright
paper\_b
\end{minipage} & \begin{minipage}[b]{\linewidth}\raggedright
mod\_b
\end{minipage} & \begin{minipage}[b]{\linewidth}\raggedright
type
\end{minipage} & \begin{minipage}[b]{\linewidth}\raggedright
strength
\end{minipage} \\
\midrule\noalign{}
\endhead
\bottomrule\noalign{}
\endlastfoot
\arxivid{2604.21910} & M1 & \arxivid{2604.21744} & M1 & composes-with & strong \\
\arxivid{2604.21744} & M1 & \arxivid{2604.20874} & M1 & validates & strong \\
\arxivid{2604.21744} & M1 & \arxivid{2601.18491} & M3 & extends & strong \\
\arxivid{2604.20874} & M1 & \arxivid{2604.21927} & M2 & same-pattern-different-domain &
strong \\
\arxivid{2604.20874} & M1 & \arxivid{2604.20920} & M6 & implements & strong \\
\arxivid{2604.21003} & M1 & \arxivid{2604.21725} & M2 & composes-with & strong \\
\arxivid{2604.20300} & M2 & \arxivid{2604.20943} & M2 & same-pattern-different-domain &
medium \\
\arxivid{2604.20300} & M2 & \arxivid{2511.11439} & M2 & extends & strong \\
\arxivid{2604.21282} & M4 & \arxivid{2604.21255} & M4 & validates & strong \\
\arxivid{2604.21794} & M4 & \arxivid{2603.22823} & M4 & cross-validates & strong \\
\arxivid{2604.21816} & M4 & \arxivid{2511.00413} & M4 & composes-with & strong \\
\arxivid{2604.20911} & M3 & \arxivid{2604.21131} & M3 & extends & strong \\
\arxivid{2604.21477} & M3 & \arxivid{2604.20994} & M3 & composes-with & strong \\
\arxivid{2604.21598} & M5 & \arxivid{2604.21505} & M5 & composes-with & strong \\
\arxivid{2511.23159} & M5 & \arxivid{2604.21187} & M5 & validates & strong \\
\arxivid{2604.20920} & M6 & \arxivid{2604.20487} & M6 & composes-with & strong \\
\arxivid{2604.12710} & M6 & \arxivid{2604.21751} & M6 & extends & medium \\
\arxivid{2604.21829} & M6 & \arxivid{2604.21090} & M3 & composes-with & strong \\
\arxivid{2510.18731} & M6 & \arxivid{2601.10003} & M6 & extends & medium \\
\arxivid{2604.20855} & M6 & \arxivid{2508.21720} & M4 & composes-with & strong \\
\arxivid{2604.21571} & M2 & \arxivid{2604.21901} & M2 & composes-with & strong \\
\arxivid{2604.21571} & M2 & \arxivid{2509.18629} & M2 & composes-with & medium \\
\arxivid{2604.21927} & M2 & \arxivid{2604.21901} & M2 & validates & strong \\
\arxivid{2604.21927} & M2 & \arxivid{2509.18629} & M2 & validates & strong \\
\arxivid{2604.21927} & M2 & \arxivid{2604.21905} & M2 & validates & medium \\
\arxivid{2604.21905} & M2 & \arxivid{2604.21571} & M2 & same-pattern-different-domain &
weak \\
\arxivid{2604.21090} & M3 & \arxivid{2604.21505} & M5 & composes-with & strong \\
\arxivid{2604.21744} & M1 & \arxivid{2604.21090} & M3 & validates & strong \\
\arxivid{2604.20874} & M1 & \arxivid{2511.11439} & M2 & implements & strong \\
\arxivid{2604.20874} & M1 & \arxivid{2604.21816} & M4 & validates & strong \\
\arxivid{2604.20911} & M3 & \arxivid{2604.20874} & M1 & validates & medium \\
\arxivid{2604.21477} & M3 & \arxivid{2604.21816} & M4 & composes-with & strong \\
\arxivid{2604.20994} & M3 & \arxivid{2604.21816} & M4 & composes-with & strong \\
\arxivid{2604.21718} & M3 & \arxivid{2604.21003} & M1 & implements & strong \\
\arxivid{2604.21718} & M3 & \arxivid{2601.07262} & M6 & same-pattern-different-domain &
medium \\
\arxivid{2604.21725} & M2 & \arxivid{2604.20943} & M2 & composes-with & medium \\
\arxivid{2604.21330} & M2 & \arxivid{2604.21003} & M1 & same-pattern-different-domain &
medium \\
\arxivid{2604.20300} & M2 & \arxivid{2604.20874} & M1 & implements & strong \\
\arxivid{2604.21282} & M4 & \arxivid{2603.22823} & M4 & validates & strong \\
\arxivid{2604.21794} & M4 & \arxivid{2604.20487} & M6 & same-pattern-different-domain &
medium \\
\arxivid{2511.00413} & M4 & \arxivid{2604.21590} & M4 & composes-with & medium \\
\arxivid{2604.21590} & M4 & \arxivid{2604.21003} & M1 & extends & medium \\
\arxivid{2604.21590} & M4 & \arxivid{2604.21725} & M2 & composes-with & medium \\
\arxivid{2604.21590} & M4 & \arxivid{2604.21718} & M3 & composes-with & medium \\
\arxivid{2604.20848} & M4 & \arxivid{2510.26615} & M4 & composes-with & strong \\
\arxivid{2604.20848} & M4 & \arxivid{2604.21444} & M4 & composes-with & strong \\
\arxivid{2604.20848} & M4 & \arxivid{2601.18491} & M3 & extends & medium \\
\arxivid{2508.21720} & M4 & \arxivid{2604.21444} & M4 & composes-with & strong \\
\arxivid{2604.21444} & M4 & \arxivid{2604.21380} & M6 & composes-with & medium \\
\arxivid{2510.26615} & M4 & \arxivid{2601.10003} & M6 & composes-with & medium \\
\arxivid{2510.26615} & M4 & \arxivid{2604.20920} & M6 & composes-with & medium \\
\arxivid{2604.20848} & M4 & \arxivid{2604.21380} & M6 & composes-with & medium \\
\arxivid{2604.21255} & M4 & \arxivid{2604.21829} & M6 & extends & medium \\
\arxivid{2604.21255} & M4 & \arxivid{2604.21816} & M4 & composes-with & weak \\
\arxivid{2604.21255} & M4 & \arxivid{2511.00413} & M4 & validates & medium \\
\arxivid{2603.01170} & M5 & \arxivid{2604.21477} & M3 & extends & strong \\
\arxivid{2603.01170} & M5 & \arxivid{2604.20994} & M3 & extends & medium \\
\arxivid{2604.18792} & M5 & \arxivid{2604.21570} & M5 & composes-with & medium \\
\arxivid{2604.21570} & M5 & \arxivid{2604.21598} & M5 & composes-with & strong \\
\arxivid{2604.14709} & M5 & \arxivid{2603.01170} & M5 & composes-with & medium \\
\arxivid{2512.03048} & M3 & \arxivid{2511.23159} & M5 & composes-with & strong \\
\arxivid{2512.03048} & M3 & \arxivid{2604.21718} & M3 & implements & medium \\
\arxivid{2512.03048} & M3 & \arxivid{2604.20874} & M1 & validates & strong \\
\arxivid{2604.20920} & M6 & \arxivid{2604.20300} & M2 & composes-with & medium \\
\arxivid{2604.20487} & M6 & \arxivid{2604.20300} & M2 & extends & medium \\
\arxivid{2604.20487} & M6 & \arxivid{2601.10003} & M6 & composes-with & medium \\
\arxivid{2604.20855} & M6 & \arxivid{2604.21380} & M6 & composes-with & medium \\
\arxivid{2604.20855} & M6 & \arxivid{2601.10003} & M6 & composes-with & medium \\
\arxivid{2510.18731} & M6 & \arxivid{2604.21380} & M6 & same-pattern-different-domain &
medium \\
\arxivid{2510.18731} & M6 & \arxivid{2604.21598} & M5 & composes-with & medium \\
\arxivid{2601.07262} & M6 & \arxivid{2604.21725} & M2 & composes-with & medium \\
\arxivid{2601.07262} & M6 & \arxivid{2604.20920} & M6 & composes-with & medium \\
\arxivid{2604.21308} & M6 & \arxivid{2604.21571} & M2 & extends & medium \\
\arxivid{2604.20904} & M6 & \arxivid{2604.21308} & M6 & composes-with & medium \\
\arxivid{2604.20904} & M6 & \arxivid{2604.21352} & M6 & same-pattern-different-domain &
medium \\
\arxivid{2604.21352} & M6 & \arxivid{2604.20996} & M6 & same-pattern-different-domain &
strong \\
\arxivid{2604.12710} & M6 & \arxivid{2604.20996} & M6 & same-pattern-different-domain &
medium \\
\arxivid{2604.21751} & M6 & \arxivid{2604.20996} & M6 & extends & medium \\
\arxivid{2604.21829} & M6 & \arxivid{2604.12710} & M6 & extends & medium \\
\arxivid{2604.21829} & M6 & \arxivid{2604.21380} & M6 & extends & medium \\
\arxivid{2604.21829} & M6 & \arxivid{2604.21477} & M3 & composes-with & strong \\
\end{longtable}

\textbf{Total rows: 81.} All six modules represented. Distribution: M1
papers in 12 connections; M2 in 24; M3 in 24; M4 in 24; M5 in 12; M6 in
38. M6 dominance reflects its 14-paper count and the
cultural-sensitivity cluster's internal density.

\begin{center}\rule{0.5\linewidth}{0.5pt}\end{center}

\subsection{§6 Concerns surfaced for Phase 802 (impact
analysis)}

The following numbered items collect concerns raised in the three track
audits and surfaced during this synthesis. Phase 802 must address each.

\begin{enumerate}
\def\labelenumi{\arabic{enumi}.}
\tightlist
\item
  \textbf{CS25-13 token-budget acceptance gate (Track B)} --- The ≥40\%
  p50 reduction target on Tool Attention depends on tool-set composition
  in the planning-bridge MCP. Phase 802 should require Phase 806 wave 1
  to produce a baseline measurement before locking the acceptance
  threshold.
\item
  \textbf{CS25-15 fail-closed semantics (Track B)} --- STD-calibration's
  fail-closed behaviour when calibration is missing collides with the
  bootstrap path (a fresh deployment has no calibration by definition).
  Phase 808 plan must specify the bootstrap exception explicitly.
\item
  \textbf{HB-08 backlog deferral framing (Track B)} --- Intake §2.4
  frames \arxivid{2604.21255} RPS/AGS as ADOPT, but task instructions
  force backlog deferral. Phase 802 must reconcile the public-framing
  implication and confirm that the metric-definition-only adoption at
  v1.49.575 is the right ceiling.
\item
  \textbf{HB-09 DryRUN architecture sketch quality (Track C)} --- The
  sketch in the ADR is a paragraph; if HB-09 promotes in v1.49.576+, the
  sketch must expand to a concrete module skeleton before implementation
  begins. Phase 802 should set the trigger condition.
\item
  \textbf{CI-Work threat-model publication boundary (Track C)} --- The
  threat-model document straddles the Fox Companies IP rule (specifics
  stay in \texttt{.planning/}). Phase 802 must confirm with maintainers
  before any document moves outside
  \texttt{.planning/fox-companies/foxfiber/}.
\item
  \textbf{Caesar adversarial-refinement loop calibration (Track C)} ---
  The loop's reframing penalty needs calibration on a known-good
  research mission before deployment; if PSC is the chosen
  proving-ground, Phase 802 must schedule a calibration pass before the
  loop runs unattended.
\item
  \textbf{Capsules KV-injection prerequisite (Track C)} --- The
  kernel-hook prerequisite may never land; the TRACK-T3 designation
  acknowledges this and is intentional. Phase 802 should confirm the
  prerequisite-tracking discipline.
\item
  \textbf{Federated-economy threat taxonomy (synthesis finding)} --- Per
  acceptance criterion, any federated-economy recommendation cites
  \arxivid{2604.21829} Black-Box Skill Stealing. CC-9
  (cultural-sensitivity Safety Warden) and CC-2 (SKILL.md authoring)
  both touch federated-skill futures; Phase 802 must verify the citation
  is present in every federated-economy adjacent recommendation that
  lands in v1.49.575 Half B.
\item
  \textbf{\arxivid{2604.21927} regime-fixed comparison protocol
  (synthesis finding)} --- Every M2 LoRA benchmark must be evaluated
  under controlled regime variables. Phase 802 must require the
  v1.49.576+ Silicon Layer benchmark plan to specify regime variables
  explicitly and propagate them to CC-5 (LoRA bench suite) acceptance
  criteria.
\item
  \textbf{MCP triple-gate sequencing (synthesis finding)} --- CC-1's
  three papers must ship together as a single Phase 805 baseline
  harness, but the engineering ticket for Tool Attention (HB-01) is
  currently scoped narrower than the harness. Phase 802 must decide
  whether to widen the HB-01 ticket or open a sibling ticket for the
  harness.
\item
  \textbf{AEL × Last Harness state-isolation invariant (synthesis
  finding)} --- CC-3's two-loop pattern depends on no state leakage
  between Worker, Evaluator, and Evolution. Track-A audit flags this as
  the operational risk; Phase 802 must require the Phase 806 plan to
  enumerate the isolation invariants and gate them under
  \texttt{cs25-26-sweep.two-loop} flag.
\item
  \textbf{M2 LoRA cluster vs Separable Expert tier (Track A)} ---
  Separable Expert is T2 architecture (FoxFiber/FoxCompute
  multi-tenant), but GiVA and HyperAdapt are T1 benchmark targets. Phase
  802 must clarify whether the v1.49.576+ Silicon Layer benchmark suite
  ships before or after the Separable Expert architecture spec.
\item
  \textbf{Track-C cultural-sensitivity gate completeness (synthesis
  finding)} --- SC-IND, SC-MED, SC-SOV all PASSED in Track C, but Phase
  802 should verify that the Phase 805 implementation tickets (when
  written) preserve the same nation-specific naming discipline applied
  in the ADRs.
\item
  \textbf{\arxivid{2604.21090} ↔ \arxivid{2604.21505} ship-together
  discipline (Track B+C)} --- CC-2 requires both linters ship at
  v1.49.575 as a single discipline upgrade. Phase 802 must confirm the
  joint-ship is the implementation reality and not a documentation
  aspiration that gets split at the engineering boundary.
\item
  \textbf{\arxivid{2604.21744} ↔ \arxivid{2604.21090} convergent-discovery
  framing (Track A+B)} --- The five principles are claimed as the formal
  mirror of GSD's existing CSO discipline. Phase 802 should require an
  explicit cross-reference in the Phase 810 ticket so the
  convergent-discovery framing is not lost in implementation.
\end{enumerate}

End of Phase 800 Wave 2.1 synthesis.


\subsection{Full Cross-Module Connection Table (xmod.csv)}

The 81 rows below are the machine-readable export from \texttt{xmod.csv}. Strength values: \texttt{strong} = gates a Phase 805--811 ticket or carries explicit composition statement in two places; \texttt{medium} = single ADR cross-reference; \texttt{weak} = bibliographic-only.

\begin{footnotesize}
\renewcommand{\arraystretch}{1.10}
\begin{longtable}{@{}L{0.09\textwidth} L{0.03\textwidth} L{0.09\textwidth} L{0.03\textwidth} L{0.13\textwidth} L{0.40\textwidth} L{0.06\textwidth}@{}}
\toprule
\textbf{\sffamily Paper A} & \textbf{\sffamily M} & \textbf{\sffamily Paper B} & \textbf{\sffamily M} & \textbf{\sffamily Type} & \textbf{\sffamily Summary} & \textbf{\sffamily Str.} \\
\midrule
\endfirsthead
\toprule
\textbf{\sffamily Paper A} & \textbf{\sffamily M} & \textbf{\sffamily Paper B} & \textbf{\sffamily M} & \textbf{\sffamily Type} & \textbf{\sffamily Summary} & \textbf{\sffamily Str.} \\
\midrule
\endhead
\bottomrule
\endfoot
\arxivid{2604.21910} & M1 & \arxivid{2604.21744} & M1 & \texttt{\small composes-with} & Skill-as-markdown is the governance layer GROUNDINGmd Hard Constraints record onto & strong \\
\arxivid{2604.21744} & M1 & \arxivid{2604.20874} & M1 & \texttt{\small validates} & GROUNDINGmd override-the-prompt invariants are the Root Theorem made formal & strong \\
\arxivid{2604.21744} & M1 & \arxivid{2601.18491} & M3 & \texttt{\small extends} & Hard Constraints flow into AgentDoG where/how/what taxonomy as runtime records & strong \\
\arxivid{2604.20874} & M1 & \arxivid{2604.21927} & M2 & \texttt{\small same-pattern-different-domain} & Bounded-channel theory C2 signal/token-count independence mirrors regime-fixed CL claim independence & strong \\
\arxivid{2604.20874} & M1 & \arxivid{2604.20920} & M6 & \texttt{\small implements} & Root theorem accumulate-compress step instantiated by gist compression & strong \\
\arxivid{2604.21003} & M1 & \arxivid{2604.21725} & M2 & \texttt{\small composes-with} & Worker/Evaluator/Evolution role split composes with AEL fast/slow timescale into single fast-Worker fast-Evaluator slow-Evolution pattern & strong \\
\arxivid{2604.20300} & M2 & \arxivid{2604.20943} & M2 & \texttt{\small same-pattern-different-domain} & FSFM defines what to forget while SCM defines when via NREM/REM consolidation & medium \\
\arxivid{2604.20300} & M2 & \arxivid{2511.11439} & M2 & \texttt{\small extends} & Selective forgetting composed with low-rank-sparse subspace via Retrofit & strong \\
\arxivid{2604.21282} & M4 & \arxivid{2604.21255} & M4 & \texttt{\small validates} & Heterogeneous experts validated by behavior-similarity AGS metric for cooperative-game complementarity & strong \\
\arxivid{2604.21794} & M4 & \arxivid{2603.22823} & M4 & \texttt{\small cross-validates} & Latent vs structured agent communication compared via five-axis matrix & strong \\
\arxivid{2604.21816} & M4 & \arxivid{2511.00413} & M4 & \texttt{\small composes-with} & Tool gating runtime complement to tree-trajectory training time pattern & strong \\
\arxivid{2604.20911} & M3 & \arxivid{2604.21131} & M3 & \texttt{\small extends} & Omission decay intra-session generalized to cross-session correlator threat & strong \\
\arxivid{2604.21477} & M3 & \arxivid{2604.20994} & M3 & \texttt{\small composes-with} & MCP attack-surface joint coverage six-attack-family plus context-agnostic hijack & strong \\
\arxivid{2604.21598} & M5 & \arxivid{2604.21505} & M5 & \texttt{\small composes-with} & Test gen without ground truth pairs with ambiguity taxonomy as spec/test envelope & strong \\
\arxivid{2511.23159} & M5 & \arxivid{2604.21187} & M5 & \texttt{\small validates} & Probable-to-provable rhetoric validated by SAT+LLM+Lean Ramsey precedent & strong \\
\arxivid{2604.20920} & M6 & \arxivid{2604.20487} & M6 & \texttt{\small composes-with} & Gist compression routes over nonparametric capsules in compress-route-unfold architecture & strong \\
\arxivid{2604.12710} & M6 & \arxivid{2604.21751} & M6 & \texttt{\small extends} & Semantic-bottleneck anchoring is the mitigation surface for CROQ-measured cultural bias & medium \\
\arxivid{2604.21829} & M6 & \arxivid{2604.21090} & M3 & \texttt{\small composes-with} & Marketplace skill-stealing threat defended by structural-completeness quality gate & strong \\
\arxivid{2510.18731} & M6 & \arxivid{2601.10003} & M6 & \texttt{\small extends} & RLAAR curriculum instruction-shards composed with SocraticKG 5W1H decomposition & medium \\
\arxivid{2604.20855} & M6 & \arxivid{2508.21720} & M4 & \texttt{\small composes-with} & Adversarial-refinement composes with intermediate-tree pattern in PosterForest+Caesar pipeline & strong \\
\arxivid{2604.21571} & M2 & \arxivid{2604.21901} & M2 & \texttt{\small composes-with} & Separable Expert per-user proxy enabled by GiVA 8x rank reduction for compactness & strong \\
\arxivid{2604.21571} & M2 & \arxivid{2509.18629} & M2 & \texttt{\small composes-with} & Separable Expert domain-expert pool benefits from HyperAdapt O(m+n) high-rank pattern & medium \\
\arxivid{2604.21927} & M2 & \arxivid{2604.21901} & M2 & \texttt{\small validates} & Regime-fixed comparison constraint applies to GiVA benchmarking against vanilla LoRA & strong \\
\arxivid{2604.21927} & M2 & \arxivid{2509.18629} & M2 & \texttt{\small validates} & Regime-fixed comparison constraint applies to HyperAdapt benchmarking & strong \\
\arxivid{2604.21927} & M2 & \arxivid{2604.21905} & M2 & \texttt{\small validates} & Regime-fixed comparison defines regime variables LoRA Redux survey catalogues & medium \\
\arxivid{2604.21905} & M2 & \arxivid{2604.21571} & M2 & \texttt{\small same-pattern-different-domain} & LoRA Redux survey catalogues Separable Expert in adapter design space & weak \\
\arxivid{2604.21090} & M3 & \arxivid{2604.21505} & M5 & \texttt{\small composes-with} & Structural completeness linter ships with Orchid four-type ambiguity taxonomy as SKILL.md authoring upgrade & strong \\
\arxivid{2604.21744} & M1 & \arxivid{2604.21090} & M3 & \texttt{\small validates} & GROUNDINGmd Hard Constraints + Convention Parameters are convergent-discovery analog of five principles & strong \\
\arxivid{2604.20874} & M1 & \arxivid{2511.11439} & M2 & \texttt{\small implements} & Root theorem 20\%/3-correction caps formalized by Retrofit bounded-tape mechanism & strong \\
\arxivid{2604.20874} & M1 & \arxivid{2604.21816} & M4 & \texttt{\small validates} & Root theorem fidelity-threshold framing is formal foundation for 70\% context-fracture threshold & strong \\
\arxivid{2604.20911} & M3 & \arxivid{2604.20874} & M1 & \texttt{\small validates} & Root theorem fidelity-threshold framing is the formal explanation for SRD & medium \\
\arxivid{2604.21477} & M3 & \arxivid{2604.21816} & M4 & \texttt{\small composes-with} & Tool Attention lazy-schema-loading is third leg of joint MCP design change with Pitfall Lab and Function Hijacking & strong \\
\arxivid{2604.20994} & M3 & \arxivid{2604.21816} & M4 & \texttt{\small composes-with} & Function hijacking joint required with Tool Attention for planning-bridge MCP design changes & strong \\
\arxivid{2604.21718} & M3 & \arxivid{2604.21003} & M1 & \texttt{\small implements} & CHAI critique-revise pattern is operational mechanism for Last Harness Evaluator role & strong \\
\arxivid{2604.21718} & M3 & \arxivid{2601.07262} & M6 & \texttt{\small same-pattern-different-domain} & CHAI critique signal mirrors ColorBrowserAgent HITL-feedback-to-domain-knowledge pipeline & medium \\
\arxivid{2604.21725} & M2 & \arxivid{2604.20943} & M2 & \texttt{\small composes-with} & SCM NREM/REM consolidation is candidate slow-timescale reflection mechanism for AEL bandit & medium \\
\arxivid{2604.21330} & M2 & \arxivid{2604.21003} & M1 & \texttt{\small same-pattern-different-domain} & TGR-MoE teacher-student dispatch pattern mirrors Last Harness Worker/Evaluator at harness level & medium \\
\arxivid{2604.20300} & M2 & \arxivid{2604.20874} & M1 & \texttt{\small implements} & Selective forgetting taxonomy implements shed step in Root Theorem accumulate-compress-rewrite-shed cycle & strong \\
\arxivid{2604.21282} & M4 & \arxivid{2603.22823} & M4 & \texttt{\small validates} & Empirical Comparison ACPs five-axis matrix is methodology for evaluating 3+1 chipset against alternatives & strong \\
\arxivid{2604.21794} & M4 & \arxivid{2604.20487} & M6 & \texttt{\small same-pattern-different-domain} & DiffMAS latent KV-cache communication uses same substrate as Knowledge Capsules KV-injection & medium \\
\arxivid{2511.00413} & M4 & \arxivid{2604.21590} & M4 & \texttt{\small composes-with} & AgenticQwen dual flywheels supply curriculum for tree-trajectory training pipeline & medium \\
\arxivid{2604.21590} & M4 & \arxivid{2604.21003} & M1 & \texttt{\small extends} & Last Harness Worker/Evaluator/Evolution split is meta-evolution complement to AgenticQwen dual flywheels & medium \\
\arxivid{2604.21590} & M4 & \arxivid{2604.21725} & M2 & \texttt{\small composes-with} & AEL fast/slow learning is bandit-tier complement to AgenticQwen flywheel discipline & medium \\
\arxivid{2604.21590} & M4 & \arxivid{2604.21718} & M3 & \texttt{\small composes-with} & CHAI pre-X/post-X capture supplies source signal AgenticQwen flywheels consume & medium \\
\arxivid{2604.20848} & M4 & \arxivid{2510.26615} & M4 & \texttt{\small composes-with} & MATRAG Item Analysis depends on SARA hybrid-RAG retrieval as backing layer & strong \\
\arxivid{2604.20848} & M4 & \arxivid{2604.21444} & M4 & \texttt{\small composes-with} & HiCrew dynamic role-selection dispatches among the four MATRAG agents & strong \\
\arxivid{2604.20848} & M4 & \arxivid{2601.18491} & M3 & \texttt{\small extends} & AgentDoG schema embeds MATRAG transparency score in BLOCK records & medium \\
\arxivid{2508.21720} & M4 & \arxivid{2604.21444} & M4 & \texttt{\small composes-with} & PosterForest Poster Tree synthesis complements HiCrew Hybrid Tree retrieval & strong \\
\arxivid{2604.21444} & M4 & \arxivid{2604.21380} & M6 & \texttt{\small composes-with} & IRAP interactive elicitation supplies question-complexity signal HiCrew planning layer reasons over & medium \\
\arxivid{2510.26615} & M4 & \arxivid{2601.10003} & M6 & \texttt{\small composes-with} & SocraticKG 5W1H QA-driven KG construction supplies structural priors SARA vector pool indexes & medium \\
\arxivid{2510.26615} & M4 & \arxivid{2604.20920} & M6 & \texttt{\small composes-with} & Gist Sparse Attention is in-context complement to SARA hybrid retrieval (compressed memory plus selective unfolding) & medium \\
\arxivid{2604.20848} & M4 & \arxivid{2604.21380} & M6 & \texttt{\small composes-with} & IRAP interactive elicitation supplies User Modeling preference profile MATRAG consumes & medium \\
\arxivid{2604.21255} & M4 & \arxivid{2604.21829} & M6 & \texttt{\small extends} & RPS/AGS behavior-similarity is marketplace-skill provenance check for federated-economy threat model & medium \\
\arxivid{2604.21255} & M4 & \arxivid{2604.21816} & M4 & \texttt{\small composes-with} & Tool Attention gated-tool-set affects action-graph the AGS metric measures & weak \\
\arxivid{2604.21255} & M4 & \arxivid{2511.00413} & M4 & \texttt{\small validates} & RPS/AGS metrics validate that subagent trees produce diverse rather than collapsed branches & medium \\
\arxivid{2603.01170} & M5 & \arxivid{2604.21477} & M3 & \texttt{\small extends} & ATLAS CWE-to-assertion pipeline maps onto MCP Pitfall Lab six-attack-family threat catalog & strong \\
\arxivid{2603.01170} & M5 & \arxivid{2604.20994} & M3 & \texttt{\small extends} & ATLAS assertion pipeline must cover Function Hijacking threat surface & medium \\
\arxivid{2604.18792} & M5 & \arxivid{2604.21570} & M5 & \texttt{\small composes-with} & DSLTrans cutoff bound composes with SpecSyn strength-scoring for bounded checks & medium \\
\arxivid{2604.21570} & M5 & \arxivid{2604.21598} & M5 & \texttt{\small composes-with} & SpecSyn synthesizes specs DryRUN synthesizes tests covering both sides of auditor & strong \\
\arxivid{2604.14709} & M5 & \arxivid{2603.01170} & M5 & \texttt{\small composes-with} & HWE-Bench repair-side benchmark complements ATLAS hardware-verification pipeline & medium \\
\arxivid{2512.03048} & M3 & \arxivid{2511.23159} & M5 & \texttt{\small composes-with} & Specification Trap empirical case + probable-to-provable formal-verification companion address spec-difficulty horn & strong \\
\arxivid{2512.03048} & M3 & \arxivid{2604.21718} & M3 & \texttt{\small implements} & CHAI critique-revise is operational mechanism for developmentally-responsive specs & medium \\
\arxivid{2512.03048} & M3 & \arxivid{2604.20874} & M1 & \texttt{\small validates} & Root Theorem supplies formal information-theoretic complement to Specification Trap philosophical argument & strong \\
\arxivid{2604.20920} & M6 & \arxivid{2604.20300} & M2 & \texttt{\small composes-with} & FSFM selective forgetting composes with gist compression to forget what is not gist-routed & medium \\
\arxivid{2604.20487} & M6 & \arxivid{2604.20300} & M2 & \texttt{\small extends} & FSFM selective forgetting becomes capsule-level eviction policy & medium \\
\arxivid{2604.20487} & M6 & \arxivid{2601.10003} & M6 & \texttt{\small composes-with} & SocraticKG 5W1H decomposition is candidate capsule construction pipeline & medium \\
\arxivid{2604.20855} & M6 & \arxivid{2604.21380} & M6 & \texttt{\small composes-with} & IRAP bounded-elicitation is elicitation half complementing Caesar synthesis half & medium \\
\arxivid{2604.20855} & M6 & \arxivid{2601.10003} & M6 & \texttt{\small composes-with} & SocraticKG 5W1H provides KG-construction step Caesar KG-driven reasoning consumes & medium \\
\arxivid{2510.18731} & M6 & \arxivid{2604.21380} & M6 & \texttt{\small same-pattern-different-domain} & RLAAR abstention-at-round-N mirrors IRAP bounded-rounds escalation & medium \\
\arxivid{2510.18731} & M6 & \arxivid{2604.21598} & M5 & \texttt{\small composes-with} & DryRUN trace-synthesis supplies verifiable-accuracy half of curriculum reward & medium \\
\arxivid{2601.07262} & M6 & \arxivid{2604.21725} & M2 & \texttt{\small composes-with} & AEL agent evolving learning consumes domain knowledge ColorBrowserAgent captures & medium \\
\arxivid{2601.07262} & M6 & \arxivid{2604.20920} & M6 & \texttt{\small composes-with} & Gist Sparse Attention is candidate compression layer for ColorBrowserAgent progressive-summarization half & medium \\
\arxivid{2604.21308} & M6 & \arxivid{2604.21571} & M2 & \texttt{\small extends} & Separable Expert composable adapters enable CI-Work multi-tenant patterns & medium \\
\arxivid{2604.20904} & M6 & \arxivid{2604.21308} & M6 & \texttt{\small composes-with} & CI-Work contextual-integrity supplies evaluation surface for Privacy via Fiction normative-simulacrum training & medium \\
\arxivid{2604.20904} & M6 & \arxivid{2604.21352} & M6 & \texttt{\small same-pattern-different-domain} & Privacy via Fiction normative-simulacrum framing pairs with CARE counselor-alignment as assist-not-replace patterns & medium \\
\arxivid{2604.21352} & M6 & \arxivid{2604.20996} & M6 & \texttt{\small same-pattern-different-domain} & CARE permission-neutral discipline mirrors AFRILANGTUTOR sovereignty-attentive discipline & strong \\
\arxivid{2604.12710} & M6 & \arxivid{2604.20996} & M6 & \texttt{\small same-pattern-different-domain} & LASA semantic-bottleneck and AFRILANGTUTOR low-resource pipeline share language-agnostic concern & medium \\
\arxivid{2604.21751} & M6 & \arxivid{2604.20996} & M6 & \texttt{\small extends} & High-resource-language bias finding pairs with AFRILANGTUTOR low-resource focus & medium \\
\arxivid{2604.21829} & M6 & \arxivid{2604.12710} & M6 & \texttt{\small extends} & LASA semantic-bottleneck anchoring defends skill-attribution rules against paraphrase attack & medium \\
\arxivid{2604.21829} & M6 & \arxivid{2604.21380} & M6 & \texttt{\small extends} & IRAP elicitation discipline reduces surface area of skill specifications visible to attackers & medium \\
\arxivid{2604.21829} & M6 & \arxivid{2604.21477} & M3 & \texttt{\small composes-with} & MCP Pitfall Lab metadata-poisoning attack family overlaps with Black-Box Skill Stealing threat model & strong \\
\end{longtable}
\end{footnotesize}


\clearpage

\section{GSD Architectural Impact Analysis (Phase 802)}

\textit{This appendix embeds the Phase 802 deliverable from \texttt{.planning/missions/cs25-26-sweep/work/synthesis/impact.md}. It satisfies CS25-05 (per-subsystem ``what changes / what does not change'' table covers every GSD subsystem affected by an ADOPT decision).}

\subsection{§1 Executive summary}

The v1.49.575 milestone shifts seven GSD subsystems with new code,
references three more with citation-only edits, and leaves the remaining
subsystems untouched. The seven shifting subsystems align with the Half
B T1 roster HB-01..HB-07. Three carry \textbf{CAPCOM HARD GATE}
designation: HB-03 STD calibration (Safety Warden BLOCK timing), HB-04
Worker/Evaluator/Evolution role split (skill-creator orchestration),
HB-07 AEL fast/slow bandit auto-load (skill-space dispatch). All seven
ship default-off behind \texttt{cs25-26-sweep.*} flags so the baseline
is byte-identical when flags are off --- invariant carried forward from
v1.49.570/.571/.572/.574.

The five backlog items (HB-08 RPS/AGS divergence; HB-09 DryRUN test
synthesis; HB-10 IRAP elicitation cap; HB-11 Black-Box Skill Stealing
threat model; HB-12 FSFM 4-class forgetting taxonomy) target subsystems
referenced but not modified in v1.49.575: chipset coherence (HB-08),
skill validation (HB-09), vision-to-mission (HB-10), Grove federated
economy (HB-11), reinforcement series (HB-12). Each carries a documented
promotion trigger.

Aggregate: \textbf{9 subsystems} addressed in §2 (orchestration, safety,
skill-creator, cartridge, Silicon Layer refs, DACP refs, College, Grove,
reinforcement) plus \textbf{3 optional} (memory stack,
cultural-sensitivity gates, vision-to-mission). \textbf{8 prior
milestones} v1.49.561/.568/.569/.570/.571/.572/.573/.574 are touched in
§3. \textbf{15 xmod concerns + 4 audit-level concerns} dispositioned in
§4. Risk register carries 13 rows. The keystone preservation rule ---
\emph{no module proposes replacing CAPCOM, no override of
SC-IND/SC-MED/SC-SOV} --- is reissued as sign-off lines in §6.

\subsection{§2 Per-subsystem impact
table}

\subsubsection{\texorpdfstring{1. Orchestration / planning-bridge MCP
---
\texttt{src/orchestration/}}{1. Orchestration / planning-bridge MCP --- src/orchestration/}}

\textbf{What changes (this milestone, v1.49.575):} - New module
\texttt{src/orchestration/tool-attention/} (HB-01 / CS25-13) per
ADR-\arxivid{2604.21816}: ISO-score selector from sentence embeddings, state-aware
gating, two-phase lazy schema loader. Default-off via
\texttt{cs25-26-sweep.tool-attention}. - Files: \texttt{iso-score.ts},
\texttt{state-gate.ts}, \texttt{lazy-loader.ts},
\texttt{budget-monitor.ts} plus tests under \texttt{\_\_tests\_\_/}. -
Cross-link with M3 attack surface: HB-01 composes with the MCP
triple-gate (\arxivid{2604.21477} Pitfall Lab + \arxivid{2604.20994}
Function Hijacking + \arxivid{2604.21816} Tool Attention) per xmod CC-1;
the harness exercises Tool Attention as runtime mitigation. -
Substrate-doc reference:
\texttt{docs/substrate-references/tool-attention.md} (NEW) names MATRAG
transparency-scoring (\arxivid{2604.20848}) and Tree-Training
(\arxivid{2604.21255} RPS/AGS HB-08 forward) as composing patterns.

\textbf{What does not change:} -
\texttt{src/orchestration/jepa-planner/}, \texttt{retrieval-loop.ts},
\texttt{selector.ts}, \texttt{sol-budget/} are byte-identical. - Tool
dispatch path with flag off is byte-identical to v1.49.574 baseline
(CS25-13 acceptance line). - CAPCOM gate placement on planning-bridge
MCP is preserved --- the Tool Attention loader sits \emph{above} the
dispatch boundary, never below it.

\textbf{Future milestones (v1.49.576+ / v1.50 / later):} - v1.49.576+:
HB-08 RPS/AGS divergence metrics (\arxivid{2604.21255}) lands as
orchestration-side chipset-coherence telemetry. - v1.50: full MCP
six-attack-family runtime harness exits Phase 805 spec into operational
runtime. - v1.50+: Mission Crew Manifest (MATRAG split deferred half
from \arxivid{2604.20848}) rejoins the orchestration retrieval loop.

\subsubsection{\texorpdfstring{2. Safety Warden / CAPCOM ---
\texttt{src/safety/}}{2. Safety Warden / CAPCOM --- src/safety/}}

\textbf{What changes (this milestone, v1.49.575):} - New module
\texttt{src/safety/agentdog/} (HB-02 / CS25-14) per ADR-\arxivid{2601.18491}:
where/how/what diagnostic taxonomy emitted with every BLOCK decision.
Default-off via \texttt{cs25-26-sweep.agentdog-schema}. Haiku-tier
sidecar deployment plan named. - New module
\texttt{src/safety/std-calibration/} (HB-03 / CS25-15, \textbf{CAPCOM
HARD GATE}) per ADR-\arxivid{2604.20911}: per-model omission-decay measurement,
calibration table at \texttt{.planning/safety/std-calibration.json},
re-injection middleware. Default-off via
\texttt{cs25-26-sweep.std-calibration}. - Phase 805 safety-test-harness
expansion adds three categories per CS25-10: STD calibration test,
where/how/what BLOCK schema test, MCP six-attack-family tests covering
at least metadata-poisoning, puppet-server, image-to-tool,
function-hijacking, plus two from the \arxivid{2604.21477} six-attack
baseline. - Cross-session correlator architecture lands per xmod CC-6
--- \texttt{src/safety/cross-session/} directory created with header
docs only; runtime defers to v1.49.576+ (depends on stable AgentDoG
records + STD calibration).

\textbf{What does not change:} - Safety Warden BLOCK authority is
unchanged --- BLOCK is still emitted exactly where it is today; HB-02
only adds a diagnostic record alongside the BLOCK. With flag off, BLOCK
emission is byte-identical. - CAPCOM HARD GATE list remains: HB-03,
HB-04, HB-07. \textbf{No subsystem proposed in v1.49.575 replaces or
bypasses CAPCOM.} This invariant is restated in §6 as a sign-off line. -
\texttt{.claude/skills/security-hygiene/SKILL.md} BLOCK doctrine
continues to override the user prompt --- the GROUNDINGmd citation
(\arxivid{2604.21744}) is the published independent rediscovery of this
rule, not a replacement.

\textbf{Future milestones (v1.49.576+ / v1.50 / later):} - v1.49.576+:
CSTM-Bench (\arxivid{2604.21131}) cross-session correlator runtime ships;
consumes AgentDoG + STD calibration. - v1.49.576+: HB-11 Black-Box Skill
Stealing threat-model document (\arxivid{2604.21829}) lands as Fox
Companies internal artifact under
\texttt{.planning/fox-companies/foxfiber/} (publication boundary
discussed in §4 disposition \#5). - v1.50: CARE counselor-alignment
ADOPT (\arxivid{2604.21352}, M6) --- assist-not-replace framing
referenced in Safety Warden module docs; no new module, citation-only.

\subsubsection{\texorpdfstring{3. skill-creator ---
\texttt{src/skill-creator/}}{3. skill-creator --- src/skill-creator/}}

\textbf{What changes (this milestone, v1.49.575):} - New module
\texttt{src/skill-creator/roles/} (HB-04 / CS25-16, \textbf{CAPCOM HARD
GATE}) per ADR-\arxivid{2604.21003}: Worker / Evaluator / Evolution role-isolated
subagents. Replaces ad-hoc 3-correction-min with adversarial-Evaluator
pattern. Default-off via \texttt{cs25-26-sweep.weler-roles}. - New file
\texttt{src/skill-creator/auto-load/bandit.ts} (HB-07 / CS25-19,
\textbf{CAPCOM HARD GATE}) per ADR-\arxivid{2604.21725}: Thompson-Sampling bandit
+ LLM-reflection slow loop. Default-off via
\texttt{cs25-26-sweep.ael-bandit}. - Composition: HB-04 × HB-07 is xmod
Composition 1 (AEL × Last Harness fast-Worker / fast-Evaluator /
slow-Evolution). The two modules ship together as a single architectural
pattern; flag toggling is independent but the semantic pair is the unit
of design. - \texttt{.claude/agents/} documentation maps existing
executor / verifier / planner subagents onto Worker / Evaluator /
Evolution per convergent-discovery §5 citation map. -
\texttt{.planning/PROJECT.md} bounded-learning-policy section gains a
citation to \arxivid{2604.21003} as the published meta-evolution
derivation (under CS25-09's edit budget).

\textbf{What does not change:} - Six-step loop (Observe → Detect →
Suggest → Manage → Auto-Load → Learn/Compose) shape is preserved; the
Worker/Evaluator/Evolution split \emph{names} what is already there
rather than restructuring. - The 20\% / 3-correction / 7-day cooldown
bounded-learning policy is unchanged --- Root Theorem
(\arxivid{2604.20874}) C1 is the formal foundation for refusing to relax
these caps with bigger context windows. - With both flags off,
skill-creator runs the existing six-step loop byte-identically (CS25-16
+ CS25-19 acceptance). - Existing executor/verifier/planner subagents
keep their current invocation paths; the new role names are additive.

\textbf{Future milestones (v1.49.576+ / v1.50 / later):} - v1.49.576+:
HB-12 FSFM 4-class forgetting taxonomy (\arxivid{2604.20300}) formalizes
the bounded-learning policy with C2/C4 corollaries from Root Theorem; no
new module, taxonomy adoption only. - v1.49.576+: HB-09 DryRUN
(\arxivid{2604.21598}) test-synthesis path lands as skill-validation
harness; sketch promotion gate noted in §4 disposition \#4. - v1.50:
Meta-Evolution loop (Last Harness outer loop) crosses skill boundaries
--- currently single-skill scope only.

\subsubsection{\texorpdfstring{4. cartridge / SKILL.md authoring ---
\texttt{src/cartridge/}}{4. cartridge / SKILL.md authoring --- src/cartridge/}}

\textbf{What changes (this milestone, v1.49.575):} - New file
\texttt{src/cartridge/linter/structural-completeness.ts} (HB-05 /
CS25-17) per ADR-\arxivid{2604.21090}: five-principle SKILL.md validator
(computability / proof-theory / Bayesian / data-classification /
assessment-rubric). Default-off via
\texttt{cs25-26-sweep.structural-completeness-lint} (warning-only when
off). - New file \texttt{src/cartridge/linter/ambiguity.ts} (HB-06 /
CS25-18) per ADR-\arxivid{2604.21505}: lexical / syntactic / semantic / vagueness
type checks. Default-off via \texttt{cs25-26-sweep.ambiguity-lint}. -
New doc \texttt{docs/skill-authoring/ambiguity-checklist.md} ships the
four-type taxonomy as authoring guidance. -
\texttt{docs/cartridge/SKILL.md} authoring guide gains GROUNDINGmd
two-class taxonomy (Hard Constraints / Convention Parameters per
\arxivid{2604.21744}) as a structural-completeness criterion. CS25-17's
negative-fixture corpus includes ``Hard Constraints declared'' cases. -
xmod CC-2 ship-together discipline: HB-05 + HB-06 ship as a single
SKILL.md authoring discipline upgrade (Phase 810 + Phase 811 in
roadmap). Convergent-discovery framing per \arxivid{2604.21744}
cross-reference is added explicitly.

\textbf{What does not change:} - The cartridge promotion pipeline still
routes \texttt{.planning/staging/inbox/} → published; the linter is
\textbf{additive} and gated default-off so existing skills neither break
nor regress. - Existing CSO discipline is unchanged ---
\arxivid{2604.21744}'s convergent-discovery framing is that the five
principles \emph{mirror} CSO, not that they replace it (Track-A audit
cross-link suggestion \#5 elaborates). - Zero false positives on
existing in-tree skills is the CS25-18 acceptance condition; if any
in-tree skill trips, the failure is treated as a fixture defect, not a
baseline regression.

\textbf{Future milestones (v1.49.576+ / v1.50 / later):} - v1.49.576+:
structural-completeness gate flips from warning-only to blocking after
corpus audit completes. Trigger: ≤2\% false-positive rate on the
v1.49.575 corpus. - v1.49.576+: HB-11 Black-Box Skill Stealing
threat-model (\arxivid{2604.21829}) integrates with cartridge linter as
marketplace-contamination defense (xmod CC-9 Federated-Skill futures);
Black-Box Skill Stealing citation gate is restated in §6. - v1.50:
federated-skill DoltHub cartridge schema requires HB-11 threat-model as
pre-rollout citation (Black-Box Skill Stealing pre-rollout-gate, see §4
disposition \#8).

\subsubsection{5. Silicon Layer / LoRA pipeline --- references only this
milestone}

\textbf{What changes (this milestone, v1.49.575):} - New doc
\texttt{docs/substrate-references/lora-cs25-26.md} (NEW, Phase 805)
catalogues the M2 LoRA dense subgraph per xmod §3: \arxivid{2604.21571}
Separable Expert (T2 multi-tenant architecture), \arxivid{2604.21901}
GiVA (T1 benchmark), \arxivid{2509.18629} HyperAdapt (T1 benchmark),
\arxivid{2604.21905} LoRA Redux (TRACK survey), \arxivid{2604.21330}
TGR-MoE (TRACK pattern). Adds \arxivid{2604.21927} regime-fixed
comparison protocol as the binding methodological constraint across all
five. - No new code in \texttt{src/} this milestone --- this is a
documentation-only adoption per Track-A audit's T2 designation for
Separable Expert. - Cross-link: bounded-tape × bounded-learning unified
foundation per xmod CC-4 --- Root Theorem (\arxivid{2604.20874}) becomes
the M1 anchor citation that grounds Retrofit (\arxivid{2511.11439}) C2/C4
corollaries.

\textbf{What does not change:} - Zero changes to running Silicon Layer
code --- there is no Silicon Layer code in \texttt{src/} to change. The
doc is a forward-pointer. - The \texttt{.planning/research-catalog.csv}
Silicon Layer entries are byte-identical. - RTX 4060 Ti VRAM tier
(Memory Arena M6, \texttt{examples/chipsets/silicon-layer/}) is
unchanged.

\textbf{Future milestones (v1.49.576+ / v1.50 / later):} - v1.49.576+:
LoRA bench suite at \texttt{examples/chipsets/silicon-layer/lora-bench/}
instantiates GiVA + HyperAdapt + Separable Expert under the regime-fixed
protocol from \arxivid{2604.21927}. Concern \#9 in §4 names this
regime-spec requirement. - v1.50: full Silicon Layer pipeline ships;
Separable Expert per-user proxy + per-domain expert pool lands as
multi-tenant deployment surface. - v1.50+: TGR-MoE learned-dispatch
milestone (TRACK promotion trigger per Track-A judgment call \#3).

\subsubsection{6. DACP fidelity levels --- references only this
milestone}

\textbf{What changes (this milestone, v1.49.575):} -
\texttt{docs/substrate/semantic-channel.md} (existing from v1.49.572)
gains citation block referencing \arxivid{2603.22823} Empirical
Comparison ACPs five-axis evaluation matrix and \arxivid{2604.21794}
DiffMAS latent-vs-structured comparison. -
\texttt{docs/substrate-references/dacp-evaluation.md} (NEW,
citation-only) names the chipset-validation harness composition per xmod
Composition 6 --- \arxivid{2604.21282} 3+1 Strategic Heterogeneous +
\arxivid{2604.21255} RPS/AGS + \arxivid{2603.22823} Empirical Comparison
ACPs as the joint evaluation tooling.

\textbf{What does not change:} - DACP three-part bundle structure
(deterministic input / agent transform / verifiable output) is preserved
verbatim. v1.49.572 GAP-6 close stands. - Semantic-channel
byte-identical guarantee from v1.49.572 G7 HARD preservation gate is
preserved. - No runtime instrumentation lands this milestone ---
chipset-validation harness defers to v1.49.576+.

\textbf{Future milestones (v1.49.576+ / v1.50 / later):} - v1.49.576+:
chipset-validation harness with the AGS metric (HB-08,
\arxivid{2604.21255}) becomes runtime instrumentation for the Opus×N +
Haiku-verifier tier-separation check. - v1.50: DACP fidelity-level
evaluation matrix runs as a continuous-integration check on every
chipset release.

\subsubsection{\texorpdfstring{7. College of Knowledge ---
\texttt{.college/}}{7. College of Knowledge --- .college/}}

\textbf{What changes (this milestone, v1.49.575):} - Citations only this
milestone. New concept seeds across \texttt{ai-computation} /
\texttt{safety-verification} / \texttt{data-science} departments per the
four anchors (per STATE.md candidate slugs and the convergent-discovery
citation map). - Concept slug additions:
\texttt{bounded-tape-root-theorem},
\texttt{groundingmd-hard-constraints},
\texttt{worker-evaluator-evolution-roles},
\texttt{skills-as-md-three-tier}, \texttt{mcp-six-attack-family},
\texttt{std-calibration-omission-decay},
\texttt{iso-score-tool-attention}, \texttt{four-type-ambiguity},
\texttt{five-principle-structural-completeness}. -
\texttt{.college/README.md} gains terminology import: semantic /
deterministic / knowledge layer naming per \arxivid{2604.21910}.

\textbf{What does not change:} - College department / panel structure is
byte-identical. Rosetta Core unchanged. - \texttt{.college/calibration/}
rubrics unchanged. CSO discipline unchanged (the five principles mirror
CSO per Track-A cross-link \#5; they do not replace it).

\textbf{Future milestones (v1.49.576+ / v1.50 / later):} - v1.49.576+:
College retrieval pipeline lands per xmod CC-8 --- SARA + MATRAG +
HiCrew + PosterForest + SocraticKG five-paper composition at
\texttt{src/skills/research-mission-generator/v2/} and
\texttt{.college/retrieval/}. - v1.50: KG construction lineage per xmod
Composition 5 --- SocraticKG → Gist → Capsules → Caesar; PosterForest as
front-end synthesis.

\subsubsection{\texorpdfstring{8. Grove format / federated economy ---
\texttt{docs/GROVE-FORMAT.md},
\texttt{src/memory/grove-format.ts}}{8. Grove format / federated economy --- docs/GROVE-FORMAT.md, src/memory/grove-format.ts}}

\textbf{What changes (this milestone, v1.49.575):} - Citations only.
\texttt{docs/GROVE-FORMAT.md} gains a forward-pointing note that any
DoltHub federated-economy adoption in v1.49.576+ must cite
\arxivid{2604.21829} Black-Box Skill Stealing as the threat-taxonomy
pre-rollout-gate (per CS25-04 acceptance language). - xmod CC-9
cultural-sensitivity Safety Warden surface integrates with Grove via the
federated-skill-future link (HB-11 backlog).

\textbf{What does not change:} - Layer A content-addressed Grove records
are byte-identical. SkillView, GroveNamespace, SkillCodebase, SkillDiff,
ActivityLog all unchanged. - 299-resource arena
(\texttt{.grove/arena.json}) loadable by any Grove-compatible tool ---
schema unchanged. - Wasteland exclusion remains absolute ---
\texttt{data/chipset/muses/} and \texttt{wasteland} path segment
continue to be blocked.

\textbf{Future milestones (v1.49.576+ / v1.50 / later):} - v1.49.576+:
HB-11 threat-model document lands at
\texttt{.planning/fox-companies/foxfiber/threat-model.md} (Fox Companies
IP --- see §4 disposition \#5). - v1.50: federated-economy DoltHub work
begins. Pre-rollout requirement: HB-11 threat-model citation
\textbf{named} in every federated-economy adjacent recommendation
(CS25-04 + xmod concern \#8 --- restated in §6 sign-off line).

\subsubsection{\texorpdfstring{9. Reinforcement series /
bounded-learning policy --- \texttt{src/reinforcement/},
\texttt{src/eligibility/}, \texttt{src/learnable-k\_h/},
etc.}{9. Reinforcement series / bounded-learning policy --- src/reinforcement/, src/eligibility/, src/learnable-k\_h/, etc.}}

\textbf{What changes (this milestone, v1.49.575):} - Citations only this
milestone. Retrofit (\arxivid{2511.11439}) cited as the parameter-space
form of the 20\% rule (xmod connection paper\_a=\arxivid{2604.20300} ↔
paper\_b=\arxivid{2511.11439} extends-strong). -
\texttt{docs/substrate/bounded-tape-root-theorem.md} (NEW per
convergent-discovery §4) becomes the formal foundation for the
bounded-learning caps under Root Theorem C1+C2+C4.

\textbf{What does not change:} - \texttt{src/reinforcement/},
\texttt{src/eligibility/}, \texttt{src/learnable-k\_h/},
\texttt{src/lyapunov/}, \texttt{src/projection/},
\texttt{src/dead-zone/}, \texttt{src/stochastic/},
\texttt{src/langevin/}, \texttt{src/temperature/},
\texttt{src/embeddings/}, \texttt{src/representation-audit/} --- all
byte-identical. - 20\% cap / 3-correction-min / 7-day cooldown
unchanged. C1 (monotone decay independent of window size) is the formal
justification for refusing to relax these caps even with a 1M-context
model.

\textbf{Future milestones (v1.49.576+ / v1.50 / later):} - v1.49.576+:
HB-12 FSFM 4-class forgetting taxonomy (\arxivid{2604.20300}) formalizes
the policy. Composes with MD-6 representation-audit. - v1.50:
bounded-learning theorem proof completion (carried forward from
v1.49.572 Half B T1d) cites Root Theorem as the empirical anchor.

\subsubsection{\texorpdfstring{10. Memory stack --- \texttt{src/graph/},
\texttt{src/traces/}, \texttt{src/branches/},
\texttt{src/cache/}}{10. Memory stack --- src/graph/, src/traces/, src/branches/, src/cache/}}

\textbf{What changes (this milestone, v1.49.575):} - Citations only.
\arxivid{2604.20920} Gist Sparse Attention and \arxivid{2604.20487}
Knowledge Capsules cited as architectural template --- recursive
gist-of-gist pattern relevant to nested College structure (xmod
Composition 5). - No code change. Sub-citation block added to
\texttt{docs/memory-arena/M9-M13.md} referencing the KB construction
lineage SocraticKG → Gist → Capsules → Caesar.

\textbf{What does not change:} - M1 SMG / M3 Decision-Trace / M4
Branch-Context / M5 cache + KVFlow predictor --- all byte-identical.
Memory Arena M1--M13 stack untouched. - Grove integration
(\texttt{RustArenaSet}, \texttt{ContentAddressedSetStore}) preserved.

\textbf{Future milestones (v1.49.576+ / v1.50 / later):} - v1.49.576+:
Capsules KV-injection prerequisite tracking (TRACK-T3 per Track-C
concern; may never land). - v1.50: Gist Sparse Attention recursive
composition adopted at the College retrieval surface.

\subsubsection{\texorpdfstring{11. Cultural-sensitivity gates ---
\texttt{docs/college/}, citation
surface}{11. Cultural-sensitivity gates --- docs/college/, citation surface}}

\textbf{What changes (this milestone, v1.49.575):} - Citations only.
\arxivid{2604.12710} LASA semantic-bottleneck cited as anchoring
template. \arxivid{2604.20996} AFRILANGTUTOR cited as sovereignty-gate
template (TRACK; SC-SOV is a hard prerequisite --- pre-implementation
consultation required per CS25-11). - \arxivid{2604.21352} CARE
counselor-alignment cited; assist-not-replace framing preserved
exclusively (no diagnostic / therapeutic claims). - xmod CC-9
cultural-sensitivity cluster integrates into Safety Warden surface
alongside \arxivid{2604.21308} CI-Work contextual-integrity benchmark.

\textbf{What does not change --- explicit invariant restatement:} -
\textbf{SC-IND, SC-MED, SC-SOV gates are preserved with no override
proposed.} Every reference to PNW Indigenous nations names them
specifically (Coast Salish, Lummi/Lhaq'temish, Tulalip, Stillaguamish,
Wanapum). No generic ``Indigenous peoples'' phrasing. CS25-11 acceptance
line is binding for Phase 805 publication and all subsequent waves. This
restatement is a sign-off line in §6.

\textbf{Future milestones (v1.49.576+ / v1.50 / later):} - v1.49.576+:
AFRILANGTUTOR PNW adaptation \textbf{gated} on nation-specific
consultation. No implementation begins until consultation completes. -
v1.50: CROQ slice (\arxivid{2604.21751}) bias-measurement instrument
feeds Safety Warden surface; \arxivid{2604.21308} CI-Work
contextual-integrity benchmark runs as joint surface.

\subsubsection{\texorpdfstring{12. vision-to-mission /
research-mission-generator --- skill-pack level
(\texttt{Tibsfox/skills})}{12. vision-to-mission / research-mission-generator --- skill-pack level (Tibsfox/skills)}}

\textbf{What changes (this milestone, v1.49.575):} - Citations only. The
skills-as-md three-layer architecture (\arxivid{2604.21910}) is the
foundational citation per convergent-discovery §2 --- vision-to-mission
= semantic layer; research-mission-generator = deterministic layer;
gsd-skill-creator = knowledge layer. - Caesar adversarial-refinement
loop (\arxivid{2604.20855}) tracked in research-mission-generator design
notes; calibration concern noted in §4 disposition \#6.

\textbf{What does not change:} - Both skills run unchanged in v1.49.575.
No version bump, no new behaviour, no flag. - Three-stage output (Vision
Guide → Research Reference → Mission Spec) preserved.

\textbf{Future milestones (v1.49.576+ / v1.50 / later):} - v1.49.576+:
HB-10 IRAP 5-round elicitation cap (\arxivid{2604.21380}) lands as soft
discipline now (vision-to-mission soft cap), hard CAPCOM-gate escalation
later. Promotion trigger: model-affinity escalation phase opens. -
v1.49.576+: Caesar adversarial-refinement TRACK promotion contingent on
calibration pass against PSC mission (§4 disposition \#6).

\subsection{§3 Cross-milestone linkage
map}

The v1.49.575 work touches \textbf{eight prior milestones} and
\textbf{one parallel track}. Linkage forms: cite (citation-only), extend
(composes with and adds to), compose (paired structural composition),
reference (forward-pointer).

\begin{itemize}
\tightlist
\item
  \textbf{v1.49.561 Living Sensoria} --- \emph{cite + reference}.
  M6/M7/M8 (\texttt{src/sensoria/}, \texttt{src/umwelt/},
  \texttt{src/symbiosis/}) cited in convergent-discovery §4 C5 ---
  umwelt's Markov-blanket boundary = Root Theorem C5
  external-verification-gate. CC-9 cultural-sensitivity Safety Warden
  surface composes with M8 symbiosis ledger. No code change.
\item
  \textbf{v1.49.568 Nonlinear Frontier} --- \emph{cite}.
  Bounded-learning anchor for Root Theorem; NLF-IDs are the external
  mathematical foundation of the 20\% / 3-correction caps now sharpened
  by \arxivid{2604.20874} C1.
\item
  \textbf{v1.49.569 Drift in LLM Systems} --- \emph{compose}. STD
  calibration HB-03 composes with drift detection: STD measurement is
  the prohibition-decay specialization of v1.49.569's drift-detection
  methodology. Phase 808 plan references drift telemetry for the
  bootstrap-exception path (§4 \#2).
\item
  \textbf{v1.49.570 Convergent Substrate} --- \emph{extend}.
  \texttt{synthesis/convergent-discovery.md} is the second-generation
  instance of the canonical pattern v1.49.570 established. Four anchors
  (\arxivid{2604.21910}, \arxivid{2604.21744}, \arxivid{2604.20874},
  \arxivid{2604.21003}) compose with the v1.49.570 convergent block.
  Two-gate / cascade-mcp-defense modules are pre-citations of CC-1 (MCP
  triple-gate).
\item
  \textbf{v1.49.571 LeJEPA / Heuristics-Free} --- \emph{compose}. Last
  Harness Worker/Evaluator/Evolution split = heuristics-free
  skill-creator pattern; the adversarial Evaluator replaces the
  heuristic 3-correction-min. HB-04 + LEJEPA-13 (skill-isotropy) +
  LEJEPA-15 (single-λ audit) form a design-discipline triad.
\item
  \textbf{v1.49.572 Math Foundations} --- \emph{compose (theorem
  stack)}. Root Theorem composes with three v1.49.572 T1 deliverables:
  Coherent Functors (T1a), Semantic Channel DACP (T1c, Root Theorem C2 =
  rate-distortion specialization), Bounded-Learning Theorem (T1d, Root
  Theorem C1 = empirical anchor for 20\% cap proof). xmod CC-4 =
  bounded-tape × bounded-learning unified foundation.
\item
  \textbf{v1.49.573 Upstream Intelligence} --- \emph{extend}. Skilldex
  Conformance Auditor (T1a) compounds with HB-05 Five-Principle linter
  --- same gate, complementary checks. v1.49.573's upstream +
  agentskills.io audit feed the same \texttt{src/cartridge/linter/}
  directory.
\item
  \textbf{v1.49.574 Megakernel} --- \emph{compose (keystone)}. SIGReg ↔
  counter-based megakernel architectural-rhyme composes with the
  bounded-tape Root Theorem keystone --- both
  productive-analogy-at-design-discipline-level findings.
  Convergent-discovery §7 calibration reused verbatim. RAVEN-INTEG
  roster is the pattern HB-01..HB-07 follows.
\end{itemize}

\textbf{Parallel track:} v1.0 PNW Living Systems Taxonomy (phases
620--635) --- \emph{cite}. College concept-seeds feed the PNW taxonomy
hub; cross-references.json gains edges from CS25-26 hub to AAR / SST /
LLM / Chipset.

\subsection{§4 Concerns surfaced +
dispositions}

The 15 concerns from xmod.md §6 plus 4 audit-level concerns are
dispositioned below. Each: concern statement → disposition.

\begin{enumerate}
\def\labelenumi{\arabic{enumi}.}
\item
  \textbf{CS25-13 Tool Attention acceptance threshold (xmod \#1; Track
  B).} ≥40\% p50 reduction depends on tool-set composition. →
  \textbf{Accept with Phase 806 wave-1 baseline.} Phase 806 produces a
  baseline before locking the threshold; if 40\% is unreachable,
  re-scope to ``p50 reduction ≥ baseline × 0.6'' in the Phase 806 plan.
\item
  \textbf{CS25-15 STD fail-closed bootstrap exception (xmod \#2; Track
  B).} Fail-closed collides with bootstrap (fresh deployments have no
  calibration). → \textbf{Defer to Phase 808.} Phase 808 plan specifies
  a calibration-acquisition mode for ≤N turns, after which fail-closed
  activates. CAPCOM HARD GATE pass requires this spec.
\item
  \textbf{HB-08 backlog deferral (xmod \#3; Track B).} Intake §2.4
  framed \arxivid{2604.21255} as ADOPT; task instructions force backlog.
  → \textbf{Defer to retrospective.} Backlog stands;
  \texttt{docs/release-notes/v1.49.575/} retrospective documents the
  divergence (T1 capacity ceiling) so the framing is not silently
  overruled.
\item
  \textbf{HB-09 DryRUN sketch quality (xmod \#4; Track C).} Sketch is
  one paragraph. → \textbf{Defer with trigger condition.} HB-09
  promotion is gated on a ≥800-word architecture-spec ADR in
  \texttt{work/modules/M5/} before the implementing phase opens.
\item
  \textbf{CI-Work threat-model publication boundary (xmod \#5; Track
  C).} Document straddles Fox Companies IP rule. → \textbf{Surface to
  user; defer publication.} Fox Companies IP review required before any
  public artifact moves outside
  \texttt{.planning/fox-companies/foxfiber/}. Abstract attack-taxonomy
  citable; specifics stay private. User-decision item.
\item
  \textbf{Caesar adversarial-refinement calibration (xmod \#6; Track
  C).} Reframing-penalty needs calibration before deployment. →
  \textbf{TRACK only this milestone; defer ADOPT.} Caesar does not run
  autonomously in v1.49.575. If PSC runs in v1.49.576+, schedule a
  calibration pass first.
\item
  \textbf{Capsules KV-injection prerequisite (xmod \#7; Track C).} May
  never land. → \textbf{Accept TRACK-T3 discipline.} Status is
  intentional; no further action.
\item
  \textbf{Federated-economy threat-taxonomy citation gate (xmod \#8).}
  Per CS25-04, federated-economy recs must cite \arxivid{2604.21829}. →
  \textbf{Accept now as pre-rollout gate.} §2 Grove (subsystem 8) and
  cartridge (subsystem 4) blocks name it. §6 sign-off restates: any
  DoltHub work in v1.49.576+ requires the citation before merge.
\item
  \textbf{\arxivid{2604.21927} regime-fixed protocol (xmod \#9).} Every
  M2 LoRA benchmark needs regime variables (rank/dataset/schedule). →
  \textbf{Accept as benchmark-spec requirement.}
  \texttt{docs/substrate-references/lora-cs25-26.md} names variables;
  v1.49.576+ Silicon Layer plan inherits as CC-5 acceptance.
\item
  \textbf{MCP triple-gate sequencing (xmod \#10).} CC-1 three papers
  ship together; HB-01 ticket is narrower. → \textbf{Accept; widen
  HB-01.} Phase 806 ticket widens to include the six-attack-family +
  function-hijacking test-fixture surface (CS25-10); one ticket, no
  siblings.
\item
  \textbf{AEL × Last Harness state-isolation (xmod \#11).} CC-3 depends
  on no Worker/Evaluator/Evolution leakage. → \textbf{Accept; enumerate
  in Phase 808 plan.} Worker writes own state only; Evaluator reads
  Worker output (read-only); Evolution reads aggregated history
  (read-only). Tests verify no cross-role write paths.
\item
  \textbf{M2 LoRA vs Separable Expert tier sequencing (xmod \#12; Track
  A).} SE is T2; GiVA + HyperAdapt are T1 benchmarks. → \textbf{Accept;
  document sequencing.} Benchmark suite ships \textbf{before} SE
  architecture spec; benchmarks inform the spec.
\item
  \textbf{Track-C cultural-sensitivity gate completeness (xmod \#13).}
  SC-IND/SC-MED/SC-SOV PASSED. → \textbf{Accept; carry to Phase 805
  tickets.} Implementation tickets preserve nation-specific naming.
  CS25-11 = verification entry.
\item
  \textbf{\arxivid{2604.21090} ↔ \arxivid{2604.21505} ship-together (xmod
  \#14).} Both linters ship as one discipline upgrade. → \textbf{Accept
  as engineering reality.} §2 subsystem 4 names joint Phase 810 + 811
  ship; phases close together.
\item
  \textbf{\arxivid{2604.21744} ↔ \arxivid{2604.21090} convergent-discovery
  framing (xmod \#15).} Five principles = formal mirror of CSO. →
  \textbf{Accept as Phase 810 cross-reference.} Phase 810 ticket
  includes the cross-ref so the mirror-of-CSO claim is not lost in
  implementation.
\end{enumerate}

\textbf{Audit-level concerns (Tracks A/B/C):}

\begin{enumerate}
\def\labelenumi{\arabic{enumi}.}
\setcounter{enumi}{15}
\item
  \textbf{Wave-1 parallelism via subagent fleets (Track A).} Three
  concurrent fleets (A/B/C) writing ADRs to the same
  \texttt{work/modules/} directory tree --- risk of merge collision. →
  \textbf{Accept; mitigated by per-module subdirectory isolation.} M1+M2
  → Track A; M3+M4 → Track B; M5+M6 → Track C. No cross-track writes.
\item
  \textbf{Track-A LoRA Redux ADOPT-vs-TRACK judgment (Track A judgment
  call \#2).} Survey paper deferred to TRACK. → \textbf{Accept;
  bibliographic citation only.}
\item
  \textbf{Track-B CS25-13 baseline measurement (overlaps xmod \#1).} →
  Same disposition as \#1.
\item
  \textbf{Track-C SC-PARA paraphrase gate held for 22 ADRs.} →
  \textbf{Accept; verification gate passed.} No remediation required.
\end{enumerate}

\subsection{§5 Risk register}

Per-subsystem integration risk aggregated from ADR Risk fields.
Mitigation owner is the phase number that closes the risk.

\begin{longtable}[]{@{}
  >{\raggedright\arraybackslash}p{(\linewidth - 6\tabcolsep) * \real{0.2500}}
  >{\raggedright\arraybackslash}p{(\linewidth - 6\tabcolsep) * \real{0.2500}}
  >{\raggedright\arraybackslash}p{(\linewidth - 6\tabcolsep) * \real{0.2500}}
  >{\raggedright\arraybackslash}p{(\linewidth - 6\tabcolsep) * \real{0.2500}}@{}}
\toprule\noalign{}
\begin{minipage}[b]{\linewidth}\raggedright
Subsystem
\end{minipage} & \begin{minipage}[b]{\linewidth}\raggedright
Aggregate risk
\end{minipage} & \begin{minipage}[b]{\linewidth}\raggedright
Dominant concern
\end{minipage} & \begin{minipage}[b]{\linewidth}\raggedright
Mitigation owner
\end{minipage} \\
\midrule\noalign{}
\endhead
\bottomrule\noalign{}
\endlastfoot
Orchestration / Tool Attention (HB-01) & Medium & Acceptance threshold
dependent on tool-set composition (xmod \#1) & Phase 806 wave 1 \\
Safety / AgentDoG (HB-02) & Low & Schema-additive only; flag-off
byte-identical & Phase 807 \\
Safety / STD calibration (HB-03, \textbf{CAPCOM HARD GATE}) & Medium &
Fail-closed bootstrap exception (xmod \#2); CAPCOM gate timing & Phase
808 + CAPCOM gate \\
skill-creator / WELER roles (HB-04, \textbf{CAPCOM HARD GATE}) &
Medium-High & State-isolation invariants (xmod \#11); CAPCOM gate &
Phase 808 + CAPCOM gate \\
skill-creator / AEL bandit (HB-07, \textbf{CAPCOM HARD GATE}) &
Medium-High & Bandit convergence + reflection-feedback loop; CAPCOM gate
& Phase 812 + CAPCOM gate \\
Cartridge / structural-completeness linter (HB-05) & Low & False
positives on existing in-tree skills (warning-only baseline mitigates) &
Phase 810 \\
Cartridge / ambiguity linter (HB-06) & Low & Author friction on first
publish & Phase 811 \\
Silicon Layer references & Low & Documentation-only this milestone &
Phase 805 \\
DACP references & Low & Documentation-only this milestone & Phase 805 \\
College of Knowledge & Low & Citation-only; concept seeds are additive &
Phase 805 \\
Grove federated economy & Medium (deferred) & HB-11 threat-model
citation gate; publication boundary (concern \#5) & v1.49.576+ \\
Reinforcement series & Low & Citation-only this milestone & Phase 805 \\
Cultural-sensitivity gates & Medium & SC-IND/SC-MED/SC-SOV preservation;
AFRILANGTUTOR sovereignty consultation & Phase 805 + v1.49.576+ \\
\end{longtable}

Three CAPCOM HARD GATES --- HB-03, HB-04, HB-07 --- are noted
explicitly. All three require human authorization at the gate; none can
pass autonomously.

\subsection{§6 Acceptance \& sign-off
lines}

Closing checklist for Phase 803 (integration specs) and Phase 805
(publication):

\begin{itemize}
\tightlist
\item[$\boxtimes$]
  \textbf{All ≥9 subsystems addressed in §2.} Twelve subsystems covered
  (9 required + 3 optional).
\item[$\boxtimes$]
  \textbf{Every milestone affected named in §3.} Eight prior milestones
  (v1.49.561 / .568 / .569 / .570 / .571 / .572 / .573 / .574) plus the
  v1.0 PNW Taxonomy parallel track.
\item[$\boxtimes$]
  \textbf{All 15+ concerns from xmod.md §6 dispositioned in §4.} 15 xmod
  concerns + 4 audit-level concerns = 19 dispositions.
\item[$\boxtimes$]
  \textbf{CAPCOM HARD GATEs at HB-03/HB-04/HB-07 noted in §5 risk
  register.} All three rows carry the explicit CAPCOM HARD GATE
  annotation.
\item[$\boxtimes$]
  \textbf{Cultural-sensitivity gates SC-IND/SC-MED/SC-SOV explicitly
  preserved (no overrides proposed).} §2 subsystem 11 names the
  invariant restatement; §4 disposition \#13 carries it forward to Phase
  805 tickets; CS25-11 is the verification line.
\item[$\boxtimes$]
  \textbf{Federated-economy citation gate (Black-Box Skill Stealing
  threat model, \arxivid{2604.21829}) named as a pre-rollout requirement
  before any DoltHub work in v1.49.576+.} §2 subsystem 8 names the gate;
  §4 disposition \#8 accepts it as pre-rollout.
\end{itemize}

\textbf{Sign-off invariants (restated for Phase 805 publication):}

\begin{enumerate}
\def\labelenumi{\arabic{enumi}.}
\tightlist
\item
  \textbf{No module proposed in v1.49.575 replaces or bypasses CAPCOM.}
  Every Half B T1 module ships default-off; baseline behaviour is
  byte-identical when flags are off.
\item
  \textbf{SC-IND / SC-MED / SC-SOV gates carry no override.} PNW
  Indigenous nations are named specifically (Coast Salish,
  Lummi/Lhaq'temish, Tulalip, Stillaguamish, Wanapum). CARE template
  uses counselor-alignment exclusively. AFRILANGTUTOR PNW adaptation
  requires nation-specific consultation BEFORE any data construction.
\item
  \textbf{\arxivid{2604.21829} Black-Box Skill Stealing citation is a
  pre-rollout gate} for any federated-economy / DoltHub work in
  v1.49.576+.
\item
  \textbf{Publication boundary preserved.} CI-Work threat-model
  specifics stay in \texttt{.planning/fox-companies/foxfiber/}; only the
  abstract attack-taxonomy level may be cited externally. Fox Companies
  IP review required before public artifacts.
\item
  \textbf{CAPCOM HARD GATEs at HB-03/HB-04/HB-07 require human
  authorization.} No autonomous pass.
\end{enumerate}

End of Phase 802 Wave 2.3 GSD architectural impact analysis. CS25-05
satisfied.


\clearpage

\section{Drop-in Integration Specs (Phase 803)}

\textit{This appendix embeds the 12 engineering-ticket-grade specs from \texttt{.planning/missions/cs25-26-sweep/work/synthesis/specs/}. Each spec \textgreater{}=250 words. HB-01..HB-07 feed Phase 806--812 implementation in v1.49.575 Half B; HB-08..HB-12 feed v1.49.576+ backlog. Three specs carry CAPCOM HARD GATE designation. It satisfies CS25-06 (drop-in integration specs for all ADOPT-decision papers).}

\subsection{Spec roster index}
This directory holds engineering-ticket-grade specs for the 12 HB-line
ADOPT decisions surfaced by Wave 1. Each spec ≥250 words
(\textasciitilde300 target). HB-01..HB-07 feed Phase 806--812
implementation in v1.49.575 Half B; HB-08..HB-12 feed v1.49.576+ backlog
(\texttt{gsd-add-backlog} integration). Three specs carry \textbf{CAPCOM
HARD GATE} designation and inherit the gate-pattern template from
v1.49.574 megakernel-substrate
\texttt{src/cartridge/megakernel/adapter-selection-schema.ts}.

\subsubsection{Spec roster}

\paragraph{T1 (v1.49.575 Half B)}

\begin{longtable}[]{@{}
  >{\raggedright\arraybackslash}p{(\linewidth - 14\tabcolsep) * \real{0.1250}}
  >{\raggedright\arraybackslash}p{(\linewidth - 14\tabcolsep) * \real{0.1250}}
  >{\raggedright\arraybackslash}p{(\linewidth - 14\tabcolsep) * \real{0.1250}}
  >{\raggedright\arraybackslash}p{(\linewidth - 14\tabcolsep) * \real{0.1250}}
  >{\raggedright\arraybackslash}p{(\linewidth - 14\tabcolsep) * \real{0.1250}}
  >{\raggedright\arraybackslash}p{(\linewidth - 14\tabcolsep) * \real{0.1250}}
  >{\raggedright\arraybackslash}p{(\linewidth - 14\tabcolsep) * \real{0.1250}}
  >{\raggedright\arraybackslash}p{(\linewidth - 14\tabcolsep) * \real{0.1250}}@{}}
\toprule\noalign{}
\begin{minipage}[b]{\linewidth}\raggedright
HB-id
\end{minipage} & \begin{minipage}[b]{\linewidth}\raggedright
Title
\end{minipage} & \begin{minipage}[b]{\linewidth}\raggedright
Tier
\end{minipage} & \begin{minipage}[b]{\linewidth}\raggedright
Phase
\end{minipage} & \begin{minipage}[b]{\linewidth}\raggedright
Target subsystem
\end{minipage} & \begin{minipage}[b]{\linewidth}\raggedright
Source ADR
\end{minipage} & \begin{minipage}[b]{\linewidth}\raggedright
Acceptance
\end{minipage} & \begin{minipage}[b]{\linewidth}\raggedright
CAPCOM HARD GATE
\end{minipage} \\
\midrule\noalign{}
\endhead
\bottomrule\noalign{}
\endlastfoot
HB-01 & Tool Attention Lazy Schema Loader & T1 & 806 &
\texttt{src/orchestration/tool-attention/} & ADR-\arxivid{2604.21816} & CS25-13 &
--- \\
HB-02 & AgentDoG Where/How/What BLOCK Schema & T1 & 807 &
\texttt{src/safety/agentdog/} & ADR-\arxivid{2601.18491} & CS25-14 & --- \\
HB-03 & Safe-Turn-Depth Calibration & T1 & 808 &
\texttt{src/safety/std-calibration/} & ADR-\arxivid{2604.20911} & CS25-15 &
\textbf{YES} \\
HB-04 & Worker / Evaluator / Evolution Role Split & T1 & 808 &
\texttt{src/skill-creator/roles/} & ADR-\arxivid{2604.21003} (composes
ADR-\arxivid{2604.21725}) & CS25-16 & \textbf{YES} \\
HB-05 & Five-Principle Structural-Completeness Linter & T1 & 810 &
\texttt{src/cartridge/linter/} & ADR-\arxivid{2604.21090} & CS25-17 & --- \\
HB-06 & Four-Type Ambiguity Checklist & T1 & 811 &
\texttt{src/cartridge/linter/} & ADR-\arxivid{2604.21505} & CS25-18 & --- \\
HB-07 & AEL Fast/Slow Bandit Auto-Load & T1 & 812 &
\texttt{src/skill-creator/auto-load/} & ADR-\arxivid{2604.21725} (composes
ADR-\arxivid{2604.21003}) & CS25-19 & \textbf{YES} \\
\end{longtable}

\begin{center}\rule{0.5\linewidth}{0.5pt}\end{center}

\paragraph{T2/T3 backlog (v1.49.576+)}

\begin{longtable}[]{@{}
  >{\raggedright\arraybackslash}p{(\linewidth - 14\tabcolsep) * \real{0.1250}}
  >{\raggedright\arraybackslash}p{(\linewidth - 14\tabcolsep) * \real{0.1250}}
  >{\raggedright\arraybackslash}p{(\linewidth - 14\tabcolsep) * \real{0.1250}}
  >{\raggedright\arraybackslash}p{(\linewidth - 14\tabcolsep) * \real{0.1250}}
  >{\raggedright\arraybackslash}p{(\linewidth - 14\tabcolsep) * \real{0.1250}}
  >{\raggedright\arraybackslash}p{(\linewidth - 14\tabcolsep) * \real{0.1250}}
  >{\raggedright\arraybackslash}p{(\linewidth - 14\tabcolsep) * \real{0.1250}}
  >{\raggedright\arraybackslash}p{(\linewidth - 14\tabcolsep) * \real{0.1250}}@{}}
\toprule\noalign{}
\begin{minipage}[b]{\linewidth}\raggedright
HB-id
\end{minipage} & \begin{minipage}[b]{\linewidth}\raggedright
Title
\end{minipage} & \begin{minipage}[b]{\linewidth}\raggedright
Tier
\end{minipage} & \begin{minipage}[b]{\linewidth}\raggedright
Phase
\end{minipage} & \begin{minipage}[b]{\linewidth}\raggedright
Target subsystem
\end{minipage} & \begin{minipage}[b]{\linewidth}\raggedright
Source ADR
\end{minipage} & \begin{minipage}[b]{\linewidth}\raggedright
Acceptance
\end{minipage} & \begin{minipage}[b]{\linewidth}\raggedright
CAPCOM HARD GATE
\end{minipage} \\
\midrule\noalign{}
\endhead
\bottomrule\noalign{}
\endlastfoot
HB-08 & AGS / RPS Chipset Divergence Metrics & T2-backlog & TBD &
\texttt{src/chipset-validation/rps-ags/} & ADR-\arxivid{2604.21255} & CS25-01 (ADR
coverage) + future runtime test & --- \\
HB-09 & DryRUN Test Synthesis & T2-backlog & TBD &
\texttt{src/cartridge/validation/synth/} & ADR-\arxivid{2604.21598} & CS25-09
(HB-09 deferred-backlog future) & --- \\
HB-10 & IRAP 5-Round Elicitation Cap & T2-backlog & TBD &
\texttt{src/model-affinity/} +
\texttt{.claude/skills/vision-to-mission/} & ADR-\arxivid{2604.21380} & CS25-10
(HB-10 deferred-backlog future) & --- \\
HB-11 & Black-Box Skill Stealing Threat Model & T2-backlog & TBD (doc
may land earlier) &
\texttt{docs/security/dolthub-federated-economy-threat-model.md} &
ADR-\arxivid{2604.21829} & CS25-11 + CS25-04 citation gate & --- \\
HB-12 & FSFM 4-Class Forgetting Taxonomy & T2-backlog & TBD &
\texttt{.planning/policies/} + \texttt{src/reinforcement/} + Memory
Arena demotion + done-retirement & ADR-\arxivid{2604.20300} & CS25-01 + CS25-05
cite & --- \\
\end{longtable}

\subsubsection{Pairings + composition
notes}

\begin{itemize}
\tightlist
\item
  \textbf{HB-04 ↔ HB-07} are a semantic pair per xmod CC-3 /
  convergent-discovery §5: HB-04 establishes the
  Worker/Evaluator/Evolution role naming + per-episode adversarial
  check; HB-07 supplies the cross-episode reflection bandit. Bandit
  lives \emph{inside} Evolution role's protocol-update path. Two flags
  toggle independently; the architectural unit is the pair. Both are
  CAPCOM HARD GATE.
\item
  \textbf{HB-05 ↔ HB-06} are a SKILL.md authoring discipline upgrade per
  xmod CC-2: structural-completeness + four-type ambiguity ship together
  as Phase 810 + 811. Both linters land in
  \texttt{src/cartridge/linter/}.
\item
  \textbf{HB-01 ↔ MCP triple-gate} (CC-1): HB-01 Phase 806 ticket is
  widened to include the six-attack-family + function-hijacking test
  fixtures from \arxivid{2604.21477} + \arxivid{2604.20994} per impact.md
  §4 \#10.
\item
  \textbf{HB-09 ↔ HB-06} (when HB-09 promotes): DryRUN test-side +
  Orchid spec-side define the spec-quality / test-quality envelope for
  skill authoring.
\item
  \textbf{HB-11 ↔ HB-05} (when HB-11 promotes): structural-completeness
  linter is the first defense against marketplace contamination; HB-11
  threat-model is the citation gate before federated promotion runs.
\end{itemize}

\subsubsection{Special-handling dispositions (impact.md
§4)}

\begin{itemize}
\tightlist
\item
  \textbf{HB-01 baseline (impact.md §4 \#1):}
  \texttt{token-budget-baseline.test.ts} is a pre-implementation step.
  Phase 806 wave-1 produces a baseline on a synthetic mock-MCP scenario
  with N tools at three complexity tiers before locking the 40\% p50
  reduction threshold; if unreachable, the threshold re-scopes to ``p50
  ≥ baseline × 0.6''.
\item
  \textbf{HB-03 fail-closed bootstrap (impact.md §4 \#2):}
  \texttt{fail-closed-bootstrap.test.ts} covers the bootstrap exception.
  On first run with no calibration: conservative STD-floor across all
  enabled models, one-time WARN, explicit re-calibration trigger
  required before fail-closed activates.
\item
  \textbf{HB-04 ↔ HB-07 composition (convergent-discovery §5):} HB-04
  spec names the integration boundary; bandit operates inside Evolution
  role; both specs cross-reference each other.
\end{itemize}

\subsubsection{Phase 805 dependencies}

The following HB-NN specs are referenced by Phase 805 publication
artefacts. For each, the artefact name is given; the spec is the
engineering-ticket-grade backing reference Phase 805 cites.

\begin{itemize}
\tightlist
\item
  \textbf{HB-01:} verification matrix entry CS25-13; substrate-doc
  reference \texttt{docs/substrate-references/tool-attention.md} (NEW);
  Phase 805 widens the HB-01 ticket to the MCP triple-gate per impact.md
  §4 \#10.
\item
  \textbf{HB-02:} verification matrix entry CS25-14; CLAUDE.md draft
  (CS25-09) --- schema cited as the BLOCK record format the v1.49.576+
  correlator consumes.
\item
  \textbf{HB-03:} verification matrix entry CS25-15;
  \texttt{.planning/safety/std-calibration.json} schema doc; CLAUDE.md
  draft cites Root Theorem C3 fidelity-threshold gating as the formal
  foundation; CAPCOM HARD GATE pass requires this spec's
  bootstrap-exception path before the gate clears.
\item
  \textbf{HB-04:} verification matrix entry CS25-16;
  \texttt{gsd-mission-management-roles.md}; CLAUDE.md ``Architectural
  Foundations'' preamble cites \arxivid{2604.21003} Last Harness as one
  of four convergent-discovery anchors (CS25-09);
  \texttt{.planning/PROJECT.md} bounded-learning-policy section cites
  \arxivid{2604.21003} as published meta-evolution derivation.
\item
  \textbf{HB-05:} verification matrix entry CS25-17; SKILL.md authoring
  guide draft addition (CS25-09) --- five-principle
  structural-completeness check; CLAUDE.md draft cites GROUNDINGmd
  \arxivid{2604.21744} as the convergent-discovery framing that the five
  principles mirror CSO discipline.
\item
  \textbf{HB-06:} verification matrix entry CS25-18; SKILL.md draft adds
  the four-type ambiguity checklist (lexical / syntactic / semantic /
  vagueness) (CS25-09);
  \texttt{docs/skill-authoring/ambiguity-checklist.md} is a Phase 805
  publication artefact.
\item
  \textbf{HB-07:} verification matrix entry CS25-19;
  \texttt{.claude/agents/} documentation update mapping
  executor/verifier/planner onto Worker/Evaluator/Evolution.
\item
  \textbf{HB-08..HB-12:} backlog citations only at v1.49.575; Phase 805
  substrate-doc references each spec as the engineering-ticket-grade
  reference for v1.49.576+ promotion; the threat-model document HB-11
  may land at v1.49.575 as a forward-looking artefact alongside the
  Phase 805 publication.
\end{itemize}

End of Phase 803 Wave 2.4 drop-in integration specs index. CS25-06
satisfied (12 spec files + index).


\subsection{HB-01 — Tool Attention Lazy Schema Loader}
\subsubsection{§Engineering ticket}

Ship a runtime mitigation for the planning-bridge MCP ``Tools Tax''
(10K--60K tokens/turn observed in the Tool Attention paper) at
\texttt{src/orchestration/tool-attention/}. The approach is three pieces
composed: (1) an Intent Schema Overlap (ISO) score computed from
sentence embeddings (Haiku-tier sidecar) that ranks loaded tools by
relevance to the current turn's intent; (2) a state-aware gate that
adjusts the top-k budget by current task phase (planning vs executing vs
verifying changes the admissible set); (3) a two-phase lazy loader that
keeps a compact pool of tool names + 1-line descriptions in context and
fetches full JSON schemas only for the top-k gated tools. A budget
monitor watches the running fraction of context occupied by tool schemas
and warns at the paper's 70\% context-fracture threshold. All four
pieces sit \emph{above} the dispatch boundary; no existing dispatch code
path is rewritten. The flag \texttt{cs25-26-sweep.tool-attention}
defaults off and the load path is byte-identical to v1.49.574 baseline
when off.

\subsubsection{§Target file paths}

\begin{itemize}
\tightlist
\item
  \texttt{src/orchestration/tool-attention/iso-score.ts} --- ISO
  embedding-overlap selector; returns a ranked tool list per turn.
\item
  \texttt{src/orchestration/tool-attention/state-gate.ts} ---
  phase-conditional top-k budget function
  (planning/executing/verifying).
\item
  \texttt{src/orchestration/tool-attention/lazy-loader.ts} --- two-phase
  loader: compact summary pool → full schemas top-k.
\item
  \texttt{src/orchestration/tool-attention/budget-monitor.ts} ---
  token-occupancy tracker with 70\% fracture-threshold alert.
\item
  \texttt{src/orchestration/tool-attention/index.ts} --- public entry
  composing the four modules behind the flag.
\item
  \texttt{src/orchestration/tool-attention/\_\_tests\_\_/iso-score.test.ts},
  \texttt{state-gate.test.ts}, \texttt{lazy-loader.test.ts},
  \texttt{budget-monitor.test.ts}, \texttt{mock-mcp-scenario.test.ts},
  \texttt{token-budget-baseline.test.ts},
  \texttt{flag-off-byte-identical.test.ts},
  \texttt{phase-pin-override.test.ts}, \texttt{fracture-alert.test.ts},
  \texttt{composition.test.ts}.
\end{itemize}

\subsubsection{§Validation tests}

\begin{enumerate}
\def\labelenumi{\arabic{enumi}.}
\tightlist
\item
  \texttt{iso-score.test.ts} --- ISO score returns identical ranks for
  identical embeddings + intent (deterministic).
\item
  \texttt{state-gate.test.ts} --- top-k changes per declared phase and
  respects phase-pinned tools.
\item
  \texttt{lazy-loader.test.ts} --- compact pool contains names + 1-line
  descriptions only; full schemas appear only for the top-k.
\item
  \texttt{budget-monitor.test.ts} --- alert fires at 70\% context
  occupancy on tool schemas; no alert at 69\%.
\item
  \texttt{mock-mcp-scenario.test.ts} --- end-to-end against a mock
  planning-bridge MCP with N=20 tools at varied complexity.
\item
  \textbf{\texttt{token-budget-baseline.test.ts}} ---
  \textbf{baseline-measurement test (CS25-13).} Records the
  all-tools-loaded p50/p95 tokens-per-turn on the synthetic mock-MCP
  scenario described below, then enables the flag and asserts ≥40\% p50
  reduction. \textbf{Pre-implementation step (per impact.md §4 \#1):}
  Phase 806 wave-1 produces this baseline first, on a synthetic scenario
  containing N tools at three complexity tiers (simple-CRUD,
  mid-complexity-with-enums, large-schema-with-nested-objects) drawn
  from real MCP server schemas; if 40\% is unreachable, the threshold is
  re-scoped to ``p50 ≥ baseline × 0.6'' in the Phase 806 plan.
\item
  \texttt{flag-off-byte-identical.test.ts} --- fixture asserting
  byte-identical dispatch path with flag off.
\item
  \texttt{phase-pin-override.test.ts} --- phase-pinned tools survive ISO
  low-rank.
\item
  \texttt{fracture-alert.test.ts} --- alert telemetry surfaces the
  offending occupancy ratio.
\item
  \texttt{composition.test.ts} --- full pipeline composes ISO → gate →
  lazy-load → monitor.
\end{enumerate}

\subsubsection{§Default-off invariant}

With \texttt{cs25-26-sweep.tool-attention=false} the planning-bridge MCP
dispatch path is byte-identical to v1.49.574 baseline. The
\texttt{flag-off-byte-identical.test.ts} fixture mirrors the v1.49.574
\texttt{disabled-byte-identical} pattern: capture the dispatch trace
under flag-off, compare against a recorded v1.49.574 reference trace,
fail on any byte diff including telemetry payloads.

\subsubsection{§Rollback criterion}

\begin{itemize}
\tightlist
\item
  Baseline-measurement deviation: if Phase 806 wave-1 records cannot
  reproduce the 40\% p50 number, the implementation is held until the
  threshold is re-scoped per impact.md §4 \#1.
\item
  ISO-score quality regression: any rare-but-needed tool gated out
  incorrectly in the existing in-tree skill corpus.
\item
  Flag-off byte-diff: any non-zero diff in the dispatch trace under
  flag-off compared to v1.49.574 baseline.
\end{itemize}

\subsubsection{§Cross-references}

\begin{itemize}
\tightlist
\item
  Source ADR:
  \texttt{.planning/missions/cs25-26-sweep/work/modules/M4/ADR-\arxivid{2604.21816}.md}.
\item
  xmod CC-1 MCP triple-gate: composes with \arxivid{2604.21477} (MCP
  Pitfall Lab) + \arxivid{2604.20994} (Function Hijacking) --- Phase 806
  ticket widens to include the six-attack-family + function-hijacking
  test fixtures per impact.md §4 \#10.
\item
  Prior milestones: v1.49.570 convergent-substrate (two-gate
  cascade-mcp-defense pre-citation), v1.49.574 megakernel substrate
  (default-off flag pattern).
\item
  v1.49.572 Root Theorem (\arxivid{2604.20874}) supplies the bounded-tape
  framing under which the 70\% threshold is a fidelity-threshold
  constant, not a capacity constant.
\end{itemize}

\subsubsection{§Open questions for implementing
engineer}

\begin{itemize}
\tightlist
\item
  ISO-score threshold tuning: fixed cosine-similarity cutoff vs adaptive
  (z-score against per-session distribution)?
\item
  Embedding model selection: Haiku-tier sidecar --- which embedding
  family (sentence-transformers vs Anthropic embeddings)?
\item
  Cache invalidation policy on schema fetches: per-turn or per-phase?
\item
  Phase-pin declaration surface: where do skills declare ``I always need
  tool X regardless of ISO rank''?
\end{itemize}


\subsection{HB-02 — AgentDoG Where/How/What BLOCK Schema}
\subsubsection{§Engineering ticket}

Ship a structured diagnostic record alongside every Safety Warden BLOCK
decision at \texttt{src/safety/agentdog/}. The approach lifts AgentDoG's
three-axis taxonomy verbatim: every BLOCK emit gains a \texttt{Where}
axis (offending component + invocation context), a \texttt{How} axis
(vulnerability vector + escalation pattern), and a \texttt{What} axis
(impacted asset + blast radius). The schema is implemented as three
TypeScript modules (one per axis) plus a composing emitter that the
existing Safety Warden hooks into when the flag is on. The 4B-tier
sidecar deployment plan is named --- a Haiku-tier inference pass that
fills in the three axes given the BLOCK context --- but the v1.49.575
ship is the schema + emitter + integration; sidecar deployment is
documented for v1.49.576+ runtime instrumentation. With
\texttt{cs25-26-sweep.agentdog-schema=false}, BLOCK emission is
byte-identical to v1.49.574: no schema field added, no extra log
entries, no telemetry envelope changes. The schema records persist to
Memory Arena M13 content-addressed Grove records as the substrate the
v1.49.576+ cross-session correlator (CSTM-Bench, \arxivid{2604.21131})
consumes.

\subsubsection{§Target file paths}

\begin{itemize}
\tightlist
\item
  \texttt{src/safety/agentdog/where.ts} --- \texttt{Where} axis:
  component + invocation context capture.
\item
  \texttt{src/safety/agentdog/how.ts} --- \texttt{How} axis:
  vulnerability vector + escalation pattern enumeration.
\item
  \texttt{src/safety/agentdog/what.ts} --- \texttt{What} axis: impacted
  asset + blast-radius classification.
\item
  \texttt{src/safety/agentdog/emitter.ts} --- composes the three axes
  into a single AgentDoG record per BLOCK.
\item
  \texttt{src/safety/agentdog/integration.ts} --- hooks into existing
  Safety Warden BLOCK paths when the flag is on.
\item
  \texttt{src/safety/agentdog/index.ts} --- public entry.
\item
  \texttt{src/safety/agentdog/\_\_tests\_\_/where.test.ts},
  \texttt{how.test.ts}, \texttt{what.test.ts}, \texttt{emitter.test.ts},
  \texttt{integration.test.ts}, \texttt{block-path-coverage.test.ts},
  \texttt{flag-off-byte-identical.test.ts},
  \texttt{grove-persistence.test.ts}, \texttt{schema-shape.test.ts},
  \texttt{sidecar-handoff.test.ts}.
\end{itemize}

\subsubsection{§Validation tests}

\begin{enumerate}
\def\labelenumi{\arabic{enumi}.}
\tightlist
\item
  \texttt{where.test.ts} --- \texttt{Where} axis records component
  identifier + invocation context (call-site stack frame, current phase,
  current skill).
\item
  \texttt{how.test.ts} --- \texttt{How} axis enumerates vulnerability
  vector (e.g., prompt-injection / metadata-poisoning /
  function-hijacking) + escalation pattern.
\item
  \texttt{what.test.ts} --- \texttt{What} axis classifies impacted asset
  (secrets / signing keys / \texttt{.git/} / wasteland exclusion / etc.)
  and blast-radius tier.
\item
  \texttt{emitter.test.ts} --- emitter composes a complete record;
  missing-field cases throw not silently fill.
\item
  \texttt{integration.test.ts} --- integration with each existing Safety
  Warden BLOCK path emits an AgentDoG record.
\item
  \texttt{block-path-coverage.test.ts} --- every existing BLOCK path in
  the repo is exercised; every path emits a record.
\item
  \texttt{flag-off-byte-identical.test.ts} --- fixture asserting
  byte-identical BLOCK emission with flag off; no schema field added at
  runtime.
\item
  \texttt{grove-persistence.test.ts} --- records persist to Memory Arena
  M13 content-addressed Grove records under the right namespace.
\item
  \texttt{schema-shape.test.ts} --- schema validates against a
  JSON-schema definition; CI fails on shape drift.
\item
  \texttt{sidecar-handoff.test.ts} --- record is shaped to be consumable
  by the v1.49.576+ Haiku-tier sidecar (forward compatibility).
\end{enumerate}

\subsubsection{§Default-off invariant}

With \texttt{cs25-26-sweep.agentdog-schema=false}, BLOCK emission is
byte-identical to v1.49.574 baseline --- no schema field added at
runtime, no extra log line, no telemetry payload change. The
\texttt{flag-off-byte-identical.test.ts} fixture mirrors the v1.49.574
\texttt{disabled-byte-identical} pattern by recording the BLOCK emission
trace under flag-off and comparing against a v1.49.574 reference trace.

\subsubsection{§Rollback criterion}

\begin{itemize}
\tightlist
\item
  Flag-off byte-diff: any non-zero diff in BLOCK emission with the flag
  off vs v1.49.574 baseline.
\item
  Schema-shape drift: any change in the JSON schema that would break the
  v1.49.576+ sidecar contract.
\item
  BLOCK-path coverage gap: any existing BLOCK path that does not emit a
  record when the flag is on.
\end{itemize}

\subsubsection{§Cross-references}

\begin{itemize}
\tightlist
\item
  Source ADR:
  \texttt{.planning/missions/cs25-26-sweep/work/modules/M3/ADR-\arxivid{2601.18491}.md}.
\item
  xmod CC-6 cross-session threat correlator: composes with
  \arxivid{2604.20911} (STD-floor breach events) + \arxivid{2604.21131}
  (CSTM-Bench correlator consumer).
\item
  Prior milestones: v1.49.561 trust-relationship (Stage 1 --- supplies
  the relationship axis the \texttt{Where} field references), Memory
  Arena M13 (content-addressed Grove records as persistence layer).
\item
  v1.49.572 GROUNDINGmd Hard Constraints (\arxivid{2604.21744}) supply
  the override-user-prompt invariants the \texttt{What} axis records.
\end{itemize}

\subsubsection{§Open questions for implementing
engineer}

\begin{itemize}
\tightlist
\item
  Sidecar latency budget: 4B model on Haiku tier --- what is the
  acceptable per-BLOCK latency before sidecar inference is gated behind
  a ``structured-only at flush'' mode?
\item
  Schema versioning: schema changes likely with v1.49.576+ correlator
  wiring --- embed a version field now or rely on Grove
  content-addressing for migration?
\item
  \texttt{Where} axis scope: should it record the full call-stack or
  only the top frame + a cryptographic hash of the rest?
\item
  Privacy boundary: BLOCK on a user's secret should not leak the
  secret's content into the AgentDoG record --- what is the redaction
  discipline?
\end{itemize}


\subsection{HB-03 — Safe-Turn-Depth Calibration}
\subsubsection{§Engineering ticket}

Ship a per-model omission-decay calibration + re-injection middleware at
\texttt{src/safety/std-calibration/}. The approach is three pieces: (1)
a calibration harness that measures per-model
prohibition-type-constraint compliance across conversation depths (Opus
/ Sonnet / Haiku), reproducing the paper's 73\%@turn-5 → 33\%@turn-16
observation in GSD's own deployment; (2) a calibration table at
\texttt{.planning/safety/std-calibration.json} recording the per-model
Safe Turn Depth (STD) --- the conversation depth before omission
compliance falls below tolerance; (3) a re-injection middleware that
fires when conversation depth approaches a model's STD floor,
re-injecting the active omission constraints from the GROUNDINGmd /
CLAUDE.md Hard Constraints set. The flag
\texttt{cs25-26-sweep.std-calibration} defaults off; flag-off
conversation flow is byte-identical to v1.49.574. \textbf{CAPCOM HARD
GATE} --- touches Safety Warden BLOCK timing; requires explicit human
authorization at the gate boundary, no autonomous pass. Records emitted
on STD-floor approach are AgentDoG-shaped (HB-02 schema). The
cross-session correlator (\arxivid{2604.21131}, v1.49.576+) consumes
STD-floor approach events as cross-session signal.

\subsubsection{§Target file paths}

\begin{itemize}
\tightlist
\item
  \texttt{src/safety/std-calibration/measure.ts} --- calibration
  harness: runs synthetic prohibition-decay scenarios per model.
\item
  \texttt{src/safety/std-calibration/table.ts} --- read/write the JSON
  calibration table at \texttt{.planning/safety/std-calibration.json}.
\item
  \texttt{src/safety/std-calibration/middleware.ts} --- re-injection
  middleware fired at depth ≈ STD.
\item
  \texttt{src/safety/std-calibration/bootstrap.ts} --- fail-closed
  bootstrap exception (see below).
\item
  \texttt{src/safety/std-calibration/index.ts} --- public entry;
  integrates with Safety Warden middleware chain.
\item
  \texttt{src/safety/std-calibration/\_\_tests\_\_/measure-opus.test.ts},
  \texttt{measure-sonnet.test.ts}, \texttt{measure-haiku.test.ts},
  \texttt{table-roundtrip.test.ts},
  \texttt{re-injection-regression.test.ts},
  \texttt{fail-closed-bootstrap.test.ts},
  \texttt{flag-off-byte-identical.test.ts},
  \texttt{capcom-gate.test.ts}.
\end{itemize}

\subsubsection{§Validation tests}

\begin{enumerate}
\def\labelenumi{\arabic{enumi}.}
\tightlist
\item
  \texttt{measure-opus.test.ts} / \texttt{measure-sonnet.test.ts} /
  \texttt{measure-haiku.test.ts} --- per-model STD measurement over a
  synthetic prohibition-decay benchmark; outputs a numerical STD that
  matches the paper's range within tolerance.
\item
  \texttt{table-roundtrip.test.ts} --- write + read of the calibration
  table preserves all fields and per-model records.
\item
  \texttt{re-injection-regression.test.ts} --- at depth = STD-1, no
  re-injection; at depth = STD, re-injection fires and emits an
  AgentDoG-shaped record.
\item
  \textbf{\texttt{fail-closed-bootstrap.test.ts}} ---
  \textbf{bootstrap-exception test (impact.md §4 \#2).} On first run
  with no calibration data: the system enters a bootstrap mode that
  defaults to a \emph{conservative STD-floor across all enabled models}
  (≤turn-5, the lowest reported STD in the paper), emits a one-time
  WARN-level diagnostic, and requires an explicit re-calibration trigger
  (CLI:
  \texttt{npx\ skill-creator\ safety\ std-calibrate\ \textless{}model\textgreater{}})
  before fail-closed activates fully. The test sets up an empty
  calibration table, asserts the conservative floor activates, asserts
  the warning emits exactly once per process start, and asserts that
  without the explicit trigger the system does not silently switch to
  autonomous fail-closed.
\item
  \texttt{flag-off-byte-identical.test.ts} --- fixture asserting
  conversation flow is byte-identical to v1.49.574 baseline with flag
  off; no re-injection fires, no calibration table is read.
\item
  \textbf{\texttt{capcom-gate.test.ts} --- CAPCOM gate test.} Without a
  recorded human authorization (no CAPCOM-pass marker in
  \texttt{.planning/safety/std-calibration.capcom}), the middleware
  refuses to engage even with the flag on; the gate fires and the
  engagement is blocked. Mirrors the v1.49.574 megakernel-substrate
  HB-04 adapter-selection-schema CAPCOM-gate-test pattern.
\item
  Integration test verifying the middleware composes with HB-02 AgentDoG
  emitter so STD events emit AgentDoG-shaped records.
\item
  Per-model boundary test: models not in the enabled set produce a clean
  error rather than silent passthrough.
\end{enumerate}

\subsubsection{§Default-off invariant}

With \texttt{cs25-26-sweep.std-calibration=false}, no constraint
re-injection fires, conversation flow is byte-identical to v1.49.574,
calibration table is not read, no STD events emit. The
\texttt{flag-off-byte-identical.test.ts} fixture mirrors v1.49.574
\texttt{disabled-byte-identical}.

\subsubsection{§Rollback criterion}

\begin{itemize}
\tightlist
\item
  CAPCOM gate bypass: any code path that engages re-injection without a
  recorded human authorization marker.
\item
  Bootstrap-exception failure: bootstrap mode silently switches to
  fail-closed without the explicit re-calibration trigger.
\item
  Re-injection placement defect: re-injection that creates a false
  confidence spike at the exact turn where decay is most acute.
\item
  Flag-off byte-diff: any non-zero diff in conversation flow under
  flag-off vs v1.49.574 baseline.
\end{itemize}

\subsubsection{§Cross-references}

\begin{itemize}
\tightlist
\item
  Source ADR:
  \texttt{.planning/missions/cs25-26-sweep/work/modules/M3/ADR-\arxivid{2604.20911}.md}.
\item
  xmod CC-6 cross-session threat correlator: composes with HB-02
  AgentDoG (\arxivid{2601.18491}) + CSTM-Bench (\arxivid{2604.21131}).
\item
  Prior milestone v1.49.569 Drift in LLM Systems --- STD calibration is
  the prohibition-decay specialization of v1.49.569's drift detection
  methodology; bootstrap-exception path references v1.49.569 drift
  telemetry per impact.md §3.
\item
  \textbf{CAPCOM gate-pattern template:} v1.49.574 megakernel-substrate
  \texttt{src/cartridge/megakernel/adapter-selection-schema.ts} --- the
  prior CAPCOM HARD GATE this spec inherits the gate structure from.
  Re-use the \texttt{disabled-byte-identical.test.ts} fixture pattern
  verbatim.
\item
  v1.49.572 Root Theorem (\arxivid{2604.20874}) C3 fidelity-threshold
  gating supplies the formal foundation: STD is the
  omission-decay-specific instance of the fidelity-threshold gate
  category.
\end{itemize}

\subsubsection{§Open questions for implementing
engineer}

\begin{itemize}
\tightlist
\item
  Calibration cadence: re-measure on every model release vs scheduled
  re-calibration (weekly / monthly)?
\item
  Bootstrap conservative-floor value: literal turn-5 vs a fraction of
  the most-conservative model's expected STD?
\item
  Re-injection content: full Hard Constraints set vs only the
  constraints flagged as recently violated?
\item
  Storage of the \texttt{.capcom} authorization marker: file with user
  signature vs git-tracked attestation?
\end{itemize}


\subsection{HB-04 — Worker / Evaluator / Evolution Role Split}
\subsubsection{§Engineering ticket}

Ship the Worker / Evaluator / Evolution naming convention +
role-isolated subagents at \texttt{src/skill-creator/roles/}. The
approach is three role-isolated subagent contracts: (1) \textbf{Worker}
generates candidate skills against the target task; (2)
\textbf{Evaluator} \emph{adversarially} diagnoses Worker output against
failure-history (replacing the current passive 3-correction-min rule
with active failure-finding) and holds BLOCK authority over Worker
outputs; (3) \textbf{Evolution} proposes loop-protocol modifications
based on cross-skill patterns. The split is \emph{additive}: it names
what skill-creator's six-step loop already does (Observe → Detect →
Suggest → Manage → Auto-Load → Learn/Compose) rather than restructuring
it. The flag \texttt{cs25-26-sweep.weler-roles} defaults off; flag-off
skill-creator runs the existing six-step loop byte-identically.
\textbf{CAPCOM HARD GATE} --- touches skill-creator orchestration;
requires explicit human authorization. \textbf{HB-04 composes with HB-07
(AEL bandit) per xmod CC-3 / convergent-discovery §5:} HB-04 establishes
the role naming and per-episode adversarial check; HB-07 supplies the
cross-episode reflection bandit. Together they form the fast-Worker /
fast-Evaluator / slow-Evolution architecture. The integration boundary:
HB-07's bandit lives \emph{inside} HB-04's role scaffolding --- the
bandit is one of Evolution's protocol-update mechanisms, not a sibling
system. Phase 808 Plan must enumerate the role-isolation invariants per
impact.md §4 \#11: Worker writes own state only; Evaluator reads Worker
output (read-only); Evolution reads aggregated history (read-only); no
cross-role write paths.

\subsubsection{§Target file paths}

\begin{itemize}
\tightlist
\item
  \texttt{src/skill-creator/roles/worker.ts} --- Worker contract:
  candidate-skill generator.
\item
  \texttt{src/skill-creator/roles/evaluator.ts} --- Evaluator contract:
  adversarial diagnosis + BLOCK authority over Worker output.
\item
  \texttt{src/skill-creator/roles/evolution.ts} --- Evolution contract:
  loop-protocol-modification proposer.
\item
  \texttt{src/skill-creator/roles/isolation.ts} --- runtime
  role-isolation enforcer (state-write boundary checks).
\item
  \texttt{src/skill-creator/roles/index.ts} --- public entry composing
  the three roles into the existing six-step loop.
\item
  \texttt{gsd-mission-management-roles.md} --- formal naming-convention
  doc.
\item
  \texttt{.claude/agents/} documentation update --- map existing
  executor/verifier/planner onto Worker/Evaluator/Evolution.
\item
  \texttt{src/skill-creator/roles/\_\_tests\_\_/worker.test.ts},
  \texttt{evaluator.test.ts}, \texttt{evolution.test.ts},
  \texttt{isolation-no-cross-write.test.ts},
  \texttt{evaluator-blocks-worker.test.ts},
  \texttt{six-step-integration.test.ts},
  \texttt{flag-off-byte-identical.test.ts},
  \texttt{capcom-gate.test.ts}, \texttt{hb07-composition.test.ts},
  \texttt{adversarial-vs-3-correction.test.ts}.
\end{itemize}

\subsubsection{§Validation tests}

\begin{enumerate}
\def\labelenumi{\arabic{enumi}.}
\tightlist
\item
  \texttt{worker.test.ts} --- Worker generates candidates against the
  target task; output shape conforms to contract.
\item
  \texttt{evaluator.test.ts} --- Evaluator adversarially diagnoses
  Worker output against failure-history; finds known-planted failure
  modes.
\item
  \texttt{evolution.test.ts} --- Evolution proposes loop-protocol
  modifications based on aggregated history.
\item
  \texttt{isolation-no-cross-write.test.ts} --- role-isolation
  invariant: any cross-role write path is rejected at runtime.
\item
  \texttt{evaluator-blocks-worker.test.ts} --- Evaluator's BLOCK
  authority over Worker outputs fires correctly; Worker cannot bypass.
\item
  \texttt{six-step-integration.test.ts} --- integration with the
  existing six-step loop is correct; outputs at each step match baseline
  shape.
\item
  \texttt{flag-off-byte-identical.test.ts} --- fixture asserting
  skill-creator runs the existing six-step loop byte-identically with
  flag off.
\item
  \textbf{\texttt{capcom-gate.test.ts} --- CAPCOM gate test.} Without a
  recorded human authorization marker, the role split refuses to engage
  even with the flag on. Mirrors the v1.49.574 megakernel-substrate
  HB-04 adapter-selection-schema gate pattern.
\item
  \textbf{\texttt{hb07-composition.test.ts}} --- verifies the
  integration boundary with HB-07: with both HB-04 and HB-07 flags on,
  the AEL bandit operates inside the Evolution role's protocol-update
  path; Evolution reads bandit history at the configured slow-timescale
  interval.
\item
  \texttt{adversarial-vs-3-correction.test.ts} --- comparative test:
  Evaluator catches a planted failure that the legacy 3-correction-min
  rule misses.
\end{enumerate}

\subsubsection{§Default-off invariant}

With \texttt{cs25-26-sweep.weler-roles=false}, skill-creator runs the
existing six-step loop byte-identically; the executor / verifier /
planner subagents continue using their current invocation paths; no
role-isolation enforcer fires. The
\texttt{flag-off-byte-identical.test.ts} fixture mirrors v1.49.574
\texttt{disabled-byte-identical}.

\subsubsection{§Rollback criterion}

\begin{itemize}
\tightlist
\item
  CAPCOM gate bypass: any code path that engages role split without a
  recorded human authorization marker.
\item
  Role-isolation invariant violation: any cross-role write path that
  survives the isolation enforcer.
\item
  Evaluator authority degradation: Evaluator's BLOCK authority becomes
  advisory-only --- at that point the role split degenerates to
  passive-counting and the upgrade is illusory (per ADR-\arxivid{2604.21003}
  Notes).
\item
  Flag-off byte-diff: any non-zero diff in six-step loop output under
  flag-off vs v1.49.574 baseline.
\end{itemize}

\subsubsection{§Cross-references}

\begin{itemize}
\tightlist
\item
  Source ADR:
  \texttt{.planning/missions/cs25-26-sweep/work/modules/M1/ADR-\arxivid{2604.21003}.md};
  \textbf{composes with}
  \texttt{.planning/missions/cs25-26-sweep/work/modules/M2/ADR-\arxivid{2604.21725}.md}
  (HB-07 / AEL).
\item
  xmod CC-3 AEL × Last Harness composed two-loop pattern
  (convergent-discovery §5 fast/slow timescale split).
\item
  Prior milestones: v1.49.571 LeJEPA / heuristics-free skill-space (Last
  Harness Worker/Evaluator/Evolution split = heuristics-free pattern;
  the adversarial Evaluator replaces the heuristic 3-correction-min);
  v1.49.574 megakernel substrate (gate pattern + flag pattern).
\item
  \textbf{CAPCOM gate-pattern template:} v1.49.574 megakernel-substrate
  \texttt{src/cartridge/megakernel/adapter-selection-schema.ts} --- the
  prior CAPCOM HARD GATE this spec inherits the gate structure from.
\item
  \texttt{.planning/PROJECT.md} bounded-learning-policy section ---
  gains a citation to \arxivid{2604.21003} as the published
  meta-evolution derivation (under CS25-09's edit budget).
\end{itemize}

\subsubsection{§Open questions for implementing
engineer}

\begin{itemize}
\tightlist
\item
  Role isolation enforced via runtime check (per-write callback) vs
  static type discipline (TypeScript role-tagged types) vs both?
\item
  Evaluator BLOCK authority surface: does Evaluator emit a BLOCK record
  into the AgentDoG schema (HB-02), or is BLOCK an internal control-flow
  signal only?
\item
  Evolution slow-timescale cadence: per-skill / per-mission /
  per-milestone interval --- and where is the cadence configured?
\item
  HB-07 integration boundary detail: does Evolution own the bandit's
  posterior storage, or is the bandit's storage shared with the
  auto-load step?
\end{itemize}


\subsection{HB-05 — Five-Principle Structural-Completeness Linter}
\subsubsection{§Engineering ticket}

Ship a SKILL.md validation linter implementing the five-principle
structural-completeness check at
\texttt{src/cartridge/linter/structural-completeness.ts}. The five
principles, lifted verbatim from \arxivid{2604.21090}, are: (1)
\textbf{computability-theory grounding} for any computational claims a
SKILL.md makes; (2) \textbf{proof-theoretic structure} for any deductive
claims; (3) \textbf{Bayesian-epistemology framing} for any uncertainty
claims; (4) \textbf{data-classification specification} for any
data-handling claims; (5) \textbf{assessment-rubric specification} for
any quality claims. The paper's empirical 37\% incompleteness rate makes
the linter operationally necessary; the principles are mechanically
checkable. Linter runs as part of the cartridge promotion pipeline and
gates promotion out of \texttt{.planning/staging/inbox/}. The flag
\texttt{cs25-26-sweep.structural-completeness-lint} defaults off →
linter runs in \textbf{warning-only} mode (does not block promotion);
when on, missing-principle records become blocking. \textbf{HB-05 ships
paired with HB-06} (Four-Type Ambiguity Checklist, \arxivid{2604.21505})
per xmod CC-2 --- the two linters land as a single SKILL.md authoring
discipline upgrade. The linter ingests the GROUNDINGmd two-class
taxonomy from \arxivid{2604.21744} as a structural-completeness
criterion: a SKILL.md must distinguish Hard Constraints
(override-the-prompt invariants the skill refuses to violate) from
Convention Parameters (defaults that allow override). This becomes the
``Hard Constraints declared'' check in the negative-fixture corpus.

\subsubsection{§Target file paths}

\begin{itemize}
\tightlist
\item
  \texttt{src/cartridge/linter/structural-completeness.ts} ---
  five-principle validator.
\item
  \texttt{src/cartridge/linter/principles/computability.ts} ---
  Principle 1 check.
\item
  \texttt{src/cartridge/linter/principles/proof-theory.ts} --- Principle
  2 check.
\item
  \texttt{src/cartridge/linter/principles/bayesian.ts} --- Principle 3
  check.
\item
  \texttt{src/cartridge/linter/principles/data-classification.ts} ---
  Principle 4 check.
\item
  \texttt{src/cartridge/linter/principles/assessment-rubric.ts} ---
  Principle 5 check.
\item
  \texttt{src/cartridge/linter/principles/hard-constraints-declared.ts}
  --- GROUNDINGmd two-class check (cross-link with \arxivid{2604.21744}).
\item
  \texttt{src/cartridge/linter/\_\_tests\_\_/structural-completeness.test.ts},
  plus per-principle test files and positive/negative fixtures under
  \texttt{src/cartridge/linter/\_\_tests\_\_/fixtures/}.
\end{itemize}

\subsubsection{§Validation tests}

\begin{enumerate}
\def\labelenumi{\arabic{enumi}.}
\tightlist
\item
  \texttt{principle-1-computability.test.ts} --- fixture: SKILL.md with
  computational claim and no computability grounding flagged; with
  grounding passes.
\item
  \texttt{principle-2-proof-theory.test.ts} --- fixture: deductive claim
  without proof structure flagged.
\item
  \texttt{principle-3-bayesian.test.ts} --- fixture: uncertainty claim
  without Bayesian framing flagged.
\item
  \texttt{principle-4-data-classification.test.ts} --- fixture:
  data-handling claim without classification flagged.
\item
  \texttt{principle-5-assessment-rubric.test.ts} --- fixture: quality
  claim without rubric flagged.
\item
  \textbf{Negative-fixture corpus} --- at least 3 SKILL.md negative
  fixtures, each missing exactly one principle; each fails the
  corresponding check and only that check.
\item
  \textbf{Positive-fixture corpus} --- at least 3 SKILL.md positive
  fixtures (drawn from existing in-tree skills selected for compliance)
  pass all five principles.
\item
  \texttt{hard-constraints-declared.test.ts} --- fixture: SKILL.md
  without a Hard-Constraints-vs-Convention-Parameters split flagged;
  cross-link test against \arxivid{2604.21744}.
\item
  \texttt{flag-off-warning-only.test.ts} --- fixture asserts that with
  the flag off, the linter emits warnings but does not block promotion
  (existing behavior preserved).
\item
  \texttt{flag-off-byte-identical.test.ts} --- promotion pipeline
  behavior is byte-identical to v1.49.574 baseline with flag off
  (warnings emit but do not affect exit code or pipeline state).
\end{enumerate}

\subsubsection{§Default-off invariant}

With \texttt{cs25-26-sweep.structural-completeness-lint=false}, the
linter runs in warning-only mode: failures emit a warning to stderr but
do not change the cartridge promotion exit code, do not block promotion,
and do not modify the pipeline state. The
\texttt{flag-off-byte-identical.test.ts} fixture mirrors v1.49.574
\texttt{disabled-byte-identical} for the promotion pipeline state.

\subsubsection{§Rollback criterion}

\begin{itemize}
\tightlist
\item
  False positives on existing in-tree skills: zero false positives is
  the CS25-18 acceptance condition. If any existing in-tree skill trips
  a check, treat as a fixture defect (refine the check), not a baseline
  regression.
\item
  Warning-only-mode breakage: the linter blocks promotion with the flag
  off --- this would regress the cartridge promotion baseline.
\item
  Mechanical-checkability defect: a principle check that requires LLM
  judgment rather than mechanical analysis (the paper's own static
  analyzer is the reference implementation pattern).
\end{itemize}

\subsubsection{§Cross-references}

\begin{itemize}
\tightlist
\item
  Source ADR:
  \texttt{.planning/missions/cs25-26-sweep/work/modules/M3/ADR-\arxivid{2604.21090}.md}.
\item
  xmod CC-2 SKILL.md authoring discipline upgrade: ships paired with
  HB-06 (\arxivid{2604.21505}).
\item
  Prior milestones: v1.49.573 Skilldex Conformance Auditor (compounds
  with this linter --- same gate, complementary checks); v1.49.574
  megakernel substrate (default-off flag pattern).
\item
  \arxivid{2604.21744} GROUNDINGmd convergent-discovery framing: the five
  principles \textbf{mirror} GSD's existing CSO discipline rather than
  introducing a new norm --- Phase 810 ticket includes the
  cross-reference per impact.md §4 \#15. (This is the
  convergent-discovery ↔ structural-completeness rhyme the published
  Phase 805 substrate doc names.)
\item
  v1.49.576+ trigger: gate flips from warning-only to blocking after
  corpus audit completes (≤2\% false-positive rate on the v1.49.575
  corpus).
\end{itemize}

\subsubsection{§Open questions for implementing
engineer}

\begin{itemize}
\tightlist
\item
  Mechanical-checkability boundary: which principles can be fully
  mechanical (computability / proof-theory / data-classification ---
  likely yes) and which need a heuristic-with-confidence pattern
  (Bayesian framing / assessment rubric)?
\item
  Negative-fixture corpus authorship: hand-author ≥3 fixtures or scrape
  from public skill marketplaces (and risk the marketplace contamination
  concern HB-11 addresses)?
\item
  Hard-Constraints-declared check signature: a YAML frontmatter block
  (\texttt{hard\_constraints:\ {[}...{]}}) vs a section heading
  (\texttt{\#\#\ Hard\ Constraints})?
\item
  Cross-link with HB-06 ambiguity linter: do the two linters share a
  common rule-engine, or do they run as independent passes?
\end{itemize}


\subsection{HB-06 — Four-Type Ambiguity Checklist}
\subsubsection{§Engineering ticket}

Ship a SKILL.md ambiguity linter implementing Orchid's four-type
taxonomy at \texttt{src/cartridge/linter/ambiguity.ts}, plus an
authoring-guidance document at
\texttt{docs/skill-authoring/ambiguity-checklist.md}. The four types,
lifted from \arxivid{2604.21505}, are: (1) \textbf{lexical} ---
terminology with multiple plausible meanings; (2) \textbf{syntactic} ---
grammar admitting multiple parses; (3) \textbf{semantic} ---
under-specified intent admitting multiple plausible interpretations; (4)
\textbf{vagueness} --- quantifier-free claims, ``appropriate'' /
``reasonable'' without operationalization. Each frontmatter capability
declaration receives four checks. A failing check blocks
\texttt{skill-creator\ cartridge\ publish} until resolved. The flag
\texttt{cs25-26-sweep.ambiguity-lint} defaults off → warnings only. The
Orchid headline finding --- \emph{advanced models suffer more from
ambiguity} --- directly justifies the chipset architecture's policy that
Opus-tier work demands stricter specs, not laxer; this is captured in
the authoring-guidance doc as the rationale section. \textbf{HB-06 ships
paired with HB-05} per xmod CC-2 (SKILL.md authoring discipline
upgrade). The two linters land together as Phase 810 + 811, and the
README index lists the pairing explicitly. The linter is a heuristic +
rule-based pass: lexical = terminology dictionary check; syntactic =
grammar-parse-ambiguity flag; semantic = quantifier presence +
reference-resolution check; vagueness = pattern-match against a corpus
of known vague terms.

\subsubsection{§Target file paths}

\begin{itemize}
\tightlist
\item
  \texttt{src/cartridge/linter/ambiguity.ts} --- four-type ambiguity
  linter entry point.
\item
  \texttt{src/cartridge/linter/ambiguity/lexical.ts} --- Type 1 check:
  terminology dictionary lookup.
\item
  \texttt{src/cartridge/linter/ambiguity/syntactic.ts} --- Type 2 check:
  grammar-parse-ambiguity flag.
\item
  \texttt{src/cartridge/linter/ambiguity/semantic.ts} --- Type 3 check:
  under-specified-intent detection.
\item
  \texttt{src/cartridge/linter/ambiguity/vagueness.ts} --- Type 4 check:
  quantifier + vague-term pattern match.
\item
  \texttt{docs/skill-authoring/ambiguity-checklist.md} ---
  authoring-guidance doc with the rationale (``advanced models suffer
  more'').
\item
  \texttt{tools/lint-skill.mjs} --- wire ambiguity checks into the
  \texttt{skill-creator\ cartridge\ publish} path.
\item
  \texttt{src/cartridge/linter/\_\_tests\_\_/ambiguity-\{lexical,syntactic,semantic,vagueness\}.test.ts},
  \texttt{ambiguity-publish-block.test.ts},
  \texttt{flag-off-byte-identical.test.ts},
  \texttt{existing-skills-zero-false-positives.test.ts}.
\end{itemize}

\subsubsection{§Validation tests}

\begin{enumerate}
\def\labelenumi{\arabic{enumi}.}
\tightlist
\item
  \texttt{ambiguity-lexical.test.ts} --- fixture: SKILL.md with
  multi-meaning term flagged; resolved-term version passes.
\item
  \texttt{ambiguity-syntactic.test.ts} --- fixture: SKILL.md with
  grammar-ambiguous claim flagged; rewritten version passes.
\item
  \texttt{ambiguity-semantic.test.ts} --- fixture: under-specified
  intent flagged; specified-intent version passes.
\item
  \texttt{ambiguity-vagueness.test.ts} --- fixture: quantifier-free
  claim (``appropriate'', ``reasonable'') flagged; operationalized
  version passes.
\item
  \texttt{ambiguity-publish-block.test.ts} ---
  \texttt{skill-creator\ cartridge\ publish} fails with one of four
  distinct error codes (one per ambiguity type) given a minimal SKILL.md
  exercising each type.
\item
  \textbf{\texttt{existing-skills-zero-false-positives.test.ts}} ---
  runs the linter against every in-tree SKILL.md in \texttt{examples/}
  and \texttt{.claude/skills/}; \textbf{CS25-18 acceptance: zero false
  positives at baseline.} Any in-tree skill that trips an ambiguity
  check is treated as a fixture defect (refine the check), not a
  baseline regression.
\item
  \texttt{flag-off-byte-identical.test.ts} --- promotion pipeline
  behavior is byte-identical to v1.49.574 baseline with flag off
  (warnings emit but no block).
\end{enumerate}

\subsubsection{§Default-off invariant}

With \texttt{cs25-26-sweep.ambiguity-lint=false}, the linter runs in
warning-only mode: failures emit warnings to stderr but do not change
exit codes or block \texttt{skill-creator\ cartridge\ publish}. The
\texttt{flag-off-byte-identical.test.ts} fixture mirrors v1.49.574
\texttt{disabled-byte-identical} for the promotion pipeline state.

\subsubsection{§Rollback criterion}

\begin{itemize}
\tightlist
\item
  Any false positive on an existing in-tree skill at baseline (CS25-18
  acceptance condition): treat as fixture defect, refine the check, do
  not regress.
\item
  Author-friction explosion: if the linter creates ≥1 hour additional
  friction per skill publish, the heuristic thresholds need
  re-calibration.
\item
  Flag-off block: any code path that blocks publish with the flag off
  would regress the cartridge baseline.
\end{itemize}

\subsubsection{§Cross-references}

\begin{itemize}
\tightlist
\item
  Source ADR:
  \texttt{.planning/missions/cs25-26-sweep/work/modules/M5/ADR-\arxivid{2604.21505}.md}.
\item
  \textbf{Pairs with HB-05} per xmod CC-2 (SKILL.md authoring discipline
  upgrade): both linters ship as Phase 810 + 811, README index lists the
  pairing.
\item
  Prior milestones: v1.49.572 output-structure frontmatter at
  \texttt{src/output-structure/} (the frontmatter is the surface the
  linter inspects); v1.49.570 reinforcement series MA-6 emitters at
  \texttt{src/reinforcement/} (each linter pass emits a reinforcement
  event so the eligibility layer can credit ambiguity-reduction skills
  retroactively).
\item
  xmod connection \arxivid{2604.21090} ↔ \arxivid{2604.21505}
  composes-with strong (table row): both run on the same
  \texttt{src/cartridge/linter/} directory.
\end{itemize}

\subsubsection{§Open questions for implementing
engineer}

\begin{itemize}
\tightlist
\item
  Lexical-dictionary source: hand-curated GSD-domain dictionary vs
  WordNet vs LLM-judged?
\item
  Syntactic-parse engine: standard NLP parser (spaCy /
  tree-sitter-markdown) vs a GSD-specific frontmatter grammar?
\item
  Semantic-check threshold: how much under-specification is tolerable
  before flagging? (e.g., is ``the planner runs over the artifacts''
  specific enough?)
\item
  Vague-term corpus maintenance: where does the list of vague terms
  live, and how is it versioned across releases?
\end{itemize}


\subsection{HB-07 — AEL Fast/Slow Bandit Auto-Load}
\subsubsection{§Engineering ticket}

Ship a Thompson Sampling bandit + LLM-reflection slow loop at
\texttt{src/skill-creator/auto-load/bandit.ts}. The approach is two
pieces: (1) \textbf{fast layer} --- a Thompson Sampling bandit selecting
which retrieval policy to apply per episode (instead of skill-creator's
current single fixed retrieval policy in the auto-load step); (2)
\textbf{slow layer} --- an LLM-driven reflection callback after N
failures that diagnoses recurring failure patterns and injects causal
insights back into the policy priors. Bandit posterior estimates feed
the existing stochastic-selection layer (\texttt{src/stochastic/}, MA-3
+ MD-2). Slow-timescale reflection becomes a reinforcement signal
feeding eligibility-trace updates over \texttt{src/learnable-k\_h/}
(MD-5 per-(skill, task-type) heads). The flag
\texttt{cs25-26-sweep.ael-bandit} defaults off; flag-off auto-load
behavior is byte-identical to v1.49.574. \textbf{CAPCOM HARD GATE} ---
touches skill-space dispatch; requires explicit human authorization.
\textbf{HB-07 lives inside HB-04's role-split scaffolding} per xmod CC-3
/ convergent-discovery §5: the bandit operates inside the
\textbf{Evolution} role's protocol-update path; Evolution reads the
bandit's posterior history at the configured slow-timescale interval and
proposes role-isolation invariants for the next epoch. The two flags are
independently toggleable but the semantic pair (HB-04 + HB-07) is the
architectural unit. Per impact.md §4 \#11, the Phase 812 plan must
enumerate state-isolation invariants: bandit writes own state only;
reflection step reads bandit history (read-only).

\subsubsection{§Target file paths}

\begin{itemize}
\tightlist
\item
  \texttt{src/skill-creator/auto-load/bandit.ts} --- Thompson Sampling
  bandit fast layer.
\item
  \texttt{src/skill-creator/auto-load/policies.ts} --- declared
  retrieval policies (the bandit's arms).
\item
  \texttt{src/skill-creator/auto-load/reflection.ts} --- LLM-driven
  slow-timescale reflection callback.
\item
  \texttt{src/skill-creator/auto-load/priors.ts} --- bandit-prior
  storage + slow-loop insight injection.
\item
  \texttt{src/skill-creator/auto-load/index.ts} --- public entry;
  integrates with skill-creator's auto-load step + HB-04 Evolution role.
\item
  \texttt{src/skill-creator/auto-load/\_\_tests\_\_/bandit-convergence.test.ts},
  \texttt{policy-switching.test.ts},
  \texttt{reflection-trigger.test.ts},
  \texttt{reflection-feedback-loop.test.ts},
  \texttt{prior-update.test.ts}, \texttt{state-isolation.test.ts},
  \texttt{flag-off-byte-identical.test.ts},
  \texttt{capcom-gate.test.ts}, \texttt{hb04-integration.test.ts},
  \texttt{eligibility-trace-integration.test.ts}.
\end{itemize}

\subsubsection{§Validation tests}

\begin{enumerate}
\def\labelenumi{\arabic{enumi}.}
\tightlist
\item
  \texttt{bandit-convergence.test.ts} --- bandit converges on a
  synthetic 3-policy benchmark within bounded episodes; convergence
  verified against the paper's Sharpe-ratio-style stationary-signal
  regime.
\item
  \texttt{policy-switching.test.ts} --- bandit switches arms when one
  policy's posterior degrades; switching latency bounded.
\item
  \texttt{reflection-trigger.test.ts} --- slow-loop reflection fires
  after N failures; N is configurable.
\item
  \texttt{reflection-feedback-loop.test.ts} --- reflection-injected
  insights update policy priors; the next episode's bandit-arm selection
  reflects the updated prior.
\item
  \texttt{prior-update.test.ts} --- prior-update path is read-only on
  bandit history (state-isolation invariant: reflection step does not
  mutate bandit posterior storage directly).
\item
  \texttt{state-isolation.test.ts} --- verifies the no-cross-write
  boundary between bandit / reflection / external skill-creator state.
\item
  \texttt{flag-off-byte-identical.test.ts} --- fixture asserting
  auto-load behavior is byte-identical to v1.49.574 baseline with flag
  off; existing single-policy retrieval runs unchanged.
\item
  \textbf{\texttt{capcom-gate.test.ts} --- CAPCOM gate test.} Without a
  recorded human authorization marker, the bandit refuses to engage even
  with the flag on; the gate fires and dispatch is blocked. Mirrors the
  v1.49.574 megakernel-substrate HB-04 adapter-selection-schema gate
  pattern.
\item
  \textbf{\texttt{hb04-integration.test.ts}} --- verifies the
  integration boundary with HB-04: bandit operates inside HB-04
  Evolution role; Evolution reads bandit history at configured
  slow-timescale interval; bandit posterior storage is owned by
  Evolution (or shared via a documented contract).
\item
  \texttt{eligibility-trace-integration.test.ts} --- slow-loop insights
  feed MA-1 eligibility traces over MD-5 learnable-k\_h heads;
  reinforcement series sees the new signal.
\end{enumerate}

\subsubsection{§Default-off invariant}

With \texttt{cs25-26-sweep.ael-bandit=false}, auto-load behavior is
byte-identical to v1.49.574: existing single-fixed-policy retrieval runs
unchanged, no bandit posterior is initialised, no reflection callback
fires. The \texttt{flag-off-byte-identical.test.ts} fixture mirrors
v1.49.574 \texttt{disabled-byte-identical}.

\subsubsection{§Rollback criterion}

\begin{itemize}
\tightlist
\item
  CAPCOM gate bypass: any code path that engages the bandit without a
  recorded human authorization marker.
\item
  State-isolation violation: any code path that allows reflection-step
  writes to bandit posterior or vice versa.
\item
  Bandit non-convergence in production: bandit underperforms a
  fixed-policy baseline because skill priors drift faster than the
  bandit converges (the paper's Risk note); flag remains default-off
  until convergence is empirically confirmed.
\item
  Flag-off byte-diff: any non-zero diff in auto-load output under
  flag-off vs v1.49.574 baseline.
\end{itemize}

\subsubsection{§Cross-references}

\begin{itemize}
\tightlist
\item
  Source ADR:
  \texttt{.planning/missions/cs25-26-sweep/work/modules/M2/ADR-\arxivid{2604.21725}.md}.
\item
  \textbf{Composes with HB-04} (\arxivid{2604.21003},
  \texttt{.planning/missions/cs25-26-sweep/work/modules/M1/ADR-\arxivid{2604.21003}.md})
  per xmod CC-3 / convergent-discovery §5.
\item
  xmod connection \arxivid{2604.21003} ↔ \arxivid{2604.21725}
  composes-with strong (table row).
\item
  Prior milestones: v1.49.570 reinforcement series MA-6 (canonical
  taxonomy + emitters); v1.49.572 MA-1 eligibility traces; v1.49.574
  megakernel substrate (default-off flag pattern + CAPCOM gate
  template).
\item
  \textbf{CAPCOM gate-pattern template:} v1.49.574 megakernel-substrate
  \texttt{src/cartridge/megakernel/adapter-selection-schema.ts} --- the
  prior CAPCOM HARD GATE this spec inherits the gate structure from.
  Re-use the \texttt{disabled-byte-identical.test.ts} fixture pattern.
\item
  Sensoria/umwelt/symbiosis (M6/M7/M8 from v1.49.561) --- symbiosis
  teaching ledger receives reflection-injected insights as durable
  across-mission artefacts.
\end{itemize}

\subsubsection{§Open questions for implementing
engineer}

\begin{itemize}
\tightlist
\item
  Thompson Sampling prior shape: Beta / Gaussian / per-arm hyperprior
  --- which is the right hyperparameter, and how is it documented to
  avoid silent regime drift between releases (per ADR Notes)?
\item
  Slow-loop cadence: failure-count threshold N --- fixed value (e.g.,
  N=5 from IRAP) vs adaptive?
\item
  Reflection-output format: structured JSON-injectable insights vs
  free-text-with-extraction?
\item
  Posterior storage: ephemeral (per-mission) vs persisted (Grove M13
  records)?
\item
  HB-04 Evolution boundary: does Evolution own the bandit's posterior,
  or does HB-07 ship its own storage with a read-contract for Evolution?
\end{itemize}


\subsection{HB-08 — RPS / AGS Cross-Chipset Divergence Audit}
\subsubsection{§Engineering ticket}

Ship two chipset-validation metrics at
\texttt{src/chipset-validation/rps-ags/}. The approach implements the
two metrics from \arxivid{2604.21255}: (1) \textbf{Response Pattern
Similarity (RPS)} --- verbal-alignment metric across paired agent
responses, computed over a tokenization that strips formatting noise;
(2) \textbf{Action Graph Similarity (AGS)} --- tool-use behavior metric
modeling agent invocations as directed graphs and comparing graph
structure across agents. The metrics' use is \textbf{strictly internal
monitoring} of the GSD chipset (Agnus/Denise/Paula/Gary →
Opus/Sonnet/Haiku), never as a public claim that any external model is
``identical to'' another (cultural-sensitivity gate per intake). Use AGS
to validate Opus-design / Sonnet-implement boundary in the chipset is
producing genuinely complementary perspectives, not collapsing to
behavior-pattern uniformity. Use RPS to detect prompt-leakage at the
Haiku tier. \textbf{HB-08 is deferred to v1.49.576+ backlog} per
impact.md §4 \#3 --- the chipset-validation harness is not yet scoped,
and ad-hoc cross-tier behavior monitoring has not yet emerged as an
operational pain point. Promotion trigger: (a) chipset architecture's
tier-separation produces a measurable semantic-drift incident, or (b)
external skill marketplaces (DoltHub, see HB-11) require behavioral
provenance verification of submitted skills.

\subsubsection{§Target file paths}

\begin{itemize}
\tightlist
\item
  \texttt{src/chipset-validation/rps-ags/rps.ts} --- Response Pattern
  Similarity metric.
\item
  \texttt{src/chipset-validation/rps-ags/ags.ts} --- Action Graph
  Similarity metric.
\item
  \texttt{src/chipset-validation/rps-ags/graph-builder.ts} ---
  directed-graph construction from chipset traces.
\item
  \texttt{src/chipset-validation/rps-ags/internal-only-guard.ts} ---
  runtime guard blocking external-claim invocations
  (cultural-sensitivity gate).
\item
  \texttt{src/chipset-validation/rps-ags/index.ts} --- public entry.
\item
  \texttt{docs/architecture/chipset.md} --- citation block for
  \arxivid{2604.21255}; metric definitions documented even before runtime
  ships.
\item
  \texttt{src/chipset-validation/rps-ags/\_\_tests\_\_/rps-paired-traces.test.ts},
  \texttt{ags-graph-similarity.test.ts},
  \texttt{internal-only-guard.test.ts},
  \texttt{chipset-tier-divergence.test.ts},
  \texttt{flag-off-byte-identical.test.ts},
  \texttt{metric-stability.test.ts},
  \texttt{kimi-k2-style-detection.test.ts},
  \texttt{cultural-sensitivity-gate.test.ts}.
\end{itemize}

\subsubsection{§Validation tests}

\begin{enumerate}
\def\labelenumi{\arabic{enumi}.}
\tightlist
\item
  \texttt{rps-paired-traces.test.ts} --- paired chipset traces yield
  reproducible RPS; cross-family pairs score lower than within-family
  per the paper's 5.9 percentage-point finding.
\item
  \texttt{ags-graph-similarity.test.ts} --- paired action-graphs yield
  reproducible AGS; planted-divergence fixture detected.
\item
  \texttt{internal-only-guard.test.ts} --- any caller passing an
  external-model trace as the ``test'' axis is rejected at runtime;
  cultural-sensitivity gate fires.
\item
  \texttt{chipset-tier-divergence.test.ts} --- Opus-design /
  Sonnet-implement boundary detection on a known-divergent fixture.
\item
  \texttt{flag-off-byte-identical.test.ts} --- chipset dispatch is
  byte-identical to v1.49.575 baseline with flag off (when promoted).
\item
  \texttt{metric-stability.test.ts} --- small input perturbations yield
  small metric changes (no metric instability).
\item
  \texttt{kimi-k2-style-detection.test.ts} --- replicates the paper's
  Kimi-K2 ↔ Claude Sonnet 4.5 result on a synthetic fixture (not against
  any real external model --- cultural-sensitivity gate).
\item
  \texttt{cultural-sensitivity-gate.test.ts} --- verifies the
  internal-only guard cannot be disabled via configuration.
\end{enumerate}

\subsubsection{§Default-off invariant}

With \texttt{cs25-26-sweep.rps-ags=false} (when promoted), chipset
dispatch is byte-identical to baseline; no metric computation runs in
the dispatch path. The \texttt{flag-off-byte-identical.test.ts} fixture
mirrors v1.49.574 \texttt{disabled-byte-identical}.

\subsubsection{§Rollback criterion}

\begin{itemize}
\tightlist
\item
  Cultural-sensitivity gate bypass: any code path that allows the
  metrics to be invoked against external models for public claims.
\item
  Metric instability: small input perturbations cause large metric
  swings, making divergence detection unreliable.
\item
  Flag-off byte-diff: any non-zero diff in chipset dispatch under
  flag-off vs baseline.
\end{itemize}

\subsubsection{§Cross-references}

\begin{itemize}
\tightlist
\item
  Source ADR:
  \texttt{.planning/missions/cs25-26-sweep/work/modules/M4/ADR-\arxivid{2604.21255}.md}.
\item
  xmod CC-6 Two-game chipset validation: composes with
  \arxivid{2604.21282} 3+1 + \arxivid{2603.22823} Empirical Comparison
  ACPs (DACP fidelity-level evaluation matrix).
\item
  Prior milestones: v1.49.572 DACP fidelity levels (semantic-channel +
  chipset-tier separation surface); v1.49.574 chipset architecture (the
  dispatch boundary the metrics measure across); Memory Arena M9-M13
  (action-graph persistence).
\item
  HB-11 (\arxivid{2604.21829} Black-Box Skill Stealing) --- promotion
  trigger (b): federated-skill marketplace provenance verification.
\end{itemize}

\subsubsection{§Open questions for implementing
engineer}

\begin{itemize}
\tightlist
\item
  Promotion-trigger detection: how is ``measurable semantic drift
  incident'' operationalized (alert threshold + on-call response)?
\item
  Trace persistence: chipset traces persisted to Grove M13 records ---
  what is the retention policy?
\item
  Cultural-sensitivity gate enforcement: runtime guard in code vs
  review-discipline at the documentation level only?
\item
  Calibration of divergence threshold: what AGS / RPS numbers indicate
  ``tier separation working'' vs ``uniformity collapse''?
\end{itemize}


\subsection{HB-09 — DryRUN Autonomous Test Synthesis}
\subsubsection{§Engineering ticket}

Ship an autonomous test-input + trace-simulation synthesizer at
\texttt{src/cartridge/validation/synth/}. The approach implements
DryRUN's pattern: given a SKILL.md frontmatter signature + the cartridge
contract, the synthesizer (1) generates candidate inputs covering the
declared capability surface; (2) simulates expected outputs against the
cartridge's declared invariants; (3) returns a trace bundle that the
cartridge validator uses as a self-correction signal \emph{before} any
human-authored ground-truth fixtures are required. The synthesizer plugs
into \texttt{skill-creator\ cartridge\ validate\ -\/-synth-trace}.
\textbf{HB-09 is deferred to v1.49.576+ backlog} per ADR Notes ---
DryRUN's value depends on an external check (the v1.50 proof-companion's
Lean verifier) closing the loop; until that check exists, a synthesized
trace can confirm a wrong invariant. \textbf{HB-09 promotion is gated on
a ≥800-word architecture-spec ADR} in \texttt{work/modules/M5/} per
impact.md §4 \#4 before the implementing phase opens --- the current
sketch is one paragraph; promotion requires the expanded spec. Promotion
trigger: v1.50 Half B chapter validation begins specifying its
inner-loop test contract.

\subsubsection{§Target file paths}

\begin{itemize}
\tightlist
\item
  \texttt{src/cartridge/validation/synth/inputs.ts} --- candidate-input
  generator from SKILL.md frontmatter.
\item
  \texttt{src/cartridge/validation/synth/trace-simulator.ts} ---
  execution-trace simulator.
\item
  \texttt{src/cartridge/validation/synth/self-correct.ts} ---
  self-correction signal feeder.
\item
  \texttt{src/cartridge/validation/synth/bundle.ts} --- trace bundle
  output format.
\item
  \texttt{src/cartridge/validation/synth/index.ts} --- public entry;
  integrates with \texttt{skill-creator\ cartridge\ validate}.
\item
  \texttt{src/cartridge/validation/synth/\_\_tests\_\_/input-coverage.test.ts},
  \texttt{trace-simulator.test.ts},
  \texttt{self-correct-signal.test.ts}, \texttt{bundle-shape.test.ts},
  \texttt{proof-companion-handoff.test.ts},
  \texttt{flag-off-byte-identical.test.ts},
  \texttt{wrong-invariant-detection.test.ts},
  \texttt{reinforcement-event-emit.test.ts}.
\end{itemize}

\subsubsection{§Validation tests}

\begin{enumerate}
\def\labelenumi{\arabic{enumi}.}
\tightlist
\item
  \texttt{input-coverage.test.ts} --- synthesized inputs cover the
  declared capability surface (positive + boundary + adversarial cases).
\item
  \texttt{trace-simulator.test.ts} --- simulated traces are
  deterministic given identical input + frontmatter.
\item
  \texttt{self-correct-signal.test.ts} --- self-correction signal
  correctly drives a re-generation loop on a fixture with a planted spec
  defect.
\item
  \texttt{bundle-shape.test.ts} --- trace bundle conforms to the
  proof-companion handoff schema.
\item
  \texttt{proof-companion-handoff.test.ts} --- bundle is consumable by
  the v1.50 proof-companion's Lean external check (the closure point).
\item
  \texttt{flag-off-byte-identical.test.ts} ---
  \texttt{skill-creator\ cartridge\ validate} is byte-identical to
  baseline with the flag off; no synth runs.
\item
  \textbf{\texttt{wrong-invariant-detection.test.ts}} --- fixture: spec
  with a wrong invariant; without the proof-companion external check,
  the synthesizer self-confirms the wrong invariant. \textbf{This is the
  rollback signal} --- if the proof-companion is not yet wired, this
  test must fail explicitly to prevent silent wrong-invariant
  confirmations.
\item
  \texttt{reinforcement-event-emit.test.ts} --- synthesized trace deltas
  emit MA-6 reinforcement events; eligibility layer credits
  trace-discovery skills retroactively.
\end{enumerate}

\subsubsection{§Default-off invariant}

With \texttt{cs25-26-sweep.dryrun-synth=false} (when promoted),
\texttt{skill-creator\ cartridge\ validate} runs the existing path with
hand-authored expectations only; no synthesizer fires. The
\texttt{flag-off-byte-identical.test.ts} fixture mirrors v1.49.574
\texttt{disabled-byte-identical}.

\subsubsection{§Rollback criterion}

\begin{itemize}
\tightlist
\item
  Proof-companion not yet wired: per ADR, do not wire DryRUN into the
  v1.49.575 cartridge harness; wait until the v1.50 proof-companion's
  external Lean check exists.
\item
  Wrong-invariant self-confirmation:
  \texttt{wrong-invariant-detection.test.ts} would silently pass --- at
  that point the synthesizer is masking spec bugs, not finding them.
\item
  Flag-off byte-diff: any non-zero diff in cartridge validator output
  under flag-off vs baseline.
\end{itemize}

\subsubsection{§Cross-references}

\begin{itemize}
\tightlist
\item
  Source ADR:
  \texttt{.planning/missions/cs25-26-sweep/work/modules/M5/ADR-\arxivid{2604.21598}.md}.
\item
  xmod connection \arxivid{2604.21598} ↔ \arxivid{2604.21505}
  composes-with strong: DryRUN test-side + Orchid spec-side define the
  spec-quality / test-quality envelope for skill authoring; HB-09 +
  HB-06 form the joint pair when HB-09 promotes.
\item
  Prior milestones: v1.49.572 output-structure validator at
  \texttt{src/output-structure/} (frontmatter contract is the
  synthesizer's natural input); v1.49.570 reinforcement series MA-6
  emitters; v1.50 proof-companion (the closure point this spec defers
  to).
\item
  xmod CC-7 M5 verification triad: \arxivid{2511.23159}
  probable-to-provable + \arxivid{2604.18792} DSLTrans +
  \arxivid{2604.21187} Doubly Saturated Ramsey --- the
  rhetorical+mechanism+precedent backbone DryRUN slots into.
\end{itemize}

\subsubsection{§Open questions for implementing
engineer}

\begin{itemize}
\tightlist
\item
  Architecture-spec authorship: who writes the ≥800-word ADR (impact.md
  §4 \#4 promotion gate) --- the implementing engineer or a separate
  Wave-1-style track?
\item
  Synthesizer model selection: small + fast (Haiku) vs large + accurate
  (Sonnet/Opus)?
\item
  Trace-bundle format: Lean-readable directly vs an intermediate JSON
  the proof-companion adapter parses?
\item
  Coverage metric: how is ``covers the declared capability surface''
  measured (line coverage on a deterministic implementation skeleton)?
\end{itemize}


\subsection{HB-10 — IRAP Interactive Requirements Elicitation}
\subsubsection{§Engineering ticket}

Ship a bounded-rounds elicitation policy at
\texttt{docs/capcom/elicitation-policy.md} plus an enforcement hook in
\texttt{src/model-affinity/}. The approach implements IRAP's five-round
elicitation cap: any CAPCOM gate clarifying dialogue that exceeds five
rounds triggers a model-tier escalation event (Haiku → Sonnet → Opus)
rather than continuing the dialogue. The five-round cap is the
operational target --- the paper reports up to 40-fold improvement on
the elicitation benchmark \emph{with} the bounded-round property --- and
frames ``cannot resolve in five rounds'' as a misshaped task or wrong
tier rather than under-promotion. The vision-to-mission skill at
\texttt{.claude/skills/vision-to-mission/} adopts the bound as soft
discipline now (no escalation, just a warning at round-5); the hard
CAPCOM-gate escalation is the v1.49.576+ landing site. \textbf{HB-10 is
deferred to v1.49.576+ backlog} per ADR --- full IRAP-based mechanism
waits for the model-affinity escalation phase to open. Promotion
trigger: model-affinity escalation policy phase opens. The
interactive-RAG mechanism specifically pairs with the College
KB-construction work (CC-8 M4 retrieval + synthesis pipeline) ---
elicitation queries the department KB, and the KB structure shapes which
clarifying questions the loop can ask.

\subsubsection{§Target file paths}

\begin{itemize}
\tightlist
\item
  \texttt{docs/capcom/elicitation-policy.md} --- bounded-rounds
  elicitation-policy doc; can land earlier as soft discipline.
\item
  \texttt{src/model-affinity/elicitation-cap.ts} --- five-round cap
  enforcement hook.
\item
  \texttt{src/model-affinity/escalation.ts} --- Haiku → Sonnet → Opus
  escalation trigger on round-5 timeout.
\item
  \texttt{src/model-affinity/round-counter.ts} --- round-count metric
  plumbed into CAPCOM dialogue trace.
\item
  \texttt{src/model-affinity/\_\_tests\_\_/cap-enforcement.test.ts},
  \texttt{escalation-trigger.test.ts}, \texttt{round-counter.test.ts},
  \texttt{vision-to-mission-soft.test.ts},
  \texttt{flag-off-byte-identical.test.ts},
  \texttt{safety-warden-metric.test.ts}.
\item
  \texttt{.claude/skills/vision-to-mission/SKILL.md} --- adds the soft
  five-round discipline note now.
\end{itemize}

\subsubsection{§Validation tests}

\begin{enumerate}
\def\labelenumi{\arabic{enumi}.}
\tightlist
\item
  \texttt{cap-enforcement.test.ts} --- five-round cap fires; round 6
  cannot proceed without an escalation event.
\item
  \texttt{escalation-trigger.test.ts} --- round-5 timeout triggers a
  model-tier escalation event; tier transition is recorded.
\item
  \texttt{round-counter.test.ts} --- round-count metric plumbed into the
  CAPCOM dialogue trace; visible in trace output.
\item
  \texttt{vision-to-mission-soft.test.ts} --- vision-to-mission emits a
  warning at round 5 when the soft discipline is engaged but no
  escalation fires.
\item
  \texttt{flag-off-byte-identical.test.ts} --- CAPCOM dialogue behavior
  is byte-identical to baseline with the flag off (when promoted); no
  escalation, no round-counter metric.
\item
  \texttt{safety-warden-metric.test.ts} --- five-round bound is reported
  as a Safety Warden metric on every CAPCOM gate dialogue.
\end{enumerate}

\subsubsection{§Default-off invariant}

With \texttt{cs25-26-sweep.irap-elicitation=false} (when promoted),
CAPCOM dialogue runs unbounded as in v1.49.574 baseline; no round
counter, no escalation, no metric. The
\texttt{flag-off-byte-identical.test.ts} fixture mirrors v1.49.574
\texttt{disabled-byte-identical} for the dialogue trace.

\subsubsection{§Rollback criterion}

\begin{itemize}
\tightlist
\item
  Hard cutoff surprise: a hard escalation that surprises users without
  UI surfacing of the round count (per ADR Risk note); mitigation is
  mandatory pre-cap UI surfacing.
\item
  Five-round cap too tight for unfamiliar domains: if measured
  cap-failures exceed an operational threshold on real CAPCOM traces,
  the cap is renegotiated (or the escalation tier is widened).
\item
  Flag-off byte-diff: any non-zero diff in CAPCOM dialogue output under
  flag-off vs baseline.
\end{itemize}

\subsubsection{§Cross-references}

\begin{itemize}
\tightlist
\item
  Source ADR:
  \texttt{.planning/missions/cs25-26-sweep/work/modules/M6/ADR-\arxivid{2604.21380}.md}.
\item
  xmod CC-8 College of Knowledge retrieval + synthesis pipeline:
  composes with \arxivid{2604.20848} MATRAG, \arxivid{2510.26615} SARA,
  \arxivid{2604.21444} HiCrew, \arxivid{2508.21720} PosterForest,
  \arxivid{2601.10003} SocraticKG.
\item
  Prior milestones: v1.49.572 output-structure frontmatter
  (per-capability elicitation expectations); v1.49.574 chipset
  architecture (treats CAPCOM as a chip with a bounded-dialogue
  contract).
\item
  xmod connection \arxivid{2604.20855} ↔ \arxivid{2604.21380}
  composes-with: Caesar adversarial-draft-refinement pairs with IRAP's
  bounded elicitation as elicitation+synthesis halves.
\item
  xmod connection \arxivid{2510.18731} ↔ \arxivid{2604.21380}
  same-pattern-different-domain: RLAAR's abstention-at-round-N mirrors
  IRAP's bounded-rounds escalation.
\end{itemize}

\subsubsection{§Open questions for implementing
engineer}

\begin{itemize}
\tightlist
\item
  Five-round cap configurability: hard-coded 5 vs configurable per task
  type?
\item
  Escalation tier policy: Haiku → Sonnet → Opus is the default --- is
  there a fast path that skips a tier on certain task signals?
\item
  Soft vs hard discipline boundary: when does vision-to-mission's soft
  warning become a hard block, and what is the trigger?
\item
  UI surfacing of round count: where does it appear in the CAPCOM gate
  UI to avoid the ``hard cutoff surprise'' Risk?
\end{itemize}


\subsection{HB-11 — Black-Box Skill-Stealing Threat Model}
\subsubsection{§Engineering ticket}

Ship a threat-model document at
\texttt{docs/security/dolthub-federated-economy-threat-model.md} that
adopts the attack taxonomy from \arxivid{2604.21829} as GSD's
defender-concern enumeration for the v1.49.576+ DoltHub federated
skill-economy work. The approach is documentation-first: enumerate each
named pipeline stage from the paper (model-generated seed prompts →
scenario rationalization with structure injection → embedding-filter
diversity step), name a corresponding detection or mitigation
requirement for each, and ship the document as a \textbf{pre-rollout
citation gate} for any DoltHub work in v1.49.576+. This is a
threat-model spec, not a code module. The market-sizing context (90,368
free-marketplace skills + 2,000+ paid listings with \$100K+ creator
earnings) informs the Wasteland-integration roadmap's scope. \textbf{Per
impact.md §4 \#5 and Fox Companies IP rule}, the document is split: an
abstract attack-taxonomy version is publicly citable; the GSD-specific
defender-attack-surface specifics live in
\texttt{.planning/fox-companies/foxfiber/threat-model.md} and require
Fox Companies IP review before any artefact moves outside
\texttt{.planning/fox-companies/foxfiber/}. Promotion trigger: DoltHub
contribution roadmap opens or first federated-skill phase begins.
Composes with HB-05 structural-completeness linter (CC-2 federated-skill
futures) --- marketplace-sourced skills must pass structural
completeness before federated promotion; structural quality is the first
defense against marketplace contamination.

\subsubsection{§Target file paths}

\begin{itemize}
\tightlist
\item
  \texttt{docs/security/dolthub-federated-economy-threat-model.md} ---
  public abstract attack-taxonomy version.
\item
  \texttt{.planning/fox-companies/foxfiber/threat-model.md} ---
  GSD-specific defender-attack-surface specifics (Fox Companies IP; not
  public).
\item
  \texttt{docs/GROVE-FORMAT.md} --- gains a forward-pointing note that
  any DoltHub federated-economy adoption in v1.49.576+ must cite
  \arxivid{2604.21829} as the threat-taxonomy pre-rollout-gate (per
  CS25-04 acceptance language).
\item
  \texttt{src/cartridge/linter/marketplace-attestation.ts} --- runtime
  hook (when promoted): refuses federated-skill publish without a
  threat-model attestation.
\item
  \texttt{docs/security/\_\_tests\_\_/citation-presence.test.ts},
  \texttt{threat-stage-coverage.test.ts},
  \texttt{attestation-required.test.ts},
  \texttt{fox-ip-publication-boundary.test.ts}.
\end{itemize}

\subsubsection{§Validation tests}

\begin{enumerate}
\def\labelenumi{\arabic{enumi}.}
\tightlist
\item
  \texttt{citation-presence.test.ts} --- every
  federated-economy-adjacent recommendation in \texttt{docs/},
  \texttt{.planning/}, and SKILL.md files cites \arxivid{2604.21829}.
  Mechanical grep + assertion (CS25-04 acceptance).
\item
  \texttt{threat-stage-coverage.test.ts} --- for each named pipeline
  stage in the paper, the threat-model doc names a defender concern + a
  corresponding detection/mitigation requirement.
\item
  \texttt{attestation-required.test.ts} (when promoted) ---
  \texttt{skill-creator\ cartridge\ publish\ -\/-federated} refuses to
  proceed without a threat-model attestation that references this
  document.
\item
  \texttt{fox-ip-publication-boundary.test.ts} --- the public version
  contains only abstract attack-taxonomy content; mechanical check that
  no \texttt{.planning/fox-companies/} content has been mirrored into
  the public doc.
\end{enumerate}

\subsubsection{§Default-off invariant}

This is a documentation artefact. The ``default-off'' semantic
translates to: in v1.49.575, only the threat-model doc lands; no DoltHub
publish path engages, no federated-economy code runs. With future
federated-economy work in v1.49.576+, the runtime attestation hook
defaults off until the threat-model doc has landed and a CS25-04
citation-presence audit has passed. The fixture pattern is the
citation-presence sweep: with the federated-skill flag off, no
federated-skill publish path can execute; flag-off behavior is
byte-identical to v1.49.574.

\subsubsection{§Rollback criterion}

\begin{itemize}
\tightlist
\item
  Fox Companies IP leak: any GSD-specific defender-attack-surface
  specifics appear in the public doc. Per impact.md §4 \#5, surface to
  user and revert immediately.
\item
  Citation-gate bypass: any v1.49.576+ federated-economy recommendation
  that does not cite \arxivid{2604.21829} (CS25-04 violation).
\item
  Pre-publish attestation bypass: any code path that allows
  \texttt{skill-creator\ cartridge\ publish\ -\/-federated} to proceed
  without a threat-model attestation.
\end{itemize}

\subsubsection{§Cross-references}

\begin{itemize}
\tightlist
\item
  Source ADR:
  \texttt{.planning/missions/cs25-26-sweep/work/modules/M6/ADR-\arxivid{2604.21829}.md}.
\item
  xmod CC-9 cultural-sensitivity Safety Warden surface: integrates with
  Grove via the federated-skill-future link.
\item
  xmod connection \arxivid{2604.21829} ↔ \arxivid{2604.21090}
  composes-with strong: Structural Quality Gaps in governance prompts
  (HB-05) provides the defender-side structural-quality gate that
  complements the attack-taxonomy.
\item
  xmod connection \arxivid{2604.21829} ↔ \arxivid{2604.21477}
  composes-with strong: MCP Pitfall Lab's six-attack family is the
  runtime-side complement.
\item
  Prior milestones: v1.49.561 trust-relationship (Stage 1 ---
  trust-relationship-provider would consume the attack taxonomy as one
  of its threat sources); Memory Arena M13
  (\texttt{docs/GROVE-FORMAT.md} --- content addressing is a natural
  defense against several attack-pipeline stages because it makes silent
  skill substitution detectable); v1.49.574 chipset architecture
  (federated skill economy is a chipset-level concern).
\item
  \textbf{Sign-off invariant (impact.md §6 \#3):} \arxivid{2604.21829} is
  a pre-rollout gate for any federated-economy / DoltHub work in
  v1.49.576+.
\end{itemize}

\subsubsection{§Open questions for implementing
engineer}

\begin{itemize}
\tightlist
\item
  Public-vs-private boundary detail: which attack-pipeline stages can be
  discussed at the abstract level publicly, and which must stay private
  until defenses ship?
\item
  Attestation format: free-form Markdown attestation vs structured
  frontmatter the publisher pipeline can validate?
\item
  Marketplace contamination defense composition: HB-11 + HB-05 +
  content-addressed Grove records --- what is the order-of-operations
  the publisher pipeline runs them in?
\item
  Sovereign-marketplace handling: AFRILANGTUTOR-style sovereignty-gated
  marketplaces --- does the threat-model treat those as a separate
  defender posture?
\end{itemize}


\subsection{HB-12 — FSFM Four-Class Forgetting Taxonomy}
\subsubsection{§Engineering ticket}

Ship the FSFM four-class forgetting taxonomy as the canonical
bounded-learning forgetting policy at
\texttt{.planning/policies/bounded-learning-forgetting-taxonomy.md}
(NEW), with citations into existing forgetting-related code paths. The
four classes from \arxivid{2604.20300} map directly onto GSD subsystems:
(1) \textbf{passive decay-based} → Memory Arena M1--M13's hot/warm/cold
tiering with importance-weighted demotion; (2) \textbf{active
deletion-based} → the \texttt{done-retirement} skill's
pipeline-retirement protocol; (3) \textbf{safety-triggered} → the Safety
Warden's BLOCK-on-content authority
(\texttt{.claude/skills/security-hygiene/}); (4) \textbf{adaptive
reinforcement-based} → the reinforcement series
(\texttt{src/reinforcement/} MA-6 + \texttt{src/eligibility/} MA-1
eligibility-trace decay + \texttt{src/temperature/} MD-4 annealed
schedule). The ``shed'' step in the Root Theorem (\arxivid{2604.20874})
accumulate-compress-rewrite-shed cycle is formally defined by this
taxonomy: shedding can be passive, active, safety-triggered, or
reinforcement-adaptive. The paper's reported numbers (+8.49\% access
efficiency, +29.2\% signal-to-noise, 100\% security-content elimination)
become Safety Warden test targets. \textbf{HB-12 is deferred to
v1.49.576+ backlog} per the tier-split directive --- adoption now is
policy-document-only at v1.49.575 (citation in §2 Subsystem 9 of
impact.md), runtime instrumentation defers. Promotion trigger:
bounded-learning theorem proof completion (carried forward from
v1.49.572 Half B T1d) cites Root Theorem as the empirical anchor, and
HB-12 lands as the operational taxonomy alongside.

\subsubsection{§Target file paths}

\begin{itemize}
\tightlist
\item
  \texttt{.planning/policies/bounded-learning-forgetting-taxonomy.md}
  --- adopts the four-class taxonomy verbatim; maps each class to a GSD
  subsystem.
\item
  \texttt{src/reinforcement/forgetting-class.ts} --- class-tagging at
  the reinforcement-event emit point.
\item
  \texttt{src/eligibility/decay-class.ts} --- eligibility-trace decay
  receives a class tag (passive vs adaptive-reinforcement).
\item
  \texttt{src/temperature/decay-rate.ts} --- annealed schedule
  modulating decay rate per class (MD-4 hookup).
\item
  Memory Arena M1--M13 demotion call sites --- tag demotion events with
  \texttt{passive-decay} class.
\item
  \texttt{done-retirement} skill --- tags pipeline-retirement events
  with \texttt{active-deletion} class.
\item
  \texttt{.claude/skills/security-hygiene/SKILL.md} --- tags
  BLOCK-on-content events with \texttt{safety-triggered} class.
\item
  \texttt{.planning/policies/\_\_tests\_\_/taxonomy-coverage.test.ts},
  \texttt{class-tag-emit.test.ts},
  \texttt{safety-triggered-100pct.test.ts},
  \texttt{signal-to-noise-improvement.test.ts},
  \texttt{flag-off-byte-identical.test.ts},
  \texttt{representation-audit-integration.test.ts}.
\end{itemize}

\subsubsection{§Validation tests}

\begin{enumerate}
\def\labelenumi{\arabic{enumi}.}
\tightlist
\item
  \texttt{taxonomy-coverage.test.ts} --- every forgetting code path in
  the repo emits an event tagged with one of the four classes.
\item
  \texttt{class-tag-emit.test.ts} --- each class-tagged event records
  the class verbatim per the FSFM taxonomy.
\item
  \texttt{safety-triggered-100pct.test.ts} --- safety-triggered class
  achieves 100\% security-content elimination on a fixture with planted
  security-tagged content (paper's headline number; Safety Warden test
  target).
\item
  \texttt{signal-to-noise-improvement.test.ts} --- synthetic benchmark
  records ≥some-threshold improvement in signal-to-noise ratio under the
  adaptive-reinforcement class (paper reports +29.2\%; threshold to be
  calibrated).
\item
  \texttt{flag-off-byte-identical.test.ts} --- forgetting code paths run
  unchanged with the flag off; class tags emit as no-ops.
\item
  \texttt{representation-audit-integration.test.ts} --- composes with
  MD-6 representation-audit (collapse detection); audit consumes class
  tags.
\end{enumerate}

\subsubsection{§Default-off invariant}

With \texttt{cs25-26-sweep.fsfm-taxonomy=false} (when promoted),
forgetting code paths run as in v1.49.574 baseline; class tags emit as
no-ops; no behavior change. The \texttt{flag-off-byte-identical.test.ts}
fixture mirrors v1.49.574 \texttt{disabled-byte-identical}. At
v1.49.575, only the policy document lands; no runtime tagging fires.

\subsubsection{§Rollback criterion}

\begin{itemize}
\tightlist
\item
  Behavior change with flag off: any forgetting code path that behaves
  differently with the flag off would regress the bounded-learning
  baseline.
\item
  Taxonomy mismatch: any forgetting code path that does not map cleanly
  into one of the four classes --- expand the taxonomy or document the
  gap.
\item
  Hippocampal-indexing scope creep: per ADR Notes, the GSD adoption is
  taxonomic, not implementation. Any code path that imports the paper's
  specific hippocampal-indexing implementation is out of scope.
\end{itemize}

\subsubsection{§Cross-references}

\begin{itemize}
\tightlist
\item
  Source ADR:
  \texttt{.planning/missions/cs25-26-sweep/work/modules/M2/ADR-\arxivid{2604.20300}.md}.
\item
  xmod connection \arxivid{2604.20300} ↔ \arxivid{2604.20874} implements
  strong: Root Theorem accumulate-compress-rewrite-shed cycle's ``shed''
  step is formally defined by this taxonomy.
\item
  xmod connection \arxivid{2604.20300} ↔ \arxivid{2511.11439} extends
  strong: Retrofit's selective-forgetting + low-rank-sparse subspace
  pairs with the four-class taxonomy.
\item
  xmod CC-4 bounded-tape × bounded-learning unified foundation.
\item
  Prior milestones: v1.49.561 sensoria/umwelt/symbiosis (M6/M7/M8) ---
  symbiosis teaching ledger uses adaptive-reinforcement decay; umwelt
  Markov-blanket boundary triggers safety-triggered forgetting on
  out-of-blanket content; v1.49.572 Bounded-Learning Theorem (T1d) ---
  HB-12 lands alongside the proof completion; v1.49.561
  trust-relationship (Stage 1) uses safety-triggered forgetting when
  trust degrades.
\end{itemize}

\subsubsection{§Open questions for implementing
engineer}

\begin{itemize}
\tightlist
\item
  Class-tag enum vs string: enum (compile-time exhaustiveness) or string
  (forward-compat with future fifth class)?
\item
  Adaptive-reinforcement decay-rate parameterisation: per-skill /
  per-task-type / per-(skill, task-type)? MD-5 already does per-(skill,
  task-type) --- does HB-12 ride on the same indexing?
\item
  Safety-triggered 100\% elimination: how is ``elimination'' verified
  (Grove content-addressed records prove absence; what is the absence
  proof for in-memory state)?
\item
  Hippocampal-indexing pilot: per ADR Notes, ``if a future College of
  Knowledge mind-body experiment wants to pilot hippocampal indexing,
  this ADR is the citation seed'' --- is there a near-term pilot in
  scope?
\end{itemize}


\clearpage

\section{ADR Index (Phases 797--799)}

\textit{This appendix indexes the 55 ADRs produced by Wave 1 (Phase 797 Module 1+2, Phase 798 Module 3+4, Phase 799 Module 5+6). Each row gives the arXiv ID, module, paper title, adoption decision (ADOPT/TRACK/REJECT), target subsystem (truncated where long), target milestone, and source ADR file path. Full ADR content is preserved at \texttt{.planning/missions/cs25-26-sweep/work/modules/M\{1--6\}/ADR-*.md}.}

\begin{footnotesize}
\renewcommand{\arraystretch}{1.15}
\begin{longtable}{@{}L{0.10\textwidth} L{0.025\textwidth} L{0.22\textwidth} L{0.06\textwidth} L{0.20\textwidth} L{0.13\textwidth} L{0.13\textwidth}@{}}
\toprule
\textbf{\sffamily arXiv ID} & \textbf{\sffamily M} & \textbf{\sffamily Title} & \textbf{\sffamily Dec.} & \textbf{\sffamily Target subsystem} & \textbf{\sffamily Target milestone} & \textbf{\sffamily Source ADR} \\
\midrule
\endfirsthead
\toprule
\textbf{\sffamily arXiv ID} & \textbf{\sffamily M} & \textbf{\sffamily Title} & \textbf{\sffamily Dec.} & \textbf{\sffamily Target subsystem} & \textbf{\sffamily Target milestone} & \textbf{\sffamily Source ADR} \\
\midrule
\endhead
\bottomrule
\endfoot
\arxivid{2604.20874} & M1 & The Root Theorem of Context Engineering & ADOPT & \texttt{docs/substrate/bounded-tape-root-theorem.md} (NEW), *The Space Between* information-theory chapter, \texttt{CLAUDE.md} design-rationale section, every CAPCOM-gate policy document, fresh-context-subagent doc\textbackslash\{\}ldots & v1.49.575 (Phase 805 substrate document + CS25-09 CLAUDE.md citation); v1.50 retrospective Half A (information-theory chapter) & \texttt{\small M1/ADR-2604.20874.md} \\
\arxivid{2604.21003} & M1 & The Last Harness You'll Ever Build & ADOPT & \texttt{src/skill-creator/roles/} (NEW), \texttt{gsd-mission-management-roles.md}, \texttt{.claude/agents/} documentation, \texttt{.planning/PROJECT.md} bounded-learning-policy section & v1.49.575 (citation in Phase 801 + role-split spec); Half B Phase 808 ships HB-04 / CS25-16 implementation & \texttt{\small M1/ADR-2604.21003.md} \\
\arxivid{2604.21744} & M1 & Agentic AI-assisted coding offers a unique opportunity to instill epistemic grounding during software development (GROUNDINGmd) & ADOPT & \texttt{CLAUDE.md}, every \texttt{SKILL.md} template (cartridge authoring guide), \texttt{.claude/skills/security-hygiene/SKILL.md}, \texttt{src/cartridge/linter/structural-completeness.ts} (Module 3 spec) & v1.49.575 (citation + linter spec); v1.50 retrospective Half A (extended writeup) & \texttt{\small M1/ADR-2604.21744.md} \\
\arxivid{2604.21910} & M1 & From Research Question to Scientific Workflow: Leveraging Agentic AI for Science Automation & ADOPT & \texttt{CLAUDE.md}, \texttt{.college/README.md}, \texttt{docs/substrate/three-layer-architecture.md} (NEW), \texttt{docs/cartridge/SKILL.md} authoring guide & v1.49.575 (citation in Phase 801 report + CS25-09 CLAUDE.md edit); v1.50 retrospective Half A (extended writeup) & \texttt{\small M1/ADR-2604.21910.md} \\
\arxivid{2509.18629} & M2 & HyperAdapt: Simple High-Rank Adaptation & ADOPT & \texttt{docs/silicon-layer/benchmark-suite.md}, Silicon Layer benchmark spec, future Silicon Layer LoRA pipeline implementation, adapter-pooling experiment design & v1.49.575 (benchmark spec); future Silicon Layer benchmark suite & \texttt{\small M2/ADR-2509.18629.md} \\
\arxivid{2511.11439} & M2 & Retrofit: Continual Learning with Controlled Forgetting for Binary Security Detection & ADOPT & \texttt{docs/silicon-layer/bounded-update-invariants.md} (NEW), bounded-learning policy doc, future Silicon Layer LoRA pipeline & v1.50 proof companion; Silicon Layer benchmarks deferred to post-v1.50 & \texttt{\small M2/ADR-2511.11439.md} \\
\arxivid{2604.20300} & M2 & FSFM: A Biologically-Inspired Framework for Selective Forgetting of Agent Memory & ADOPT & \texttt{.planning/policies/bounded-learning-forgetting-taxonomy.md} (NEW), Safety Warden BLOCK paths, Memory Arena demotion logic, \texttt{done-retirement} skill, reinforcement-series decay parameters & v1.49.575 (taxonomy document); v1.50 proof companion (mathematical formalisation of shed) & \texttt{\small M2/ADR-2604.20300.md} \\
\arxivid{2604.20943} & M2 & SCM: Sleep-Consolidated Memory with Algorithmic Forgetting for LLMs & TRACK & \texttt{.college/experiments/scm-wave-consolidation/} (NEW pilot); future v1.51 wave-architecture refactor if pilot succeeds & v1.49.575 (TRACK status logged); pilot post-v1.50; ADOPT decision in v1.51 if pilot succeeds & \texttt{\small M2/ADR-2604.20943.md} \\
\arxivid{2604.21330} & M2 & Teacher-Guided Routing for Sparse Vision Mixture-of-Experts & TRACK & \texttt{docs/chipset-architecture/future-work-notes.md}, future learned-dispatch-layer design (post-v1.50) & v1.49.575 (future-work note); revisit if/when learned dispatch is proposed & \texttt{\small M2/ADR-2604.21330.md} \\
\arxivid{2604.21571} & M2 & Separable Expert Architecture: Composable Adapters and Deletable User Proxies & ADOPT & \texttt{docs/silicon-layer/separable-expert-architecture.md} (NEW), Silicon Layer spec, FoxFiber/FoxCompute multi-tenant design, FoxData privacy whitepaper & v1.49.575 (architecture spec); future Silicon Layer pipeline implementation & \texttt{\small M2/ADR-2604.21571.md} \\
\arxivid{2604.21725} & M2 & AEL: Agent Evolving Learning for Open-Ended Environments & ADOPT & \texttt{src/skill-creator/auto-load/bandit.ts} (NEW), reinforcement-series eligibility-trace integration & v1.49.575 spec; Half B Phase 812 ships HB-07 / CS25-19 implementation & \texttt{\small M2/ADR-2604.21725.md} \\
\arxivid{2604.21901} & M2 & GiVA: Gradient-Informed Bases for Vector-Based Adaptation & ADOPT & \texttt{docs/silicon-layer/benchmark-suite.md} (NEW), Silicon Layer benchmark spec, future Silicon Layer LoRA pipeline implementation, \texttt{src/learnable-k\_h/} integration & v1.49.575 (benchmark spec); future Silicon Layer benchmark suite implementation & \texttt{\small M2/ADR-2604.21901.md} \\
\arxivid{2604.21905} & M2 & Low-Rank Adaptation Redux for Large Models & TRACK & \texttt{docs/silicon-layer/lora-bibliography.md}, future Silicon Layer pipeline rank-selection documentation & v1.49.575 (bibliography entry); referenced as needed in future Silicon Layer milestones & \texttt{\small M2/ADR-2604.21905.md} \\
\arxivid{2604.21927} & M2 & Fine-Tuning Regimes Define Distinct Continual Learning Problems & ADOPT & \texttt{docs/silicon-layer/test-harness-spec.md} (NEW), \texttt{src/ab-harness/} ME-3 A/B harness, \texttt{src/representation-audit/} MD-6 audit trail, reinforcement-series test harness & v1.49.575 (test harness spec); v1.50 Silicon Layer benchmark suite (regime-fixed protocol enforced) & \texttt{\small M2/ADR-2604.21927.md} \\
\arxivid{2512.03048} & M3 & The Specification Trap: Why Static Value Alignment Alone Is Insufficient for Robust Alignment & ADOPT & \texttt{docs/substrate/specification-trap.md} + citations in The Space Between + College of Knowledge philosophy module & v1.49.575 (citation); v1.50 retrospective Half A (extended treatment) & \texttt{\small M3/ADR-2512.03048.md} \\
\arxivid{2601.18491} & M3 & AgentDoG: A Diagnostic Guardrail Framework for AI Agent Safety and Security & ADOPT & \texttt{src/safety/agentdog/} (new); integrates with existing Safety Warden BLOCK paths & v1.49.575 (HB-02, Phase 807) & \texttt{\small M3/ADR-2601.18491.md} \\
\arxivid{2604.20911} & M3 & Omission Constraints Decay While Commission Constraints Persist in Long-Context LLM Agents & ADOPT & \texttt{src/safety/std-calibration/} (new); calibration table at \texttt{.planning/safety/std-calibration.json} & v1.49.575 (HB-03, Phase 808) & \texttt{\small M3/ADR-2604.20911.md} \\
\arxivid{2604.20994} & M3 & Breaking MCP with Function Hijacking Attacks: Novel Threats for Function Calling and Agentic Models & ADOPT & \texttt{src/orchestration/planning-bridge/access-control/} + \texttt{docs/security/mcp-threat-model.md} & v1.49.575 (threat model + design); v1.49.576+ (runtime binding implementation) & \texttt{\small M3/ADR-2604.20994.md} \\
\arxivid{2604.21090} & M3 & Structural Quality Gaps in Practitioner AI Governance Prompts & ADOPT & \texttt{src/cartridge/linter/structural-completeness.ts} + cartridge promotion pipeline at \texttt{src/cartridge/} & v1.49.575 (HB-05, Phase 810) & \texttt{\small M3/ADR-2604.21090.md} \\
\arxivid{2604.21131} & M3 & Cross-Session Threats in AI Agents: Benchmark, Evaluation, and Algorithms & ADOPT & \texttt{src/safety/cross-session/} (architecture spec); \texttt{docs/security/kill-chain-taxonomy.md} & v1.49.575 (threat model + spec); v1.49.576+ (runtime correlator) & \texttt{\small M3/ADR-2604.21131.md} \\
\arxivid{2604.21477} & M3 & MCP Pitfall Lab: Exposing Developer Pitfalls in MCP Tool Server Security under Multi-Vector Attacks & ADOPT & \texttt{src/orchestration/planning-bridge/mcp-security/} + \texttt{src/cartridge/linter/mcp-pitfalls.ts} & v1.49.575 (Phase 805) & \texttt{\small M3/ADR-2604.21477.md} \\
\arxivid{2604.21718} & M3 & Building a Precise Video Language with Human-AI Oversight (CHAI) & ADOPT & \texttt{docs/capcom/} (terminology) + \texttt{src/reinforcement/} (pre-X/post-X capture pipeline) & v1.49.575 (documentation + signal pipe) & \texttt{\small M3/ADR-2604.21718.md} \\
\arxivid{2508.21720} & M4 & PosterForest: Hierarchical Multi-Agent Collaboration for Scientific Poster Generation & ADOPT & \texttt{src/skills/research-mission-generator/v2/} (design) + \texttt{docs/skills/mission-tree.md} & v1.49.575 (design + spec); v1.49.576+ (v2 implementation) & \texttt{\small M4/ADR-2508.21720.md} \\
\arxivid{2510.26615} & M4 & SARA / SlideAgent: Hybrid RAG with Compression Vectors & ADOPT & \texttt{src/college-of-knowledge/search/sara/} (design) + \texttt{docs/college-of-knowledge/retrieval.md} & v1.49.575 (design); v1.49.576+ (implementation) & \texttt{\small M4/ADR-2510.26615.md} \\
\arxivid{2511.00413} & M4 & Tree Training: Accelerating Agentic LLMs Training via Shared Prefix Reuse & ADOPT & \texttt{docs/silicon-layer/tree-training.md} (future-work notes) + v1.50 retrospective Half A & v1.49.575 (citation + future-work notes); future Silicon Layer milestone (implementation) & \texttt{\small M4/ADR-2511.00413.md} \\
\arxivid{2603.22823} & M4 & Empirical Comparison of Agent Communication Protocols for Task Orchestration & ADOPT & \texttt{docs/architecture/dacp.md} (citations) + \texttt{.planning/dacp/comparative-benchmark-plan.md} & v1.49.575 (methodology); v1.49.576+ (first comparative-benchmark run) & \texttt{\small M4/ADR-2603.22823.md} \\
\arxivid{2604.20848} & M4 & MATRAG: Multi-Agent Transparent Retrieval-Augmented Generation & ADOPT & \texttt{src/safety/capcom/transparency-score.ts} (immediate) + \texttt{src/orchestration/mission-crew/} (design) & v1.49.575 (transparency score); v1.49.576+ (Mission Crew Manifest implementation) & \texttt{\small M4/ADR-2604.20848.md} \\
\arxivid{2604.21255} & M4 & When Agents Look the Same: Quantifying Distillation-Induced Similarity in Tool-Use Behaviors & TRACK & \texttt{src/chipset-validation/rps-ags/} (TRACK design) + \texttt{docs/architecture/chipset.md} (citation) & TRACK / backlog; promotion-trigger conditions noted & \texttt{\small M4/ADR-2604.21255.md} \\
\arxivid{2604.21282} & M4 & Strategic Heterogeneous Multi-Agent Architecture for Cost-Effective Code Vulnerability Detection & ADOPT & \texttt{docs/architecture/chipset.md} (citations) + v1.50 retrospective Half A & v1.49.575 (citations); v1.50 retrospective Half A (extended treatment) & \texttt{\small M4/ADR-2604.21282.md} \\
\arxivid{2604.21444} & M4 & HiCrew: Hierarchical Reasoning for Long-Form Video Understanding & ADOPT & \texttt{src/orchestration/activation-profile/} (design) + \texttt{src/college-of-knowledge/search/} (design) & v1.49.575 (pattern reference); v1.49.576+ (implementation as part of activation-profile and College-of-Knowledge work) & \texttt{\small M4/ADR-2604.21444.md} \\
\arxivid{2604.21590} & M4 & AgenticQwen: Training Small Agentic Language Models with Dual Data Flywheels & ADOPT & \texttt{docs/architecture/dual-flywheel.md} + v1.50 retrospective Half A & v1.49.575 (citation); v1.50 retrospective Half A (extended treatment) & \texttt{\small M4/ADR-2604.21590.md} \\
\arxivid{2604.21794} & M4 & Learning to Communicate: Toward End-to-End Optimization of Multi-Agent Language Systems (DiffMAS) & TRACK & \texttt{docs/architecture/dacp.md} (citation as upper-fidelity reference) & TRACK; promotion-trigger conditions noted above & \texttt{\small M4/ADR-2604.21794.md} \\
\arxivid{2604.21816} & M4 & Tool Attention Is All You Need: Dynamic Tool Gating and Lazy Schema Loading & ADOPT & \texttt{src/orchestration/tool-attention/} & v1.49.575 (HB-01, Phase 806) & \texttt{\small M4/ADR-2604.21816.md} \\
\arxivid{2511.23159} & M5 & AI for Software Engineering: From Probable to Provable & ADOPT & \texttt{docs/substrate/the-space-between.md}, \texttt{docs/substrate/zfc-auditor.md} (citation only) & v1.49.575 & \texttt{\small M5/ADR-2511.23159.md} \\
\arxivid{2603.01170} & M5 & ATLAS: AI-Assisted Threat-to-Assertion Learning for SoC Security Verification & TRACK & \texttt{src-tauri/src/mcp/security/} (forthcoming) design doc; planning-bridge MCP design & v1.49.576+ (after M3 MCP-taxonomy consolidates) & \texttt{\small M5/ADR-2603.01170.md} \\
\arxivid{2604.14709} & M5 & HWE-Bench: Benchmarking LLM Agents on Real-World Hardware Bug Repair Tasks & TRACK & Future \texttt{examples/chipsets/hardware-flow/}; NASA Mission Series hardware-template doc & v1.49.576+ or future NASA hardware mission & \texttt{\small M5/ADR-2604.14709.md} \\
\arxivid{2604.18792} & M5 & Tractable Verification of Model Transformations: A Cutoff-Theorem Approach for DSLTrans & ADOPT & \texttt{docs/substrate/zfc-auditor.md} (citation + framing); future \texttt{coprocessors/math/symbex/} integration sketch & v1.49.575 & \texttt{\small M5/ADR-2604.18792.md} \\
\arxivid{2604.21187} & M5 & Doubly Saturated Ramsey Graphs: A Case Study in Computer-Assisted Mathematical Discovery & ADOPT & \texttt{docs/substrate/the-space-between.md}; \texttt{coprocessors/math/} integration sketch & v1.49.575 & \texttt{\small M5/ADR-2604.21187.md} \\
\arxivid{2604.21505} & M5 & Assessing the Impact of Requirement Ambiguity on LLM-based Function-Level Code Generation (Orchid) & ADOPT & \texttt{tools/lint-skill.mjs}; \texttt{.claude/skills/skill-integration/SKILL.md} template; \texttt{src/cartridge/validation/} & v1.49.575 (Phase 811) & \texttt{\small M5/ADR-2604.21505.md} \\
\arxivid{2604.21570} & M5 & SpecSyn: LLM-based Synthesis and Refinement of Formal Specifications & TRACK & \texttt{docs/substrate/zfc-auditor.md} (design doc); future \texttt{src/cartridge/validation/spec-synth/} & v1.50 (auditor scaffold), reference now & \texttt{\small M5/ADR-2604.21570.md} \\
\arxivid{2604.21598} & M5 & DryRUN: On the Role of Public Tests in LLM-Driven Code Generation & TRACK & \texttt{src/cartridge/validation/} synthesizer plug-in; v1.50 proof-companion test-generator & v1.49.576+ (HB-09 backlog slot) & \texttt{\small M5/ADR-2604.21598.md} \\
\arxivid{2510.18731} & M6 & Mitigating Lost-in-Multi-turn Conversation via Curriculum RL with Verifiable Accuracy and Abstention Rewards (RLAAR) & ADOPT & \texttt{docs/pedagogy/teacher-student-support-curriculum.md}; \texttt{src/reinforcement/} (abstention channel); \texttt{src/model-affinity/} & v1.49.575 (design doc); v1.49.576+ (channel wiring) & \texttt{\small M6/ADR-2510.18731.md} \\
\arxivid{2601.07262} & M6 & ColorBrowserAgent: Complex Long-Horizon Browser Agent with Adaptive Knowledge Evolution & ADOPT & \texttt{docs/capcom/feedback-capture.md}; future \texttt{src/capcom/feedback-capture/}; \texttt{src/output-structure/} & v1.49.575 (design doc); v1.49.576+ (implementation) & \texttt{\small M6/ADR-2601.07262.md} \\
\arxivid{2601.10003} & M6 & SocraticKG: Knowledge Graph Construction via QA-Driven Fact Extraction & ADOPT & \texttt{.claude/skills/research-mission-generator/SKILL.md}; future \texttt{.claude/skills/college-kb-builder/SKILL.md} & v1.49.575 & \texttt{\small M6/ADR-2601.10003.md} \\
\arxivid{2604.12710} & M6 & LASA: Language-Agnostic Semantic Alignment at the Semantic Bottleneck for LLM Safety & TRACK & Future \texttt{tests/cultural-sensitivity/}; cultural-sensitivity chip design & v1.49.576+ (focused experiment) & \texttt{\small M6/ADR-2604.12710.md} \\
\arxivid{2604.20487} & M6 & Knowledge Capsules: Structured Nonparametric Memory Units for LLMs & TRACK & \texttt{.college/} design doc; future \texttt{src/college/capsule/} & v1.51+ (after Silicon Layer kernel hooks) & \texttt{\small M6/ADR-2604.20487.md} \\
\arxivid{2604.20855} & M6 & Caesar: Deep Agentic Web Exploration for Creative Answer Synthesis & ADOPT & \texttt{.claude/skills/research-mission-generator/v2-design.md}; future \texttt{coprocessors/research/} & v1.49.575 (v2 design); v1.49.576+ (adversarial-refinement phase) & \texttt{\small M6/ADR-2604.20855.md} \\
\arxivid{2604.20904} & M6 & Reinforcing Privacy Reasoning in LLMs via Normative Simulacra from Fiction & TRACK & \texttt{.planning/the-hundred-voices/privacy-experiment.md}; \texttt{.planning/fox-companies/foxdata/} & v1.50+ (cross-project collaboration phase) & \texttt{\small M6/ADR-2604.20904.md} \\
\arxivid{2604.20920} & M6 & Forget, Then Recall: Learnable Compression and Selective Unfolding via Gist Sparse Attention & TRACK & \texttt{.college/} design doc; future \texttt{src/orchestration/college-retrieval/} & v1.50+ (College retrieval phase) & \texttt{\small M6/ADR-2604.20920.md} \\
\arxivid{2604.20996} & M6 & AFRILANGTUTOR: Advancing Language Tutoring and Culture Education in Low-Resource Languages with LLMs & TRACK & \texttt{.planning/missions/foxfire-heritage/}; future \texttt{examples/chipsets/heritage-language/} & v1.50+ (gated on sovereignty consultation) & \texttt{\small M6/ADR-2604.20996.md} \\
\arxivid{2604.21308} & M6 & CI-Work: Benchmarking Contextual Integrity in Enterprise LLM Agents & ADOPT & \texttt{.planning/fox-companies/foxfiber/threat-model.md}; future \texttt{src/safety-warden/} & v1.49.575 (threat-model doc); v1.49.576+ (metric scaffolding) & \texttt{\small M6/ADR-2604.21308.md} \\
\arxivid{2604.21352} & M6 & CARE: Counselor-Aligned Response Engine for Online Mental-Health Support & TRACK & \texttt{.college/departments/mind-body/} design doc; Safety Warden interaction policy & v1.50+ (mind-body chip phase) & \texttt{\small M6/ADR-2604.21352.md} \\
\arxivid{2604.21380} & M6 & Conjecture and Inquiry: Quantifying Software Performance Requirements via Interactive RAG (IRAP) & TRACK & \texttt{.claude/skills/vision-to-mission/}; \texttt{docs/capcom/elicitation-policy.md}; \texttt{src/model-affinity/} & v1.49.576+ (HB-10 backlog slot); vision-to-mission discipline now & \texttt{\small M6/ADR-2604.21380.md} \\
\arxivid{2604.21751} & M6 & Why are all LLMs Obsessed with Japanese Culture? On the Hidden Cultural and Regional Biases of LLMs & ADOPT & Future \texttt{tests/cultural-sensitivity/croq-pacific-rim.csv}; cultural-sensitivity chip & v1.49.575 (data drafting); v1.49.576+ (harness wiring) & \texttt{\small M6/ADR-2604.21751.md} \\
\arxivid{2604.21829} & M6 & Black-Box Skill Stealing Attack from Proprietary LLM Agents & TRACK & \texttt{docs/skill-economy/dolthub-threat-model.md}; \texttt{src/memory/grove-format.ts} integration & v1.49.576+ (HB-11 backlog slot); threat-model doc may land earlier & \texttt{\small M6/ADR-2604.21829.md} \\
\end{longtable}
\end{footnotesize}



\end{document}
